\documentclass{article}
\usepackage[margin=1in]{geometry}
\usepackage{times}
\usepackage{hyperref}
\begin{document}
\title{CultivaX Consolidated Documentation}
\maketitle
\tableofcontents
\newpage
\section{MSDD}
\begin{verbatim}
# FILE: CultivaX MSDD.pdf
# PAGES: 188



===== PAGE 1 =====
CULTIVAX – MASTER SYSTEM DESIGN DOCUMENT
(MSDD)
SECTION 1 — CROP TIMELINE INTELLIGENCE
SYSTEM (CTIS)
(Fully Corrected & Stabilized Engineering Specification)
1.1 Purpose
The Crop Timeline Intelligence System (CTIS) is the deterministic computational core of CultivaX.
CTIS is responsible for:
• Modeling crop lifecycle progression.
• Maintaining chronological integrity.
• Tracking all farmer actions immutably.
• Adjusting timeline elastically.
• Integrating weather and market risk.
• Supporting offline logging and replay.
• Enabling statistical adaptation without corrupting history.
• Preserving farmer authority.
CTIS is the authoritative state engine of each CropInstance.
1.2 Core Engineering Guarantees
CTIS guarantees:
1. Deterministic replay for identical action logs.
2. Chronological ordering preserved even after long offline periods.
3. No silent historical mutation.
4. Version-aware rule execution.
5. Separation of baseline vs actual progression.
6. State transitions explicitly logged.
7. Multi-crop isolation.
1.2.1 Chronological Invariant Definition
⚙ Define:
CTIS must enforce the following invariants:
1. action_effective_date >= sowing_date
2. action_effective_date <= server_time + tolerance_limit

===== PAGE 2 =====
3. 4. 5. 6. local_monotonic_counter must be strictly increasing per device session
No action may reference a stage not yet reachable in baseline progression
Effective dates must be strictly ordered per (date, local_counter)
No state transition may occur without event log record
1.3 CropInstance Model
Each CropInstance represents one crop grown on a specific land portion.
1.3.1 CropInstance Attributes (Revised)
Each instance must contain:
• crop_instance_id (UUID)
• farmer_id
• crop_type
• land_area
• region
• sowing_date
• seasonal_window_category
• current_day_number
• baseline_day_number
• current_growth_stage
• baseline_growth_stage
• state
• risk_index
• stress_score
• last_recalculated_timestamp
• rule_version_applied
• created_at
• updated_at
• is_archived
1.3.2 Baseline vs Actual Tracking
Two parallel references must always exist:
Field Meaning
baseline_day_number Ideal agronomic progression
current_day_number Progress after deviations
Field Meaning
baseline_growth_stage Stage per ideal template
current_growth_stage Stage after modifications
This separation prevents:
• Loss of agronomic reference

===== PAGE 3 =====
• Confusion in What-If simulations
• Incorrect regional comparisons
1.4 CropRuleTemplate Versioning
Each crop template must include:
• crop_type
• version_id
• effective_from_date
• stage_definitions
• risk_parameters
• irrigation_windows
• fertilizer_windows
• harvest_windows
When replaying historical logs:
CTIS must use the rule version active at the time of action.
This ensures:
• Historical determinism.
• No retroactive distortion when admin modifies rules.
If rule version changes:
• Active CropInstances retain rule_version_applied
• New rule only applies to:
o Future CropInstances
o Or manual replay triggered by admin
Add:
Rule change impact preview (admin-only simulation)
1.5 State Machine Definition (Formalized)
States and transition triggers:
State Trigger
Created CropInstance created
Active Sowing completed
Delayed Stage drift > delay threshold AND stress < risk threshold
AtRisk stress_score > risk threshold OR risk_index > threshold
ReadyToHarvest Within harvest window AND stress acceptable
Harvested Yield submitted
Closed Post-harvest services completed
Archived Manually archived or soft deleted
State transitions must:

===== PAGE 4 =====
• Be logged.
• Be event-triggered.
• Never be manually mutated without admin log.
• Maximum allowed stage delay
• Maximum early advancement
• Drift tolerance threshold
1.5.1 Stage Drift Constraints
⚙ Define:
CTIS must define:
• Maximum delay threshold per stage
• Maximum early advancement threshold
• Elasticity factor (configurable per crop)
• Hard cap on cumulative deviation days
Example:
Wheat:
• Max stage drift: ±14 days
• Stress acceleration limit: 1.5x baseline
1.6 Action Logging System
Each action is immutable and append-only.
1.6.1 ActionLog Structure
Each ActionLog must include:
• action_id
• crop_instance_id
• action_type
• action_subtype
• action_effective_date
• device_timestamp
• server_timestamp
• local_monotonic_counter
• metadata (JSON)
• source (web / whatsapp / offline / voice)
• is_override
• rule_version_at_action

===== PAGE 5 =====
• is_synced
Replay must use:
• action_effective_date
• not device_timestamp
Device timestamp is only for sync validation.
1.6.3 Action Impact Categories
⚙ Define:
Each action must be categorized as:
Category Description Replay Impact
Structural Alters timeline root (e.g., sowing date) Full replay required
Stage-Affecting May shift progression Partial replay
Stress-Affecting Alters stress_score Stress recalculation
Informational Notes / metadata only No timeline shift
Override Explicit farmer decision Overrides recommendation logic
Add field:
action_impact_type
1.7 Offline Handling & Timestamp Integrity
Offline actions may accumulate indefinitely.
On sync:
1. Validate extreme future timestamps.
2. Clamp impossible timestamps (log adjustment).
3. Sort by:
a. action_effective_date
b. local_monotonic_counter
4. Trigger replay from earliest unsynced action.
No rejection based solely on offline duration.
1.7.1 Temporal Anomaly Detection
⚙ Define:
System must detect:
• Excessive back-dated actions
• Excessive future-dated actions
• Large batch submission anomalies
• Monotonic counter resets

===== PAGE 6 =====
Flag to admin if anomaly score > threshold.
1.8 Replay Engine (Revised Deterministic Model)
Replay must be triggered when:
• Offline logs sync.
• Sowing date edited.
• Yield submitted.
• Admin modifies rule.
• Major override logged.
• Every N actions create replay snapshot
• Store compressed instance state
• Use incremental replay optimization
1.8.1 Replay Algorithm (Formal)
Step 1:
Identify earliest impacted effective_date.
Step 2:
Reset instance state to sowing_date.
Step 3:
Apply rule_version relevant at each action timestamp.
Step 4:
Apply each ActionLog in chronological order.
Step 5:
At each step:
• Update baseline progression.
• Apply deviation impact.
• Update stress score.
• Update risk index.
Step 6:
Finalize stage, risk, and projected harvest window.
Replay must be O(n) relative to number of actions.
No nested stage loops permitted.
1.8.2 Replay Snapshot Strategy
⚙ Define:
CTIS must create:

===== PAGE 7 =====
• Snapshot after every N actions (e.g., 50)
• Snapshot stores:
o current_stage
o baseline_stage
o stress_score
o risk_index
o deviation state
• Replay begins from nearest snapshot
Snapshots stored in:
crop_instance_snapshots table
1.9 Seasonal Window Categorization
Instead of single deviation formula, use:
seasonal_window_category ∈ {Early, Optimal, Late}
Each category has:
• Separate risk multiplier.
• Separate yield sensitivity.
Seasonal risk modifies:
• Stress sensitivity.
• Yield projection.
• Risk threshold.
1.9.1 Deviation Memory Buffer
⚙ Define:
CTIS must track:
• Consecutive deviation count
• Deviation trend slope
• Recurring pattern threshold
Stored in:
deviation_profile per crop instance
This feeds:
Behavioral adaptation model (Section 4)
1.10 Stress Score Model (Formalized)
Stress score must use exponential smoothing:

===== PAGE 8 =====
new_stress =
α × current_signal +
(1 − α) × previous_stress
Where:
α = configurable constant (e.g., 0.4)
Stress reset rules:
• At stage transition → reduce weight slightly.
• After corrective action → decay factor applied.
Stress ∈ [0, 1]
1.11 Risk Index
SystemRiskIndex =
(Weather_API_Risk × 0.7) +
(Farmer_Input_Risk × 0.3)
Risk thresholds configurable per crop.
Risk influences:
• State transition to AtRisk.
• Alert urgency.
• Recommendation intensity.
1.11.1 Signal Arbitration Hierarchy
⚙ Define Order of Authority:
When conflicting signals exist:
1. Backend ML Inference
2. Weather Risk
3. Deviation Penalty
4. Edge AI Signal
5. Farmer Manual Override (displayed but not forced)
If conflict occurs:
Higher authority signal weighted more heavily in stress integration.
1.12 Yield Submission
When yield submitted:
1. Log yield.
2. Trigger replay.
3. Compute final stress.
4. Compute yield verification score.

===== PAGE 9 =====
5. 6. Apply ML yield cap if necessary.
Update regional cluster (prospective only).
7. Set state = Harvested.
Regional updates must not retroactively alter past timeline.
Farmer Truth = UI Displayed
ML Truth = Training Reference
Regional Truth = Capped Aggregate
1.12.1 Yield Truth Separation Model
⚙ Define:
Yield stored in three distinct conceptual layers:
1. Farmer Truth
→ reported_yield
→ Always shown in UI
→ Never modified
2. ML Truth
→ ml_yield_value
→ Used for training only
→ Capped if exceeding biological limit
3. Regional Truth
→ Aggregated yield per cluster
→ Derived only from verified ML yield
Explicit statement:
ML logic must NEVER override Farmer Truth.
1.13 Regional Learning Stability Rule
Regional adaptation:
• Applied only to future CropInstances.
• Never retroactively alters current or past instance.
• Requires minimum sample threshold.
• Requires yield verification threshold.
Prevents infinite replay loops.
1.14 What-If Engine (Deep Copy Enforcement)
What-If must clone:
• Entire CropInstance state
• Stress history
• Deviation history

===== PAGE 10 =====
• Seasonal category
• Market snapshot
• Weather snapshot
Simulation occurs in isolated memory context.
No live mutation allowed.
1.15 Multi-Crop Isolation
Event dispatcher partitions queue by crop_instance_id.
Replay and risk computation fully isolated.
No cross-instance state contamination.
1.16 Deletion & Archival Policy
Crop deletion must:
• Soft delete (is_archived = true)
• Preserve action logs
• Preserve anonymized ML data
• Remove UI visibility
Hard deletion prohibited.
1.17 Sowing Date Modification Policy
Sowing date:
• Editable anytime by farmer.
• Must trigger full replay.
• Logged explicitly.
• Admin can override with log entry.
1.18 Failure Handling
If replay fails:
• Revert to last stable checkpoint.
• Log error.
• Notify admin.
• Revert to last valid snapshot

===== PAGE 11 =====
• Lock crop_instance
• Set state = “RecoveryRequired”
• Notify admin
• Prevent further action logging until resolved
Timeline must never enter undefined state.
1.19 Performance Requirements
Replay complexity: O(n)
Max supported actions per crop:
500–1000
Replay latency target:
< 2 seconds for 500 actions.
Async UI loading required.
1.20 Security Constraints
• No direct DB mutation bypassing event dispatcher.
• All state changes event-driven.
• All admin changes logged.
• No silent ML-based stage changes.
SECTION 2 — SERVICE ORCHESTRATION
ECOSYSTEM (SOE)
(Full Engineering Specification)
2.1 Purpose
The Service Orchestration Ecosystem (SOE) is responsible for:
• Connecting farmers to execution resources.
• Maintaining service provider trust scoring.
• Allowing independent operation outside CTIS.
• Integrating with CTIS when triggered by timeline needs.
SOE is facilitation-only, not enforcement.
It does not:

===== PAGE 12 =====
• Handle payments (prototype phase).
• Guarantee task completion.
• Enforce contracts.
• Assume liability.
2.2 Core Design Principles
SOE must:
1. Remain logically independent from CTIS.
2. Operate without corrupting crop timeline.
3. Maintain transparent trust computation.
4. Prevent provider manipulation.
5. Support admin moderation.
6. Scale independently if extracted into microservice.
2.3 SOE Subsystems
SOE consists of:
1. Provider Management Module
2. Equipment Management Module
3. Labor Management Module
4. Storage & Transport Suggestion Module
5. Trust Score Engine
6. Provider Verification System
7. Admin Moderation Layer
8. Request Handling System
2.4 Provider Model
2.4.1 ServiceProvider Entity
Attributes:
• provider_id (UUID)
• name
• contact_info
• region
• service_types (equipment / labor / transport / storage)
• verification_status
• trust_score
• complaint_count
• completion_count

===== PAGE 13 =====
• resolution_score
• rating_history
• created_at
• updated_at
• is_flagged
• is_suspended
2.4.2 Verification Status
Possible values:
• Pending
• Verified
• Suspended
• Rejected
Providers cannot appear in marketplace unless:
verification_status = Verified
Admin approval required.
2.5 Equipment Model
2.5.1 Equipment Entity
Attributes:
• equipment_id
• provider_id
• equipment_type
• availability_status
• hourly_rate (optional)
• location
• condition_notes
• is_flagged
• created_at
• updated_at
Equipment may be:
• Listed
• Temporarily unavailable
• Flagged by admin
• Deleted (soft delete)

===== PAGE 14 =====
2.6 Labor Model
Attributes:
• labor_id
• provider_id
• labor_type
• available_units
• daily_rate (optional)
• region
• is_flagged
• created_at
• updated_at
2.7 Service Request System
Farmers may initiate:
• Equipment inquiry
• Labor inquiry
• Transport inquiry
• Storage inquiry
Each request generates:
ServiceRequest entity:
• request_id
• farmer_id
• crop_instance_id (optional)
• service_type
• region
• timestamp
• status
• provider_contacted_list
• outcome_status
• feedback_submitted
2.7.1 Request States
• Initiated
• ProviderContacted
• FarmerAccepted
• Completed
• Cancelled
• Closed
SOE does not enforce these states.
They are advisory tracking only.

===== PAGE 15 =====
2.8 Trust Score Engine
TrustScore computed as:
TrustScore =
0.30 × CompletionRate
• 0.30 × ComplaintFrequencyInverse
• 0.20 × ResolutionQuality
• 0.15 × RatingStability
• 0.05 × ResponseTime
2.8.1 Metric Definitions
CompletionRate =
completed_requests / total_requests
ComplaintFrequencyInverse =
1 − (complaints / total_requests)
ResolutionQuality =
average resolution rating
RatingStability =
|initial_rating − revised_rating| penalty
ResponseTime =
normalized response latency
Weights configurable by admin.
2.9 Trust Safeguards
To prevent manipulation:
• One review per completed request.
• Farmer must have valid ServiceRequest to submit review.
• Self-rating prohibited.
• Provider cannot delete negative reviews.
• Admin can moderate fraudulent feedback.
2.10 Admin Governance in SOE
Admin can:
• Verify providers.

===== PAGE 16 =====
• Suspend providers.
• Delete equipment.
• Flag suspicious listings.
• Adjust trust weights.
• Override trust score.
• Delete users from SOE.
• View complaint logs.
All actions logged in AdminAuditLog.
2.11 Integration with CTIS
CTIS may trigger SOE suggestion when:
• Irrigation due.
• Fertilizer required.
• Harvest approaching.
• Stress high and labor recommended.
Integration must follow:
CTIS → SuggestServiceEvent
SOE → Return provider list
SOE does not alter timeline state.
2.12 Privacy in SOE
• No financial data stored (prototype).
• Contact details shared only upon farmer initiation.
• Complaint details anonymized for ML usage.
• Provider data removable by admin (soft delete).
2.13 Failure Handling
If no providers found:
• Display “No verified providers available.”
• Offer manual contact option.
If provider flagged:
• Hide from listing immediately.
If trust score computation fails:
• Use last stable score.

===== PAGE 17 =====
2.14 Deletion & Archival Policy
Providers:
• Soft delete only.
• Historical requests preserved.
• Trust history preserved.
Equipment:
• Soft delete.
• Historical logs preserved.
2.15 Security Constraints
• No provider can modify trust score.
• No direct DB update allowed.
• All review submissions validated against ServiceRequest.
• Admin changes logged.
2.16 Performance Constraints
System designed for:
• 100–200 users.
• Moderate request volume.
• No real-time bidding.
• No payment gateway.
Trust recomputation must be O(1) per update.
2.17 SOE Stability Guarantee
SOE must never:
• Modify crop timeline state.
• Trigger replay.
• Influence stress score.
• Affect ML training directly.
SOE is a separate pillar.

===== PAGE 18 =====
SECTION 3 — EVENT DISPATCHER & SYSTEM
COORDINATION LAYER
(This is the backbone that connects everything safely)
If this layer is weak, replay breaks, ordering breaks, and cross-module corruption happens.
So we go very deep here.
SECTION 2 — SOE ENHANCEMENTS
ENHANCEMENT 1 — Marketplace Fairness & Exposure
Control
Add Under: 2.8 Trust Score Engine
Add subsection:
2.8.3 Provider Exposure Fairness Model
⚙ Define:
When returning provider list:
Final Ranking Score =
(TrustScore × 0.85) +
(RandomExposureFactor × 0.10) +
(RegionalRelevanceWeight × 0.05)
Add constraint:
• Top 3 providers cannot occupy >70% exposure over rolling 30 days.
• If provider consistently dominates visibility, slight exposure decay applied.
Why:
Prevents:
• Winner-takes-all dominance
• Lock-in bias
• Monopoly formation
This makes your marketplace ethically mature.

===== PAGE 19 =====
ENHANCEMENT 2 — Trust Score Temporal Decay
Add Under: 2.8 Trust Score Engine
Add subsection:
2.8.4 Trust Decay Mechanism
⚙ Define:
If provider inactive for X days (e.g., 90 days):
trust_score_adjusted =
trust_score × decay_factor
Decay factor example:
0.98 per month inactivity.
CompletionRate must be weighted by recency.
Why:
Prevents stale providers from permanently dominating rankings.
Keeps marketplace dynamic.
ENHANCEMENT 3 — Complaint Escalation Logic
Add Under: 2.10 Admin Governance in SOE
Add subsection:
2.10.1 Automated Complaint Escalation
⚙ Define:
If:
complaint_count / total_requests > threshold (e.g., 20%)
Then:
• provider auto-flagged
• admin review required
• listing temporarily hidden
Escalation stages:
1. Warning
2. Temporary suspension
3. Permanent suspension
All logged in AdminAuditLog.

===== PAGE 20 =====
Why:
Adds safety and abuse resistance.
ENHANCEMENT 4 — Structured Complaint Taxonomy
Add Under: 2.7 Service Request System
Add:
Each complaint must include:
• complaint_category (delay, quality, non-response, fraud, pricing issue)
• severity_level
• resolution_status
Stored separately from rating.
Why:
Trust score becomes multi-dimensional instead of simplistic average.
ENHANCEMENT 5 — Service Provider Behavioral Stability
Metric
Add Under: 2.8.1 Metric Definitions
Add metric:
ConsistencyScore =
1 − variance(completion_time, resolution_score)
Weighted lightly (5–10%).
Why:
Consistency often matters more than occasional high rating.
ENHANCEMENT 6 — SOE–CTIS Interaction Safeguard
Enhance: 2.11 Integration with CTIS
Add subsection:
2.11.1 Recommendation Isolation Guarantee
Define:

===== PAGE 21 =====
SOE suggestions must:
• Never auto-trigger action logging.
• Never auto-update stress.
• Never auto-confirm task completion.
CTIS may suggest SOE, but SOE cannot feed back into CTIS without farmer-confirmed action.
Why:
Prevents hidden timeline mutation.
Preserves autonomy.
ENHANCEMENT 7 — Service Request Lifecycle Audit Trail
Add Under: 2.7 Service Request System
Add:
Every state transition in ServiceRequest must generate:
ServiceRequestEvent:
• event_type
• previous_state
• new_state
• timestamp
• actor_role
Stored in:
service_request_events table.
Why:
Improves governance & debugging.
ENHANCEMENT 8 — Anti-Fake Review Protection
Enhance: 2.9 Trust Safeguards
Add:
Review eligibility requires:
• ServiceRequest.status = Completed
• Minimum interaction time threshold
• One review per request
• Time window limit for review submission
Optional:
Require confirmation checkbox: “Service was actually provided.”

===== PAGE 22 =====
Why:
Reduces manipulation and fake boosting.
ENHANCEMENT 9 — Trust Score Transparency Layer
Add After: 2.8 Trust Score Engine
Add subsection:
2.8.5 Trust Explanation Model
Define:
Farmer UI may show:
TrustScore Breakdown:
• Completion Rate
• Complaint Ratio
• Resolution Quality
• Stability
• Response Pattern
But hide weight formula constants.
Why:
Transparency builds trust.
ENHANCEMENT 10 — Provider Data Retention & Privacy
Rule
Add Under: 2.12 Privacy in SOE
Define:
If provider requests deletion:
• Mark is_deleted = true
• Hide from listing
• Preserve historical request logs
• Preserve trust history anonymized
Hard deletion prohibited.
Why:
Aligns with your soft-delete philosophy from CTIS.

===== PAGE 23 =====
ENHANCEMENT 11 — Regional Saturation Control
Add Under: 2.4 Provider Model
Define:
If too many providers in one micro-region:
• Rank by trust
• But ensure minimum exposure rotation
Prevents overcrowding bias.
ENHANCEMENT 12 — Marketplace Risk Flag Integration
Add Under: 2.11 Integration with CTIS
Define:
If crop in AtRisk state:
SOE suggestions may highlight:
• Urgent labor
• Emergency irrigation equipment
But must remain advisory.
Why:
Better integration without state mutation.
SECTION 3 — EVENT DISPATCHER & SYSTEM
COORDINATION LAYER
3.1 Purpose
The Event Dispatcher is the coordination backbone of CultivaX.
It ensures:
• Deterministic execution order.
• Chronological integrity.
• No cross-module race conditions.
• Controlled communication between CTIS and SOE.

===== PAGE 24 =====
• Safe replay triggering.
• Isolation between CropInstances.
It replaces the need for heavy external message brokers in prototype scale.
3.2 Architectural Position
Logical placement:
User Interaction
→ Controller Layer
→ Event Dispatcher
→ Target Module (CTIS / SOE / ML / Admin)
No module directly mutates another module’s state.
All state changes must flow through events.
3.3 Event Model
Every state-changing action must generate an Event object.
3.3.1 Event Structure
Each Event contains:
• event_id (UUID)
• event_type
• crop_instance_id (nullable)
• module_target
• payload (JSON)
• created_at
• processed_at
• status
• priority
• rule_version_context
3.4 Event Categories
Events are categorized as:
1. CTIS Events
• CropCreated
• ActionLogged
• OverrideApplied
• SowingDateModified
• YieldSubmitted
• ReplayTriggered

===== PAGE 25 =====
2. ML Events
• MediaAnalyzed
• RiskUpdated
• RegionalClusterUpdated
3. SOE Events
• ServiceRequestCreated
• ProviderVerified
• TrustScoreUpdated
• ProviderFlagged
4. Admin Events
• RuleModified
• TrustWeightsModified
• UserDeleted
• ProviderSuspended
5. System Events
• WeatherUpdated
• MarketUpdated
• OfflineSyncCompleted
3.5 Partitioned Event Queue Model
Critical rule:
The event queue must be partitioned by:
partition_key = crop_instance_id
If crop_instance_id is null:
Use module-specific partition.
3.5.1 Why Partitioning Is Mandatory
Without partitioning:
• Events from Crop A may interleave with Crop B.
• Replay of one crop may conflict with another.
• Race conditions possible.
Partitioning ensures:
FIFO order within each crop instance.
3.6 Event Processing Order Guarantee
Within each partition:

===== PAGE 26 =====
Events must execute in strict FIFO order.
Execution sequence per event:
1. Validate event payload.
2. Lock partition.
3. Execute state mutation.
4. Trigger dependent computations.
5. Log execution.
6. Release partition lock.
No parallel execution within same crop instance.
Parallel allowed across different crop instances.
3.7 Replay Trigger Coordination
ReplayTriggered event must:
• Halt subsequent CTIS events for that partition.
• Recompute entire timeline.
• Resume queued events after completion.
Replay is atomic per crop instance.
3.8 Event Priority Model
Some events must have higher priority:
High Priority:
• ReplayTriggered
• RuleModified
• SowingDateModified
Medium Priority:
• ActionLogged
• YieldSubmitted
Low Priority:
• Notification events
• Trust recalculation
Priority applies only within same partition.
3.9 Failure Handling
If event processing fails:
• Mark event status = Failed
• Log error
• Retry up to configurable threshold
• If still failing:

===== PAGE 27 =====
o Escalate to Admin
o Do not proceed to next event in partition
This prevents silent corruption.
3.10 Concurrency Control
Concurrency policy:
• One worker per partition at a time.
• Global thread pool for multiple partitions.
• No shared mutable state outside partition context.
This ensures replay safety.
3.11 Deadlock Prevention Rule
Events must never:
• Generate synchronous nested events inside same partition.
• Call replay recursively.
Instead:
Nested effects must enqueue new events.
RegionalClusterUpdated must not auto-trigger ReplayTriggered for past instances.
3.12 Idempotency Requirement
All events must be idempotent.
Meaning:
If same event processed twice:
System state must remain consistent.
This protects against:
• Network retries.
• Crash recovery.
3.13 Offline Sync Handling
OfflineSyncCompleted event must:
• Insert unsynced actions.
• Trigger ReplayTriggered.
• Lock partition until replay completes.

===== PAGE 28 =====
No intermediate UI state exposed.
3.14 System Isolation Rules
SOE events must:
• Never enter CTIS partition unless explicitly integrated.
MarketUpdated must:
• Trigger CTIS recalculation only if harvest window impacted.
WeatherUpdated must:
• Trigger RiskUpdated event only.
No uncontrolled propagation.
3.15 Performance Constraints
Expected load:
• 100–200 users.
• ~5–20 events per user per day.
Queue must handle:
• ~2000 events/day comfortably.
Replay must not starve other partitions.
3.16 Crash Recovery
If system crashes mid-event:
• Use persisted event log.
• Resume from last unprocessed event.
• Ensure idempotent execution.
Event log must be durable.
3.17 Security Constraints
• No module may bypass dispatcher.
• No direct database writes allowed for state mutation.
• Admin interface must also generate events.
• All events must be authenticated.

===== PAGE 29 =====
SECTION 3 — EVENT DISPATCHER
ENHANCEMENTS
Structure:
• What to Add
• Where to Add It
• ⚙ What Exactly to Define
• Why It Strengthens the System
ENHANCEMENT 1 — Formal Event
Lifecycle State Machine
Add After: 3.3 Event Model
Add subsection:
3.3.2 Event Lifecycle Definition
⚙ Define Event States:
Each event must transition through:
1. Created
2. Persisted
3. Queued
4. Processing
5. Processed
6. Failed
7. Retried
8. DeadLettered
Add fields to event_log:
• retry_count
• failure_reason
• last_retry_at
Why:
Currently events have status, but lifecycle is not formalized.
This improves:
• Crash recovery
• Debugging
• Deterministic recovery

===== PAGE 30 =====
ENHANCEMENT 2 — Idempotency
Enforcement Strategy
Strengthen: 3.12 Idempotency Requirement
Add subsection:
3.12.1 Idempotency Enforcement Mechanism
⚙ Define:
For every event:
• Compute event_hash = hash(event_type + crop_instance_id + payload + rule_version_context)
• Store unique constraint on event_hash
• If duplicate event received:
o Do not re-execute mutation
o Return previous execution result
Additionally:
For ActionLogged:
• Unique constraint on action_id
Why:
You say “must be idempotent” — but not HOW.
Now it is enforceable.
ENHANCEMENT 3 — Dead Letter Queue
(DLQ) Strategy
Add Under: 3.9 Failure Handling
Add subsection:
3.9.1 Dead Letter Partition
⚙ Define:
If retry_count exceeds threshold:
• Move event to dead_letter_events table
• Mark partition as “RecoveryRequired”
• Notify admin
Dead letter event must not block other partitions.

===== PAGE 31 =====
Why:
Prevents infinite retry loop.
Improves resilience.
ENHANCEMENT 4 — Event Backpressure
Control
Add Under: 3.15 Performance Constraints
Add subsection:
3.15.1 Partition Backpressure Policy
⚙ Define:
If:
• One partition has excessive replay time
• Event queue length > threshold
Then:
• Reduce event fetch rate
• Prioritize high-priority events
• Prevent starvation of smaller partitions
Add max concurrent replay limit (e.g., 5).
Why:
Prevents large crop instance from blocking system.
ENHANCEMENT 5 — Atomic Replay Guard
Strengthen: 3.7 Replay Trigger Coordination
Add subsection:
3.7.1 Replay Guard Lock
⚙ Define:
During replay:
• Partition lock must block:
o ActionLogged
o YieldSubmitted

===== PAGE 32 =====
o SowingDateModified
• Queue events but do not process.
Replay must:
• Use transaction boundary
• Commit state only after full success
If replay fails:
• Revert to snapshot
• Mark replay_failed flag
Why:
Prevents half-applied timeline state.
ENHANCEMENT 6 — Event Dependency
Declaration
Add After: 3.4 Event Categories
Add subsection:
3.4.1 Event Dependency Map
Define allowed trigger chains:
Example:
ActionLogged → ReplayTriggered → RiskUpdated → NotificationSent
But prohibit:
RegionalClusterUpdated → ReplayTriggered (retroactively)
Define allowed graph explicitly.
Why:
Prevents accidental cyclic event design later.
ENHANCEMENT 7 — Event Integrity Hash
Chain
Add Under: 3.16 Crash Recovery
Add subsection:
3.16.1 Event Chain Integrity Verification

===== PAGE 33 =====
⚙ Define:
For each crop_instance:
• Maintain rolling hash:
event_chain_hash_n = hash(previous_hash + current_event_hash)
Store latest hash in crop_instance.
On replay:
• Recalculate chain
• Compare with stored
If mismatch:
• Flag tampering
• Lock instance
Why:
This is advanced but powerful.
Prevents silent event corruption.
Very impressive academically.
ENHANCEMENT 8 — Cross-Module
Isolation Enforcement
Strengthen: 3.14 System Isolation Rules
Add:
Each module must declare:
allowed_event_types
Dispatcher must validate:
• Target module matches allowed list
• Event cannot mutate non-owned tables
Example:
SOE cannot emit event targeting CTIS tables directly.
Why:
Prevents architectural drift.

===== PAGE 34 =====
ENHANCEMENT 9 — Replay Throttling
Protection
Add Under: 3.5 Partitioned Queue Model
Define:
If replay triggered repeatedly within short time:
• Merge replay events
• Debounce logic (e.g., 300ms window)
Prevents replay storm.
ENHANCEMENT 10 — System-Wide Kill
Switch
Add Under: 3.17 Security Constraints
Define:
Admin may:
• Disable ML-triggered events
• Disable clustering-triggered events
• Disable SOE-triggered recommendations
Without disabling core CTIS.
Why:
Critical control mechanism.
Especially for prototype instability.
SECTION 4 — MACHINE LEARNING &
STATISTICAL ARCHITECTURE
(Full Engineering Specification)
This section defines:
• What ML does
• What ML does NOT do
• Where ML sits in authority hierarchy

===== PAGE 35 =====
• How ML interacts with CTIS
• How regional learning works safely
• How yield verification influences adaptation
• How edge AI integrates without corruption
We will be extremely precise.
4.1 ML Design Philosophy
ML in CultivaX is:
• Assistive
• Risk-predictive
• Pattern-adaptive
• Region-sensitive
ML is NOT:
• Autonomous
• Prescriptive
• Deterministic stage controller
• Medical-grade disease diagnostician
Authority Model:
Timeline Engine > ML Risk Engine > Edge AI
Backend ML always overrides Edge AI.
4.2 ML Subsystems
The ML architecture consists of 4 layers:
Layer 1 — Deterministic Rule Engine
This is not statistical ML.
It includes:
• Stage-based recommendations
• Irrigation windows
• Fertilizer timing logic
• Harvest readiness window
• Pest alert thresholds
• Seasonal deviation multipliers
This layer never learns.
It is version-controlled via CropRuleTemplate.

===== PAGE 36 =====
Layer 2 — Risk Prediction Model
This layer predicts probability of stress escalation or stage instability.
Model Type:
Logistic Regression (prototype phase)
Inputs:
• stress_score
• seasonal_risk_index
• delay_days
• weather_risk
• recent deviation count
• stage position
Output:
• probability_of_risk (0–1)
This probability influences:
• AtRisk state transition
• Alert urgency
• Recommendation tone
This model must be:
• Retrainable
• Replaceable
• Versioned
Layer 3 — Behavioral Adaptation Model
This layer detects farmer-specific patterns.
Examples:
• Farmer consistently delays irrigation by 2 days.
• Farmer applies fertilizer earlier than recommended.
Mechanism:
Track deviation frequency.
If deviation pattern consistent over threshold:
Apply pre-emptive recommendation shift.
This does NOT change baseline template.
It modifies recommendation timing only.
Behavioral offset stored per farmer.
Layer 4 — Regional Clustering Model
This is the adaptive regional intelligence system.
Cluster key:

===== PAGE 37 =====
(crop_type + region + stage_group + seasonal_window_category)
Features for clustering:
• average_delay
• stress_trend
• yield_category
• weather_index_average
• deviation_pattern
Algorithm:
K-Means (prototype scale)
Cluster update condition:
• Sample size ≥ defined threshold
• YieldVerificationScore ≥ threshold
Important Rule:
Regional adjustments apply prospectively only.
Never retroactively modify completed instances.
4.3 Yield Verification & ML Interaction
Hybrid Yield Verification Model:
Inputs:
1. Self-reported yield
2. Quantitative yield (if provided)
3. Optional market sale confirmation
4. Stress history consistency
5. Seasonal risk context
Compute:
YieldVerificationScore (0–1)
If:
YieldVerificationScore ≥ threshold
→ Eligible for regional adaptation
If reported_yield > biological_limit:
→ ml_yield_value = biological_limit
→ reported_yield remains unchanged in UI
ML uses ml_yield_value only.
4.4 Stress Score Integration
Stress score sources:
1. Backend ML inference (video frame analysis)
2. Edge AI inference (low data mode)
3. Weather stress signal
4. Deviation penalty
5. Pest probability score

===== PAGE 38 =====
Stress smoothing formula:
new_stress =
α × signal +
(1 − α) × previous_stress
α configurable (default 0.4)
Stress affects:
• Risk model input
• State transition
• Alert system
4.5 Edge AI Architecture
Edge AI runs only when:
• Low-data mode enabled
• Device offline
Edge AI performs:
• Image-based stress classification
• Basic pest detection
• Color anomaly detection
Edge AI outputs structured summary:
• stress_estimate
• pest_probability
• confidence
Backend re-evaluates upon sync.
Edge AI never has final authority.
4.6 Backend Media Intelligence Pipeline
Video Processing Pipeline:
1. Extract frames (interval-based).
2. Run CNN classifier for stress detection.
3. Run object detection (YOLO-lite) for pest detection.
4. Aggregate frame-level results.
5. Compute:
a. average_stress
b. pest_density_index
c. anomaly_flag
Result stored as structured features.
Raw video retained 3 months.

===== PAGE 39 =====
4.7 Market Intelligence Modeling
Two prediction layers:
Short-Term:
• ARIMA
• 7-day forecast
Long-Term:
• Historical seasonal band
• Volatility estimation
Long-term prediction shown as range only.
Market model never alters timeline.
It influences recommendation messaging.
4.8 Weather Intelligence
SystemRiskIndex =
(Weather_API_Risk × 0.7) +
(Farmer_Input_Risk × 0.3)
Weather API failures:
Fallback to historical baseline.
Weather never forces state transition.
Only influences risk score.
4.9 ML Versioning
Each ML model must have:
• model_version
• training_date
• training_dataset_reference
• evaluation_metrics
When model updated:
• Only new CropInstances use new version.
• Historical instances not retroactively recalculated unless replay explicitly triggered.
4.10 ML Safety Constraints
ML must never:
• Directly modify sowing date.

===== PAGE 40 =====
• Directly skip stages.
• Directly delete action logs.
• Override farmer input.
ML influences only:
• risk_score
• recommendation timing
• alert intensity
4.11 Data Used for Training
Permitted:
• Anonymized stress history
• Capped yield values
• Deviation patterns
• Weather aggregates
Not permitted:
• Raw images beyond 3 months
• Identifiable user data
• Direct personal metadata
4.12 Performance Constraints
Risk inference latency:
< 200 ms
Video processing:
Async background job
Clustering recalculation:
Batch job (daily or threshold-based)
4.13 Failure Handling
If ML service unavailable:
• Fall back to rule engine.
• Use last known stress.
• Mark prediction as degraded.
System must remain usable.

===== PAGE 41 =====
SECTION 4 — ML & STATISTICAL
ARCHITECTURE ENHANCEMENTS
Structure:
• What to Enhance
• Where to Add It
• ⚙ What Exactly to Define
• Why It Matters
ENHANCEMENT 1 — Formal Authority
Hierarchy Table
Strengthen: 4.1 ML Design Philosophy
Add subsection:
4.1.1 Intelligence Authority Hierarchy
⚙ Define Explicit Authority Order:
1. Farmer Confirmed Action (UI truth)
2. CTIS Deterministic Rule Engine
3. Backend ML Inference
4. Weather Risk Model
5. Behavioral Adaptation Offset
6. Edge AI Inference
Define:
Lower authority signals may influence risk_score but cannot override higher authority signals.
Why:
Right now authority described verbally.
Formalizing hierarchy removes ambiguity.

===== PAGE 42 =====
ENHANCEMENT 2 — Confidence
Propagation Model
Add Under: 4.2 ML Subsystems
Add subsection:
4.2.1 Confidence & Data Sufficiency Propagation
⚙ Define:
Every ML output must include:
• prediction_value
• confidence_score (0–1)
• data_sufficiency_index
• model_version_used
Example:
{
"risk_probability": 0.68,
"confidence": 0.74,
"data_sufficiency": 0.52,
"model_version": "risk_v1.2"
}
Confidence affects:
risk_adjusted =
prediction × confidence
Low confidence → softer recommendation tone.
Why:
Prevents overconfident ML behavior.
Improves explainability.

===== PAGE 43 =====
ENHANCEMENT 3 — Signal Arbitration
Function
Strengthen: 4.4 Stress Score Integration
Add subsection:
4.4.1 Multi-Signal Arbitration Function
⚙ Define Weighted Integration Model:
stress_signal =
(w_backend × backend_ml) +
(w_weather × weather_signal) +
(w_deviation × deviation_penalty) +
(w_edge × edge_signal)
Where:
w_backend > w_weather > w_deviation > w_edge
Weights configurable per crop.
Why:
You combine signals — but don’t define weighting structure.
Now it's mathematically grounded.
ENHANCEMENT 4 — Long-Horizon
Forecast Model
Expand: 4.7 Market Intelligence Modeling
Add subsection:
4.7.1 Long-Term Forecast Confidence Band
⚙ Define:
When farmer asks at sowing time:
System must return:
• Expected harvest price band

===== PAGE 44 =====
• Volatility index
• Confidence label (Low / Medium / High)
• Data coverage ratio
Displayed as:
“Projected Harvest Price Range (Low Confidence)”
Why:
You mentioned harvest-at-sowing forecasting earlier.
Now formalize it.
ENHANCEMENT 5 — Model Drift
Monitoring System
Add Under: 4.9 ML Versioning
Add subsection:
4.9.1 Model Drift Detection Policy
⚙ Define:
After model update:
Compare:
• Accuracy change
• Precision change
• False positive rate change
If performance drop > threshold:
• Flag model_version
• Allow rollback
Store evaluation metrics in ml_models table.
Why:
Prevents silent ML degradation.

===== PAGE 45 =====
ENHANCEMENT 6 — Behavioral Adaptation
Safeguard
Strengthen: 4.2 Layer 3 — Behavioral Adaptation Model
Add:
Behavioral offset must:
• Never exceed ±X days
• Never alter baseline template
• Be reversible
• Expire after season ends
Why:
Prevents behavior from gradually distorting agronomic model.
ENHANCEMENT 7 — Cluster Stability
Validation
Strengthen: 4.2 Layer 4 — Regional Clustering Model
Add subsection:
4.2.4 Cluster Stability Check
⚙ Define:
Before applying new cluster parameters:
Check:
• Sample size ≥ threshold
• Variance within cluster ≤ threshold
• Yield verification average ≥ threshold
• Drift from previous cluster ≤ tolerance
If unstable:
Do not update cluster.
Why:
Prevents noisy clustering in small datasets.

===== PAGE 46 =====
ENHANCEMENT 8 — ML Feature
Traceability
Add Under: 4.11 Data Used for Training
Add subsection:
4.11.1 Feature Lineage Tracking
Define:
For each model training:
• Log dataset source
• Log date range
• Log feature schema
• Log preprocessing version
Stored in:
ml_training_audit table.
Why:
Research-grade reproducibility.
ENHANCEMENT 9 — Edge AI Confidence
Degradation Rule
Strengthen: 4.5 Edge AI Architecture
Add:
If:
confidence_score < threshold:
• Do not update stress directly
• Mark as “Low Confidence Edge Inference”
• Await backend confirmation
Why:
Prevents noisy offline inference from influencing timeline too aggressively.

===== PAGE 47 =====
ENHANCEMENT 10 — ML Kill Switch
Add Under: 4.10 ML Safety Constraints
Add:
Admin may disable:
• Risk prediction influence
• Behavioral adaptation
• Regional clustering influence
CTIS must fall back to deterministic rule engine.
Why:
Critical operational safeguard.
SECTION 5 — DATA MODEL &
DATABASE ARCHITECTURE
(Full Engineering Specification)
5.1 Database Philosophy
CultivaX uses:
• Relational core database (PostgreSQL recommended)
• JSON fields where flexibility required
• Strict referential integrity
• Soft-delete model for historical preservation
• Event log persistence
Primary goals:
• Deterministic replay
• Historical traceability

===== PAGE 48 =====
• ML-safe storage
• Governance transparency
5.2 Core Entity Groups
Database divided into logical groups:
1. CTIS Core Tables
2. SOE Tables
3. ML & Analytics Tables
4. Event Dispatcher Tables
5. Governance Tables
6. Media Storage Metadata
7. Configuration Tables
5.3 CTIS Core Tables
5.3.1 farmers
Field Type Notes
farmer_id UUID (PK)
name TEXT
phone TEXT unique
region TEXT
preferred_language TEXT
accessibility_settings JSONB
created_at TIMESTAMP
updated_at TIMESTAMP
is_deleted BOOLEAN soft delete
5.3.2 crop_instances
Field Type
crop_instance_id UUID (PK)
farmer_id FK

===== PAGE 49 =====
crop_type TEXT
land_area FLOAT
region TEXT
sowing_date DATE
seasonal_window_category TEXT
current_day_number INT
baseline_day_number INT
current_growth_stage TEXT
baseline_growth_stage TEXT
state TEXT
stress_score FLOAT
risk_index FLOAT
rule_version_applied TEXT
is_archived BOOLEAN
created_at TIMESTAMP
updated_at TIMESTAMP
Indexes required:
• farmer_id
• crop_type
• region
• state
5.3.3 action_logs
Field Type
action_id UUID (PK)
crop_instance_id FK
action_type TEXT
action_subtype TEXT
action_effective_date DATE
device_timestamp TIMESTAMP
server_timestamp TIMESTAMP
local_monotonic_counter INT
metadata JSONB
source TEXT
is_override BOOLEAN
rule_version_at_action TEXT
is_synced BOOLEAN
Critical index:
(crop_instance_id, action_effective_date, local_monotonic_counter)
Replay depends on this ordering.

===== PAGE 50 =====
5.3.4 yield_records
Field Type
yield_id UUID (PK)
crop_instance_id FK
reported_yield FLOAT
ml_yield_value FLOAT
yield_unit TEXT
verification_score FLOAT
submitted_at TIMESTAMP
5.4 CropRuleTemplate Tables
5.4.1 crop_rule_templates
Field Type
template_id UUID (PK)
crop_type TEXT
version_id TEXT
effective_from_date DATE
stage_definitions JSONB
irrigation_windows JSONB
fertilizer_windows JSONB
harvest_windows JSONB
risk_parameters JSONB
created_at TIMESTAMP
5.5 SOE Tables
5.5.1 service_providers
Field Type
provider_id UUID (PK)
name TEXT
region TEXT
verification_status TEXT

===== PAGE 51 =====
trust_score FLOAT
complaint_count INT
completion_count INT
resolution_score FLOAT
is_flagged BOOLEAN
is_suspended BOOLEAN
created_at TIMESTAMP
5.5.2 equipment
Field Type
equipment_id UUID
provider_id FK
equipment_type TEXT
availability_status TEXT
hourly_rate FLOAT
is_flagged BOOLEAN
5.5.3 service_requests
Field Type
request_id UUID
farmer_id FK
crop_instance_id FK (nullable)
service_type TEXT
region TEXT
status TEXT
created_at TIMESTAMP
closed_at TIMESTAMP
5.6 ML & Analytics Tables
5.6.1 stress_history
Field Type
stress_id UUID
crop_instance_id FK
stress_score FLOAT
stage TEXT
computed_at TIMESTAMP

===== PAGE 52 =====
5.6.2 regional_clusters
Field Type
cluster_id UUID
crop_type TEXT
region TEXT
stage_group TEXT
seasonal_window_category TEXT
avg_delay FLOAT
avg_yield FLOAT
sample_size INT
last_updated TIMESTAMP
5.6.3 ml_models
Field Type
model_id UUID
model_type TEXT
version TEXT
training_date TIMESTAMP
metrics JSONB
active BOOLEAN
5.7 Event Dispatcher Tables
5.7.1 event_log
Field Type
event_id UUID
event_type TEXT
crop_instance_id UUID (nullable)
payload JSONB
priority INT
status TEXT
created_at TIMESTAMP
processed_at TIMESTAMP
Index:
(crop_instance_id, created_at)

===== PAGE 53 =====
5.8 Governance Tables
5.8.1 admin_audit_log
Field Type
audit_id UUID
admin_id UUID
action_type TEXT
entity_type TEXT
entity_id UUID
before_value JSONB
after_value JSONB
timestamp TIMESTAMP
Immutable table.
5.9 Media Metadata Tables
5.9.1 media_records
Field Type
media_id UUID
crop_instance_id FK
media_type TEXT
file_path TEXT
uploaded_at TIMESTAMP
scheduled_deletion_at TIMESTAMP
extracted_features JSONB
Deletion cron job required.
5.10 Soft Delete Policy
No hard delete allowed for:
• crop_instances
• action_logs

===== PAGE 54 =====
• service_providers
• yield_records
Use:
is_deleted or is_archived flag.
This preserves replay and audit integrity.
5.11 Indexing Strategy
Mandatory indexes:
• action_logs (crop_instance_id + action_effective_date)
• event_log (crop_instance_id)
• crop_instances (farmer_id)
• regional_clusters (crop_type + region)
This ensures replay and clustering performance.
5.12 Data Retention Enforcement
Background job must:
• Delete raw media after 3 months.
• Retain extracted_features.
• Log deletion in audit.
5.13 Data Integrity Constraints
Foreign key constraints mandatory:
• crop_instance_id references crop_instances
• farmer_id references farmers
• provider_id references service_providers
No orphan records allowed.

===== PAGE 55 =====
5.14 Performance Considerations
Expected scale:
• 100–200 users
• ~500 actions per crop
• < 50k action_logs total
Database optimized for write-heavy action logs.
SECTION 5 — DATABASE & DATA
ARCHITECTURE ENHANCEMENTS
ENHANCEMENT 1 — Replay Snapshot Table
Add Under: 5.3 CTIS Core Tables
Add new table:
5.3.6 crop_instance_snapshots
Fields:
• snapshot_id (UUID)
• crop_instance_id (FK)
• snapshot_index (INT)
• stage
• baseline_stage
• stress_score
• risk_index
• deviation_profile (JSONB)
• rule_version_applied
• created_at
Index:
(crop_instance_id, snapshot_index)

===== PAGE 56 =====
Why:
Supports incremental replay optimization.
Prevents O(n) explosion at scale.
ENHANCEMENT 2 — Event Hash & Chain
Integrity Fields
Modify: 5.7.1 event_log
Add fields:
• event_hash (TEXT)
• chain_hash (TEXT)
• retry_count (INT)
• failure_reason (TEXT)
Add unique index:
(event_hash)
Why:
Enforces idempotency + tamper detection.
Strengthens crash recovery.
ENHANCEMENT 3 — Temporal Composite
Index Optimization
Strengthen: 5.11 Indexing Strategy
Add:
CREATE INDEX idx_action_replay
ON action_logs(crop_instance_id, action_effective_date DESC, local_monotonic_counter DESC);
Also:

===== PAGE 57 =====
CREATE INDEX idx_event_partition
ON event_log(crop_instance_id, priority, created_at);
Why:
Optimizes replay and partition processing.
ENHANCEMENT 4 — ML Training Audit
Table
Add Under: 5.6 ML & Analytics Tables
Add:
5.6.4 ml_training_audit
Fields:
• training_id (UUID)
• model_version
• dataset_reference
• feature_schema_version
• training_start
• training_end
• evaluation_metrics (JSONB)
• created_at
Why:
Supports ML reproducibility and version traceability.
ENHANCEMENT 5 — Deviation Profile
Table
Add Under: 5.3 CTIS Core Tables
Add:

===== PAGE 58 =====
5.3.7 deviation_profiles
Fields:
• profile_id (UUID)
• crop_instance_id (FK)
• consecutive_deviation_count
• deviation_trend_slope
• recurring_pattern_flag
• last_updated
Why:
Separates behavioral memory from action logs.
Keeps CTIS clean.
ENHANCEMENT 6 — Soft Delete Metadata
Standardization
Strengthen: 5.10 Soft Delete Policy
Add fields to all soft-deletable tables:
• deleted_at (TIMESTAMP)
• deleted_by (UUID nullable)
Why:
Improves audit trail clarity.
ENHANCEMENT 7 — Data Migration
Strategy Definition
Add After: 5.14 Performance Considerations
Add subsection:
5.15 Schema Migration Policy

===== PAGE 59 =====
Define:
• Use migration tool (Alembic recommended)
• Versioned migration scripts
• No direct schema edits in production
• Test migration on staging first
• Maintain backward compatibility for at least one API version
Why:
Prevents schema chaos later.
ENHANCEMENT 8 — Data Retention Policy
Clarification
Strengthen: 5.12 Data Retention Enforcement
Add:
Retention classes:
Data Type Retention Deletion Type
Raw media 3 months Hard delete
Extracted features Permanent Retained
Action logs Permanent Soft only
Event logs Permanent Soft only
Admin audit Permanent Immutable
Why:
Clarifies long-term data behavior.
ENHANCEMENT 9 — Partition-Aware Table
Strategy
Add Under: 5.1 Database Philosophy
Define:
At >50k action_logs:

===== PAGE 60 =====
• Enable PostgreSQL table partitioning by crop_instance_id
OR
• Time-based partitioning (by season/year)
Why:
Future scalability path documented.
ENHANCEMENT 10 — Read Replica
Readiness
Add Under: 5.14 Performance Considerations
Define:
If scale increases:
• Read-only queries (dashboard view)
• Served from read replica
• Writes remain primary DB only
Why:
Prepares architecture for scale.
ENHANCEMENT 11 — Structured JSON
Validation Rules
Add Under: 5.13 Data Integrity Constraints
Define:
JSONB fields must:
• Use Pydantic schema validation before insert
• Enforce schema version
Especially for:
• metadata

===== PAGE 61 =====
• extracted_features
• risk_parameters
Why:
Prevents JSON chaos over time.
SECTION 6 — MEDIA INTELLIGENCE
PIPELINE
(Edge AI + Backend Video Processing Architecture)
6.1 Purpose
The Media Intelligence Pipeline enables farmers to upload:
• Images
• Videos
• Voice notes
And extracts structured signals from media to support:
• Stress estimation
• Pest detection probability
• Growth anomaly detection
• Risk modeling in CTIS
Media analysis is assistive, not authoritative.
6.2 Authority Model
Hierarchy:
Backend ML → CTIS
Edge AI → Preliminary only
Raw media → Evidence only

===== PAGE 62 =====
Edge AI results are advisory.
Backend ML results are authoritative.
Timeline engine makes final decision.
6.3 Media Input Types
Supported:
1. Image (multiple angles allowed)
2. Short video clip
3. Voice input (for contextual notes)
4. Text description
Each media item stored in:
media_records table.
6.4 Edge AI Architecture
Edge AI activates when:
• Device offline
• Low-data mode enabled
• User manually selects low-data mode
6.4.1 Edge AI Capabilities
Edge AI performs:
• Lightweight image classification
• Stress probability estimation
• Basic pest probability
• Leaf discoloration detection
Edge AI does NOT:
• Diagnose specific disease
• Override backend decision
• Directly modify timeline

===== PAGE 63 =====
6.4.2 Edge AI Model
Model Type:
Lightweight CNN (MobileNet-based)
Inputs:
• Resized image
• Color normalized
Outputs:
• stress_estimate (0–1)
• pest_probability (0–1)
• anomaly_flag (boolean)
• confidence_score
Output format:
{
"stress_estimate": 0.62,
"pest_probability": 0.31,
"confidence": 0.78
}
6.4.3 Edge AI Data Handling
Edge AI generates:
• Structured feature JSON
• Optional compressed image (WebP)
Structured data sent immediately if possible.
Raw image uploaded when connection restored.
6.5 Backend Media Intelligence Pipeline
Backend processing is authoritative.
6.5.1 Image Processing Flow
1. Receive image.
2. Preprocess:

===== PAGE 64 =====
a. Resize
b. Normalize
3. Run CNN classifier.
4. Generate:
a. stress_probability
b. pest_probability
c. growth_stage_estimate
5. Store structured features.
6. Emit MediaAnalyzed event.
6.5.2 Video Processing Flow
Video processing must be asynchronous.
Steps:
1. Receive video.
2. Store raw video.
3. Extract frames at defined interval (e.g., 1 frame/sec).
4. For each frame:
a. Run stress classifier.
b. Run pest detection model (YOLO-lite).
5. Aggregate results:
a. average_stress
b. max_stress
c. pest_density_index
d. anomaly_duration
6. Store structured summary.
7. Emit MediaAnalyzed event.
Video processing must not block UI.
6.6 Feature Aggregation Logic
If multiple images uploaded:
Aggregate:
• Average stress
• Maximum stress
• Variance of stress
• Pest probability average
High variance indicates:

===== PAGE 65 =====
• Uneven crop health.
Aggregation must occur before updating CTIS.
6.7 Stress Integration into CTIS
When MediaAnalyzed event triggered:
CTIS must:
1. Update stress_score using smoothing formula.
2. Recompute risk_index.
3. Check threshold transitions.
4. Trigger alert if necessary.
Media analysis must not:
• Directly change growth stage.
• Skip timeline progression.
6.8 Voice Media Handling
Voice input:
• Converted to text (via external API or browser STT).
• Stored as transcript.
• May be parsed for keywords.
• Used as contextual metadata.
Voice does not directly influence stress score.
6.9 Media Storage Policy
Raw media stored for:
• 3 months from upload.
Fields:
• scheduled_deletion_at
After 3 months:

===== PAGE 66 =====
• Raw file deleted.
• extracted_features retained.
Deletion must be logged.
6.10 ML Dataset Extraction Policy
Only extracted structured features may be used for training:
Allowed:
• stress_probability
• pest_probability
• growth_stage_estimate
• yield_verification_score
• seasonal_window_category
Not allowed:
• raw images after retention period
• personal identifiers
• voice recordings
All ML datasets must be anonymized.
6.11 Failure Handling
If Edge AI fails:
• Allow upload without inference.
If Backend ML fails:
• Use last known stress value.
• Mark media analysis as pending.
• Notify admin if repeated failures.
If video processing fails:
• Log failure.
• Do not update stress score.
System must never crash timeline due to media failure.

===== PAGE 67 =====
6.12 Performance Constraints
Image processing:
< 500ms per image
Video processing:
Async background job
Max video length:
Defined limit (e.g., 30 seconds)
Max concurrent video jobs:
Configurable threshold
6.13 Security Constraints
• Media file paths must not expose internal storage.
• Access must require authentication.
• No direct public URL for raw media.
• Media deletion must be irreversible after retention period.
6.14 Limitations
Prototype does not:
• Perform full disease classification.
• Identify specific pest species.
• Provide medical-grade crop diagnosis.
System reports:
• Risk probability
• Stress estimation
• Pest likelihood
Clear disclaimer required in UI.

===== PAGE 68 =====
SECTION 6 — MEDIA INTELLIGENCE
PIPELINE ENHANCEMENTS
ENHANCEMENT 1 — Multi-Angle
Consistency Analysis
Add Under: 6.6 Feature Aggregation Logic
Add subsection:
6.6.1 Multi-Angle Variance Detection
⚙ Define:
When multiple images uploaded:
Compute:
• mean_stress
• max_stress
• stress_variance
• pest_variance
If stress_variance > threshold:
• Set uneven_growth_flag = true
• Increase risk_index slightly
• Display: “Field health uneven across area.”
Why:
You allow multi-angle uploads — but not leveraging variance insight.
This makes multi-angle meaningful.

===== PAGE 69 =====
ENHANCEMENT 2 — Image Quality
Validation Layer
Add Under: 6.5 Backend Media Intelligence Pipeline
Add subsection:
6.5.3 Image Quality Filtering
⚙ Define:
Before inference:
Check:
• Blur score (Laplacian variance)
• Brightness level
• Overexposure
• Motion distortion
• Image resolution minimum
If quality < threshold:
• Lower confidence_score
• Mark analysis = low_reliability
Do NOT reject automatically.
Why:
Improves ML signal quality without harming usability.
ENHANCEMENT 3 — Environmental Noise
Filtering
Add Under: 6.5 Video Processing Flow
Define:
During frame extraction:
Discard frames where:

===== PAGE 70 =====
• Brightness extreme
• Frame too dark
• Motion blur high
• Frame duplicate similarity > threshold
Aggregate only clean frames.
Why:
Video is noisy. This prevents polluted stress signals.
ENHANCEMENT 4 — Growth Stage Cross-
Validation
Add Under: 6.5.1 Image Processing Flow
Add:
growth_stage_estimate from CNN must be:
• Compared with CTIS current_growth_stage
• If difference > allowed threshold:
o Flag stage_discrepancy
o Notify farmer for confirmation
o Do NOT auto-shift stage
Why:
Prevents ML overriding deterministic timeline.
ENHANCEMENT 5 — Stress Escalation
Guardrail
Strengthen: 6.7 Stress Integration into CTIS
Add:
Stress update must respect:
• Maximum daily increase limit

===== PAGE 71 =====
• Maximum single-event jump threshold
• Confidence-weighted smoothing
Example:
new_stress =
α × (signal × confidence) +
(1 − α) × previous_stress
Why:
Prevents stress spikes from one bad image.
ENHANCEMENT 6 — Media Confidence
Propagation
Add Under: 6.7 Stress Integration
Define:
Each MediaAnalyzed event must include:
• stress_probability
• pest_probability
• confidence_score
• image_quality_score
• source (backend / edge)
CTIS must weight stress influence by confidence.
Why:
Makes ML behavior bounded and explainable.
ENHANCEMENT 7 — Edge AI Confidence
Degradation
Strengthen: 6.4 Edge AI Architecture
Add:

===== PAGE 72 =====
If:
confidence < threshold (e.g., 0.6)
Then:
• Do not modify stress_score
• Mark inference as provisional
• Await backend re-evaluation
Why:
Prevents noisy offline inference from influencing timeline strongly.
ENHANCEMENT 8 — Pest Density
Escalation Model
Add Under: 6.5 Video Processing Flow
Define:
If pest_density_index > threshold:
• Increase pest_alert_level
• Recommend inspection action
• Do NOT automatically change stage
Add pest_alert_history table.
Why:
Separates pest probability from stress probability.
Adds nuance.

===== PAGE 73 =====
ENHANCEMENT 9 — Media Anomaly
Persistence Tracking
Add Under: 6.6 Feature Aggregation
Define:
Track anomaly persistence:
If anomaly detected across consecutive uploads:
• Increase risk escalation weight
• Reduce smoothing dampening
Why:
Repeated signals matter more than isolated signals.
ENHANCEMENT 10 — Video Job Isolation
Safeguard
Add Under: 6.12 Performance Constraints
Define:
Video worker must:
• Run in isolated container
• Use CPU limit
• Max concurrent jobs configurable
• Not share DB connection pool with API service
Why:
Prevents video overload from affecting timeline determinism.

===== PAGE 74 =====
ENHANCEMENT 11 — ML Dataset Filtering
Gate
Strengthen: 6.10 ML Dataset Extraction Policy
Add:
Only include media features in training if:
• confidence_score ≥ threshold
• image_quality_score ≥ threshold
• YieldVerificationScore ≥ threshold
Why:
Prevents garbage-in garbage-out ML.
ENHANCEMENT 12 — Evidence-Based UI
Transparency
Add Under: 6.14 Limitations
Define:
UI must display:
• “Stress estimated from image”
• Confidence level
• Source of analysis (Edge / Backend)
Never display:
“Disease confirmed”
Only:
“High stress probability detected”
Why:
Protects ethically and legally.

===== PAGE 75 =====
SECTION 7 — ACCESSIBILITY & HUMAN-
CENTERED INTERACTION ARCHITECTURE
(Full Engineering Specification)
7.1 Purpose
The Accessibility Architecture ensures that CultivaX remains:
• Usable in outdoor sunlight conditions
• Usable by low-literacy farmers
• Usable on low-end devices
• Usable in low-bandwidth regions
• Usable by users with visual or cognitive difficulty
• Usable across language preferences
Accessibility must:
• Never alter system logic.
• Never bypass validation.
• Never corrupt timeline data.
• Remain UI-layer scoped only.
7.2 Architectural Placement
Accessibility Layer sits between:
User Interface
→ Accessibility Layer
→ Application Logic (CTIS / SOE)
Accessibility modifies:
• Presentation
• Input methods
• Rendering behavior
Accessibility does NOT modify:

===== PAGE 76 =====
• Timeline state
• ML logic
• Trust score
• Replay behavior
7.3 Accessibility Modes
Accessibility must support independent toggles:
1. High Contrast Mode
2. Sunlight Mode
3. Lite Mode
4. Low Literacy Mode
5. Cognitive Simplified Mode
6. Typography Adjustment Mode
7. Voice Assist Mode
8. Motion Reduction Mode
9. Structured WhatsApp Mode
Each toggle stored per-user.
7.4 High Contrast Mode
Purpose:
Improve visibility under varying light conditions.
Behavior:
• Increase contrast ratio > WCAG AA standard.
• Replace gradients with solid backgrounds.
• Increase border thickness.
• Remove subtle shadows.
• Increase icon weight.
Does not change layout structure.
7.5 Sunlight Mode
Purpose:
Outdoor farming usage.

===== PAGE 77 =====
Additional behaviors:
• Boost brightness theme.
• Remove animations.
• Increase button size.
• Increase padding.
• Replace small icons with large ones.
• Use strong edge outlines.
Sunlight Mode may auto-enable when device brightness high.
7.6 Lite Mode (Low Bandwidth / Low Device
Mode)
Purpose:
Support rural network conditions.
When enabled:
• Disable background animations.
• Disable auto-refresh.
• Preload only essential timeline data.
• Compress image uploads.
• Prefer text over images.
• Defer heavy charts.
Lite Mode affects:
• Media resolution.
• Image preview size.
• Notification polling frequency.
Lite Mode does NOT affect:
• Timeline calculations.
• Risk computations.
7.7 Low Literacy Mode
Purpose:
Support farmers with limited reading proficiency.

===== PAGE 78 =====
Changes:
• Replace long text with icons + short labels.
• Use color-coded stage indicators.
• Use visual progress bars instead of paragraph explanation.
• Show only essential daily action.
Low Literacy Mode hides:
• Advanced analytics.
• Deep statistical panels.
• Regional cluster charts.
No logic change occurs.
7.8 Cognitive Simplified Mode
Purpose:
Reduce cognitive overload.
UI adjustments:
• Show only:
o Today’s task
o Current risk
o Weather alert
o Market alert
• Hide What-If simulation.
• Hide regional comparison graphs.
This mode simplifies decision-making.
7.9 Typography Adjustment Mode
User may control:
• Font size
• Line height
• Letter spacing
• Word spacing
• Font weight
Settings stored in accessibility_settings JSON.

===== PAGE 79 =====
Must re-render dynamically without page reload.
7.10 Motion Reduction Mode
When enabled:
• Disable timeline animation.
• Disable stage transition animation.
• Disable loading animation loops.
• Replace animation with static indicators.
Helps users sensitive to motion.
7.11 Voice Assist Mode
Voice Assist includes:
• Voice-to-text input
• Text-to-speech playback
Voice Assist does NOT:
• Override text validation.
• Execute command without confirmation.
Voice inputs must convert to structured text before action logging.
7.12 Structured WhatsApp Mode
WhatsApp mode must:
• Present structured menu.
• Use numeric input options.
• Avoid free-text ambiguity.
• Confirm actions before logging.
Example:
Reply:
1 — Completed

===== PAGE 80 =====
2 — Delayed
3 — Upload photo
All WhatsApp actions must pass through same validation logic as web input.
7.13 Multi-Language Support
System must support:
• Hindi
• English
• Expandable architecture for additional languages
Language design:
• No hard-coded strings.
• All UI text stored in translation files.
• Dynamic language switching.
No ML or timeline logic language-dependent.
7.14 Accessibility Data Storage
In farmers table:
accessibility_settings JSON includes:
{
"contrast_mode": true,
"sunlight_mode": false,
"lite_mode": true,
"low_literacy_mode": false,
"font_scale": 1.2,
"motion_reduction": true,
"voice_assist": false
}
Must load on login.

===== PAGE 81 =====
7.15 Accessibility Testing Requirements
Each major UI action must be testable in:
• Normal mode
• High contrast
• Lite mode
• WhatsApp mode
All core functions must remain accessible in Lite mode.
7.16 Failure Handling
If accessibility setting conflicts:
Example:
Sunlight + Lite + Motion Reduction
System must:
• Merge behaviors safely.
• Never disable input validation.
• Never disable confirmation prompts.
7.17 Non-Functional Constraints
Accessibility must:
• Not degrade performance beyond acceptable limits.
• Not cause layout break on small devices.
• Not interfere with offline logging.
7.18 Security Considerations
Accessibility must not:
• Bypass confirmation dialogs.
• Bypass override logging.
• Execute command without structured validation.

===== PAGE 82 =====
Voice commands must always require structured confirmation before action logging.
7.19 Separation from Core Logic
Accessibility Layer must:
• Be implemented in frontend.
• Never modify backend business logic.
• Never write directly to CTIS tables.
All user inputs pass through same validation pipeline.
SECTION 6 — MEDIA INTELLIGENCE
PIPELINE ENHANCEMENTS
ENHANCEMENT 1 — Multi-Angle
Consistency Analysis
Add Under: 6.6 Feature Aggregation Logic
Add subsection:
6.6.1 Multi-Angle Variance Detection
⚙ Define:
When multiple images uploaded:
Compute:
• mean_stress
• max_stress
• stress_variance
• pest_variance
If stress_variance > threshold:
• Set uneven_growth_flag = true

===== PAGE 83 =====
• Increase risk_index slightly
• Display: “Field health uneven across area.”
Why:
You allow multi-angle uploads — but not leveraging variance insight.
This makes multi-angle meaningful.
ENHANCEMENT 2 — Image Quality
Validation Layer
Add Under: 6.5 Backend Media Intelligence Pipeline
Add subsection:
6.5.3 Image Quality Filtering
⚙ Define:
Before inference:
Check:
• Blur score (Laplacian variance)
• Brightness level
• Overexposure
• Motion distortion
• Image resolution minimum
If quality < threshold:
• Lower confidence_score
• Mark analysis = low_reliability
Do NOT reject automatically.
Why:
Improves ML signal quality without harming usability.

===== PAGE 84 =====
ENHANCEMENT 3 — Environmental Noise
Filtering
Add Under: 6.5 Video Processing Flow
Define:
During frame extraction:
Discard frames where:
• Brightness extreme
• Frame too dark
• Motion blur high
• Frame duplicate similarity > threshold
Aggregate only clean frames.
Why:
Video is noisy. This prevents polluted stress signals.
ENHANCEMENT 4 — Growth Stage Cross-
Validation
Add Under: 6.5.1 Image Processing Flow
Add:
growth_stage_estimate from CNN must be:
• Compared with CTIS current_growth_stage
• If difference > allowed threshold:
o Flag stage_discrepancy
o Notify farmer for confirmation
o Do NOT auto-shift stage
Why:
Prevents ML overriding deterministic timeline.

===== PAGE 85 =====
ENHANCEMENT 5 — Stress Escalation
Guardrail
Strengthen: 6.7 Stress Integration into CTIS
Add:
Stress update must respect:
• Maximum daily increase limit
• Maximum single-event jump threshold
• Confidence-weighted smoothing
Example:
new_stress =
α × (signal × confidence) +
(1 − α) × previous_stress
Why:
Prevents stress spikes from one bad image.
ENHANCEMENT 6 — Media Confidence
Propagation
Add Under: 6.7 Stress Integration
Define:
Each MediaAnalyzed event must include:
• stress_probability
• pest_probability
• confidence_score
• image_quality_score
• source (backend / edge)
CTIS must weight stress influence by confidence.

===== PAGE 86 =====
Why:
Makes ML behavior bounded and explainable.
ENHANCEMENT 7 — Edge AI Confidence
Degradation
Strengthen: 6.4 Edge AI Architecture
Add:
If:
confidence < threshold (e.g., 0.6)
Then:
• Do not modify stress_score
• Mark inference as provisional
• Await backend re-evaluation
Why:
Prevents noisy offline inference from influencing timeline strongly.
ENHANCEMENT 8 — Pest Density
Escalation Model
Add Under: 6.5 Video Processing Flow
Define:
If pest_density_index > threshold:
• Increase pest_alert_level
• Recommend inspection action
• Do NOT automatically change stage
Add pest_alert_history table.

===== PAGE 87 =====
Why:
Separates pest probability from stress probability.
Adds nuance.
ENHANCEMENT 9 — Media Anomaly
Persistence Tracking
Add Under: 6.6 Feature Aggregation
Define:
Track anomaly persistence:
If anomaly detected across consecutive uploads:
• Increase risk escalation weight
• Reduce smoothing dampening
Why:
Repeated signals matter more than isolated signals.
ENHANCEMENT 10 — Video Job Isolation
Safeguard
Add Under: 6.12 Performance Constraints
Define:
Video worker must:
• Run in isolated container
• Use CPU limit
• Max concurrent jobs configurable
• Not share DB connection pool with API service

===== PAGE 88 =====
Why:
Prevents video overload from affecting timeline determinism.
ENHANCEMENT 11 — ML Dataset Filtering
Gate
Strengthen: 6.10 ML Dataset Extraction Policy
Add:
Only include media features in training if:
• confidence_score ≥ threshold
• image_quality_score ≥ threshold
• YieldVerificationScore ≥ threshold
Why:
Prevents garbage-in garbage-out ML.
ENHANCEMENT 12 — Evidence-Based UI
Transparency
Add Under: 6.14 Limitations
Define:
UI must display:
• “Stress estimated from image”
• Confidence level
• Source of analysis (Edge / Backend)
Never display:
“Disease confirmed”
Only:

===== PAGE 89 =====
“High stress probability detected”
Why:
Protects ethically and legally.
SECTION 8 — API CONTRACTS &
INTERFACE DESIGN
(Full Engineering Specification)
8.1 Purpose
The API layer provides structured communication between:
• Web frontend (PWA)
• WhatsApp bot
• Admin frontend
• Backend CTIS
• SOE
• ML services
• Event Dispatcher
APIs must:
• Be versioned
• Be authenticated
• Be idempotent where required
• Not expose internal state mutation endpoints
8.2 API Design Principles
1. RESTful structure
2. JWT-based authentication
3. Role-based access control (Farmer / Provider / Admin)
4. Versioned routes: /api/v1/

===== PAGE 90 =====
5. 6. No direct database write endpoints
All state-changing requests generate events
8.3 Authentication APIs
8.3.1 Register Farmer
POST /api/v1/auth/register
Request:
{
"name": "Ramesh",
"phone": "XXXXXXXXXX",
"region": "Haryana",
"language": "hi"
Response:
}
{
"farmer_id": "uuid",
"token": "jwt_token"
}
8.3.2 Login
POST /api/v1/auth/login
Returns JWT token.
JWT must include:
• user_id
• role
• expiration
8.4 Crop Timeline APIs (CTIS)

===== PAGE 91 =====
8.4.1 Create Crop Instance
POST /api/v1/crops
Request:
{
"crop_type": "wheat",
"land_area": 2.5,
"region": "Haryana",
"sowing_date": "2026-11-15"
}
Action:
→ Generates CropCreated event.
8.4.2 Get Crop Timeline
GET /api/v1/crops/{crop_instance_id}
Returns:
• current_stage
• baseline_stage
• stress_score
• risk_index
• seasonal_risk
• projected_harvest_window
• recent_actions
8.4.3 Log Action
POST /api/v1/crops/{id}/actions
Request:
{
"action_type": "irrigation",
"effective_date": "2026-12-10",
"metadata": {}
}
Triggers:
ActionLogged event.

===== PAGE 92 =====
No direct timeline mutation allowed.
8.4.4 Submit Yield
POST /api/v1/crops/{id}/yield
Triggers:
YieldSubmitted event.
Returns:
YieldVerificationScore.
8.4.5 Modify Sowing Date
PUT /api/v1/crops/{id}/sowing-date
Triggers:
SowingDateModified event.
Replay mandatory.
8.5 Media APIs
8.5.1 Upload Image
POST /api/v1/crops/{id}/media/image
Returns:
Upload confirmation + processing status.
Triggers:
MediaUploaded event.
8.5.2 Upload Video
POST /api/v1/crops/{id}/media/video

===== PAGE 93 =====
Async processing.
Returns:
Processing ID.
8.5.3 Get Media Analysis Result
GET /api/v1/crops/{id}/media/{media_id}
Returns structured features only.
Raw media not directly exposed.
8.6 What-If API
POST /api/v1/crops/{id}/what-if
Request:
{
"hypothetical_action": "delay_irrigation",
"delay_days": 3
}
Returns:
Projected timeline shift, risk change, yield impact.
No mutation occurs.
8.7 Service Orchestration APIs (SOE)
8.7.1 Create Service Request
POST /api/v1/services/request
Triggers:
ServiceRequestCreated event.

===== PAGE 94 =====
8.7.2 Get Providers
GET /api/v1/services/providers?region=X&type=equipment
Returns verified providers only.
8.7.3 Submit Review
POST /api/v1/services/review
Validates:
Must have completed ServiceRequest.
Triggers:
TrustScoreUpdated event.
8.8 Admin APIs
Admin endpoints must require:
• Admin JWT
• Role verification
8.8.1 Modify Crop Rule
PUT /api/v1/admin/crop-rules
Triggers:
RuleModified event.
Replay may be triggered prospectively.
8.8.2 Verify Provider
PUT /api/v1/admin/providers/{id}/verify

===== PAGE 95 =====
8.8.3 Suspend Provider
PUT /api/v1/admin/providers/{id}/suspend
8.8.4 Adjust Trust Weights
PUT /api/v1/admin/trust-config
8.8.5 View Audit Logs
GET /api/v1/admin/audit
8.9 WhatsApp Integration API
WhatsApp Webhook Endpoint:
POST /api/v1/whatsapp/webhook
Process:
1. Validate signature.
2. Parse structured reply.
3. Convert to action.
4. Trigger ActionLogged event.
WhatsApp cannot directly mutate DB.
8.10 Event Log API (Internal)
Internal-only endpoints:
POST /internal/events
External clients must never access this.

===== PAGE 96 =====
8.11 Error Handling Standard
All API responses must follow:
{
"success": false,
"error_code": "INVALID_ACTION",
"message": "Human readable message"
}
No stack traces returned.
8.12 Rate Limiting
To prevent abuse:
• Action logging limited per minute.
• Media upload size limit enforced.
• WhatsApp spam protection.
8.13 Idempotency Requirements
Endpoints that must be idempotent:
• Yield submission
• Action logging
• Media upload retry
Use idempotency_key header.
8.14 Security Constraints
• HTTPS mandatory.
• JWT expiration enforced.
• Role-based route protection.
• Admin endpoints separate namespace.
• Media access requires authentication.

===== PAGE 97 =====
SECTION 8 — API CONTRACT &
INTERFACE ENHANCEMENTS
ENHANCEMENT 1 — API Version
Deprecation Policy
Add After: 8.2 API Design Principles
Add subsection:
8.2.1 API Version Lifecycle Policy
⚙ Define:
• All APIs under /api/v1/
• Breaking changes require /api/v2/
• Maintain backward compatibility for minimum one major release
• Deprecated endpoints marked with:
o X-API-Deprecated: true header
o Sunset date
Why:
Prevents chaos when schema evolves.
ENHANCEMENT 2 — Bulk Offline Sync
Endpoint
Add Under: 8.4 Crop Timeline APIs
Add:
8.4.6 Sync Batch Actions
POST /api/v1/crops/{id}/sync-batch
Request:

===== PAGE 98 =====
{
"actions": [
{ action_payload },
{ action_payload }
],
"device_session_id": "uuid",
"last_sync_at": "timestamp"
}
Response:
• accepted_actions_count
• rejected_actions
• replay_status
Why:
Currently offline sync is conceptual, not contract-defined.
This formalizes it.
ENHANCEMENT 3 — Idempotency Header
Standardization
Strengthen: 8.13 Idempotency Requirements
Add:
All state-changing endpoints must require:
Header:
Idempotency-Key: <UUID>
Server must:
• Store key in request_log table
• Reject duplicate within defined time window
Applies to:
• Action logging
• Yield submission
• Media upload
• Service request creation

===== PAGE 99 =====
Why:
Prevents double submissions during poor connectivity.
ENHANCEMENT 4 — Response Envelope
Standardization
Strengthen: 8.11 Error Handling Standard
Define universal response structure:
{
"success": true,
"data": {},
"meta": {
"timestamp": "",
"version": "v1",
"request_id": ""
}
}
Errors:
{
"success": false,
"error": {
"code": "",
"message": "",
"details": ""
}
}
Why:
Uniform structure improves maintainability and debugging.

===== PAGE 100 =====
ENHANCEMENT 5 — Event Trigger
Declaration in API Docs
Add Under Each State-Changing API
For example:
Under:
8.4.3 Log Action
Add:
Triggers:
• ActionLogged event
• ReplayTriggered event (conditional)
Under:
8.4.4 Submit Yield
Triggers:
• YieldSubmitted
• ReplayTriggered
• RegionalClusterUpdated (conditional)
Why:
Makes API layer explicitly event-driven.
Improves clarity.
ENHANCEMENT 6 — Role Scope
Enforcement Matrix
Add After: 8.2 API Design Principles
Add subsection:
8.2.2 Role Scope Enforcement Table
Define allowed roles per endpoint in matrix form.
Include:

===== PAGE 101 =====
• Farmer
• Provider
• Admin
• Internal Worker
Why:
Prevents future endpoint creep.
ENHANCEMENT 7 — Rate Limit Tiering
Strengthen: 8.12 Rate Limiting
Define tiers:
Farmer:
• X requests/min
Provider:
• Y requests/min
Admin:
• Higher threshold
Media uploads:
• Separate rate limit
WhatsApp webhook:
• IP-based throttling
Why:
Prevents DOS via action spam.

===== PAGE 102 =====
ENHANCEMENT 8 — Long-Running
Operation Status API
Add Under: 8.5 Media APIs
Add:
GET /api/v1/operations/{operation_id}
Used for:
• Video processing status
• Replay processing status
• Batch clustering status
Response:
• status (pending / processing / completed / failed)
• progress_percentage
Why:
Prevents frontend polling chaos.
ENHANCEMENT 9 — Audit Trace Exposure
(Read-Only)
Add Under: 8.8 Admin APIs
Add:
GET /api/v1/admin/events/{crop_instance_id}
Returns:
• Event history
• Replay triggers
• State transitions
Read-only.

===== PAGE 103 =====
Why:
Improves transparency and debugging.
ENHANCEMENT 10 — Strict Input Schema
Validation
Strengthen: 8.14 Security Constraints
Define:
• All requests validated via Pydantic schemas
• Reject unknown fields
• Enforce strict type validation
• Validate JSON schema version
Why:
Prevents malformed JSON from corrupting logic.
ENHANCEMENT 11 — API Health &
Readiness Endpoints
Add New Section: 8.15 System Health Endpoints
Define:
GET /api/v1/health
Returns:
• DB connectivity
• Event queue status
• ML service status
• Media worker status
Why:
Essential for deployment and monitoring.

===== PAGE 104 =====
ENHANCEMENT 12 — Soft Delete Visibility
Policy
Add Under: 8.4 Crop APIs
Define:
Archived crops:
• Not returned in default GET
• Available via:
?include_archived=true
Why:
Keeps UI clean while preserving audit.
SECTION 9 — DEPLOYMENT &
INFRASTRUCTURE ARCHITECTURE
(Full Engineering Specification)
9.1 Deployment Philosophy
CultivaX is designed as:
• Modular monolith backend
• Separate admin frontend
• PWA frontend
• Optional WhatsApp integration service
• ML-enabled backend processing
• Async video processing
• Event-driven internal architecture

===== PAGE 105 =====
Prototype scale target:
100–200 concurrent users.
Infrastructure must support:
• Deterministic replay
• Media processing
• Background jobs
• Secure storage
• Crash recovery
9.2 High-Level Infrastructure Layout
User Devices
→ PWA Frontend
→ Backend API Service
→ Event Dispatcher
→ Database
→ ML Processing Worker
→ Media Storage
→ Admin Frontend
→ WhatsApp Webhook Service
9.3 Frontend Deployment
PWA (Farmer Interface)
Technology:
• Next.js (TypeScript)
• Hosted on:
o Google Cloud Run OR
o Vercel (prototype)
Must support:
• Service Workers
• Offline caching
• IndexedDB storage
• Background sync
No direct database access.
All via API.

===== PAGE 106 =====
Admin Frontend
Separate deployment instance.
Access restricted by:
• Admin JWT
• IP filtering (optional)
• Role validation
Must not share session cookies with farmer frontend.
9.4 Backend API Service
Technology:
• FastAPI (Python)
Hosted on:
• Google Cloud Run
Responsibilities:
• CTIS
• SOE
• ML orchestration
• Event dispatcher
• Replay engine
• API layer
Must run in:
Stateless mode.
All state stored in DB.
9.5 Database Deployment
Technology:
• PostgreSQL (Cloud SQL on GCP recommended)

===== PAGE 107 =====
Must support:
• 50k+ action_logs
• Partitioned indexes
• JSONB support
• ACID transactions
Backups:
• Daily snapshot
• Point-in-time recovery enabled
9.6 Event Dispatcher Infrastructure
Implementation inside backend service.
Uses:
• In-memory queue for prototype
OR
• Redis (optional improvement)
Critical requirement:
Event log persisted in database before processing.
Ensures crash recovery.
9.7 ML Processing Infrastructure
Two components:
9.7.1 Risk & Statistical Model
Runs inside main backend process.
Lightweight inference.

===== PAGE 108 =====
9.7.2 Video Processing Worker
Separate async worker process.
Deployed as:
• Separate Cloud Run service
OR
• Background worker container
Consumes:
• MediaUploaded events.
Processes video frames asynchronously.
Updates structured features.
Must not block main API.
9.8 Media Storage Infrastructure
Raw media stored in:
• Google Cloud Storage bucket
Policy:
• Private bucket.
• Signed URLs for temporary access.
• Automatic lifecycle rule:
o Delete after 3 months.
Structured extracted features stored in PostgreSQL.
9.9 WhatsApp Integration Deployment
Webhook endpoint:
• Deployed as separate lightweight service
OR
• Separate route in backend
Must:

===== PAGE 109 =====
• Validate Meta signature.
• Parse structured input.
• Generate ActionLogged event.
No direct DB writes allowed.
9.10 Security Infrastructure
Mandatory:
• HTTPS only.
• JWT authentication.
• Role-based route protection.
• CORS restricted.
• Input validation enforced.
• Rate limiting enabled.
Admin routes isolated under:
/api/v1/admin
9.11 Background Jobs
Required scheduled jobs:
1. Media deletion (3-month retention)
2. Regional cluster recalculation (batch)
3. Trust score recalculation (optional batch)
4. Inactive crop archival (optional)
Can use:
• Cloud Scheduler
• Or background cron service
9.12 Scaling Strategy (Prototype Phase)
Expected usage:
• 100–200 farmers
• < 20 concurrent sessions

===== PAGE 110 =====
• < 50 video uploads/day
Scaling approach:
• Horizontal scaling of backend container.
• Separate video worker container.
• DB vertical scaling if required.
Event partitions prevent cross-crop blocking.
9.13 Logging & Monitoring
Enable:
• Structured application logs
• Error logs
• Replay failure logs
• Admin action logs
• ML inference logs
Monitoring tools:
• Cloud Logging
• Error alerting via email/webhook
9.14 Crash Recovery Strategy
If backend crashes:
1. Event log persists in DB.
2. On restart:
a. Resume unprocessed events.
3. Idempotent event handling prevents duplication.
If video worker crashes:
• Unprocessed media events remain in queue.
• Retry on restart.

===== PAGE 111 =====
9.15 Disaster Recovery
Database:
• Automated backups.
• Manual restore testing recommended.
Media storage:
• Object storage redundancy.
• Lifecycle enforcement.
9.16 Performance Constraints
Replay time target:
< 2 seconds per 500 actions.
Video processing:
Async, not blocking.
API response latency target:
< 300ms average.
9.17 Deployment Separation Summary
Component Deployment Unit
PWA Frontend Separate
Admin Frontend Separate
Backend API Separate
Video Worker Separate
Database Managed Service
Media Storage Cloud Object Storage
SECTION 9 — DEPLOYMENT &
INFRASTRUCTURE ENHANCEMENTS

===== PAGE 112 =====
ENHANCEMENT 1 — Environment
Separation Architecture
Add After: 9.1 Deployment Philosophy
Add subsection:
9.1.1 Environment Segmentation Policy
⚙ Define:
CultivaX must maintain:
1. Development Environment
2. Staging Environment
3. Production Environment
Rules:
• Separate databases
• Separate storage buckets
• Separate service accounts
• No shared secrets
• No cross-environment data mixing
Why:
Prevents accidental data corruption during testing.
ENHANCEMENT 2 — Infrastructure as Code
(IaC) Policy
Add Under: 9.4 Backend API Service
Add subsection:
9.4.1 Infrastructure Definition Strategy
Define:
• Use Terraform (preferred) or GCP Deployment Manager
• Define:
o Cloud Run service

===== PAGE 113 =====
o Cloud SQL
o Storage bucket
o Service accounts
o IAM roles
• No manual infrastructure provisioning in production
Why:
Prevents drift and misconfiguration.
ENHANCEMENT 3 — Replay Resource
Isolation
Add Under: 9.6 Event Dispatcher Infrastructure
Add:
Replay operations must:
• Run within controlled thread pool
• Enforce CPU usage cap
• Timeout threshold defined
• Not block health endpoint
Why:
Prevents replay from exhausting container resources.
ENHANCEMENT 4 — Resource Budget
Definition
Add Under: 9.12 Scaling Strategy
Add subsection:
9.12.1 Resource Budget Planning
Define expected usage:
Backend API:

===== PAGE 114 =====
• 1–2 vCPU
• 512MB–1GB RAM
Video Worker:
• 2–4 vCPU
• 2GB RAM
Database:
• Moderate IOPS
• 50k–100k action_logs capacity
Why:
Makes infrastructure realistic and defendable.
ENHANCEMENT 5 — Media Storage
Security Hardening
Strengthen: 9.8 Media Storage Infrastructure
Add:
• Bucket must be private
• Signed URLs valid < 15 minutes
• No public ACLs
• Lifecycle deletion policy enforced
• Audit logs enabled on bucket
Why:
Prevents accidental exposure of crop images.
ENHANCEMENT 6 — Network Isolation
Model
Add Under: 9.10 Security Infrastructure
Define:

===== PAGE 115 =====
• Cloud SQL must use private IP
• Backend connects via VPC connector
• Admin frontend restricted via firewall rule or IP allowlist
• No direct DB public exposure
Why:
Strengthens security posture significantly.
ENHANCEMENT 7 — Worker Isolation
Strategy
Strengthen: 9.7.2 Video Processing Worker
Add:
Video worker must:
• Run as separate Cloud Run service
• Use separate service account
• Use independent DB connection pool
• Use task queue trigger (Pub/Sub optional in future)
Why:
Prevents media overload affecting CTIS.
ENHANCEMENT 8 — Background Job
Isolation
Strengthen: 9.11 Background Jobs
Add:
Jobs must:
• Be idempotent
• Log execution status
• Run in separate container or scheduled worker

===== PAGE 116 =====
• Timeout protected
Jobs include:
• Media cleanup
• Cluster recalculation
• Trust recalculation
Why:
Prevents cron job failures from going unnoticed.
ENHANCEMENT 9 — Monitoring & Alert
Thresholds
Strengthen: 9.13 Logging & Monitoring
Add:
Define alert triggers:
• Replay failure rate > threshold
• ML service downtime > X minutes
• Video queue backlog > threshold
• DB connection errors spike
• Dead letter events present
Alerts must notify admin via email/webhook.
Why:
Operational maturity.
ENHANCEMENT 10 — Disaster Recovery
Drill Policy
Strengthen: 9.15 Disaster Recovery
Add subsection:
9.15.1 Recovery Validation Policy

===== PAGE 117 =====
Define:
• Quarterly restore test from backup
• Test replay from restored DB
• Verify event chain integrity
Why:
Backup without testing = illusion of safety.
ENHANCEMENT 11 — Horizontal Scaling
Constraint
Add Under: 9.12 Scaling Strategy
Define:
Cloud Run must:
• Allow min instances = 1
• Max instances = configurable
• Stateless containers only
• No in-memory state relied upon
Why:
Ensures container scaling doesn’t break event order.
ENHANCEMENT 12 — Feature Flag
Framework
Add Under: 9.4 Backend API Service
Define:
Feature flags stored in:
feature_flags table
Flags for:

===== PAGE 118 =====
• ML risk influence
• Regional clustering
• Edge AI integration
• SOE integration
Allows runtime toggling without redeploy.
Why:
Operational safety.
SECTION 10 — END-TO-END SYSTEM
FLOW & LIFECYCLE WALKTHROUGH
(Full System Operational Narrative)
10.1 Purpose
This section describes:
• How CultivaX operates from start to finish.
• How CTIS, SOE, ML, Event Dispatcher, and Accessibility interact.
• How offline usage is handled.
• How yield influences learning.
• How admin actions affect system.
• How failure conditions are handled.
This is the operational glue of the entire system.
10.2 Primary Lifecycle — Normal Online Flow
Step 1: Farmer Registration
1. Farmer registers.

===== PAGE 119 =====
2. Accessibility preferences stored.
3. JWT issued.
4. No CTIS involvement yet.
Step 2: Create Crop Instance
Farmer creates new crop.
Flow:
Frontend → API → EventDispatcher
→ CropCreated event
→ CTIS initializes instance
→ Baseline stage calculated
→ Seasonal window category assigned
→ State set to Active
Response sent to frontend.
Step 3: Timeline Viewing
Farmer opens crop dashboard.
API returns:
• current_stage
• baseline_stage
• stress_score
• risk_index
• recommended_action
• projected_harvest_window
No event triggered.
Read-only operation.
Step 4: Logging an Action
Example: Farmer irrigates.
Frontend → Log Action API
→ ActionLogged event created
→ EventDispatcher processes
→ CTIS stores action

===== PAGE 120 =====
→ ReplayTriggered event auto-generated
→ Replay engine recalculates timeline
→ Stress recalculated
→ Risk recalculated
→ State updated
Response returned to UI.
10.3 Offline Logging Flow
Scenario: Farmer offline for 5 days.
Step 1: Offline Action
Action stored locally with:
• action_effective_date
• local_monotonic_counter
No server communication.
Step 2: Sync on Connectivity
Device reconnects.
Batch upload:
• Unsynced actions sent.
Backend:
→ OfflineSyncCompleted event
→ Insert actions
→ Trigger ReplayTriggered
→ Full recomputation
→ Update current_stage
→ Update stress_score
Farmer sees updated timeline.
No data loss.
No timeline corruption.

===== PAGE 121 =====
10.4 Media Upload Flow
Case A: Online Image Upload
Farmer uploads image.
Flow:
→ MediaUploaded event
→ Backend ML processes image
→ MediaAnalyzed event
→ CTIS updates stress
→ Risk recalculated
Alert triggered if threshold crossed.
Case B: Offline + Edge AI
Farmer offline, low-data mode enabled.
Edge AI runs locally.
Structured features generated.
Stored locally.
On sync:
→ MediaUploaded
→ Backend re-evaluates
→ Backend result overrides Edge AI
→ MediaAnalyzed event
→ CTIS updates stress
Authority preserved.
10.5 Video Processing Flow
Farmer uploads video.

===== PAGE 122 =====
Flow:
→ MediaUploaded event
→ Video Worker picks job
→ Frame extraction
→ Frame inference
→ Aggregate results
→ MediaAnalyzed event
→ CTIS updates stress
UI shows:
“Processing in background…”
No blocking.
10.6 What-If Simulation Flow
Farmer simulates delay.
Flow:
→ WhatIf API
→ Clone CropInstance
→ Apply hypothetical action
→ Run replay in isolated memory
→ Return projected changes
No database mutation.
No event stored.
10.7 Yield Submission Flow
At harvest:
Farmer submits yield.
Flow:
→ YieldSubmitted event
→ Replay triggered
→ YieldVerificationScore computed
→ ml_yield_value capped if needed
→ RegionalClusterUpdated event (if eligible)
→ State set to Harvested

===== PAGE 123 =====
Yield displayed exactly as entered.
Regional update applies only to future crops.
10.8 Regional Learning Update Flow
When threshold met:
Batch job runs clustering.
Flow:
→ RegionalClusterUpdated event
→ Update cluster parameters
→ Store new regional offsets
Important:
No retroactive replay for existing crops.
Applies prospectively only.
10.9 Service Orchestration Flow
CTIS suggests irrigation.
Farmer selects “Find Equipment.”
Flow:
→ ServiceRequestCreated event
→ SOE returns verified providers
→ Farmer contacts provider
→ Completion recorded
→ Farmer submits review
→ TrustScoreUpdated event
CTIS unaffected.
10.10 Admin Rule Modification Flow
Admin modifies crop rule template.

===== PAGE 124 =====
Flow:
→ RuleModified event
→ New rule version created
→ Existing CropInstances:
• No automatic retroactive change
→ Future instances use new version
Optional:
Admin can trigger Replay for selected instance manually.
All actions logged in audit log.
10.11 Market Update Flow
Market API updates price.
Flow:
→ MarketUpdated event
→ Affects only harvest-stage projections
→ Recommendation messaging updated
Does not alter stage or stress.
10.12 Weather Update Flow
Weather API refreshes data.
Flow:
→ WeatherUpdated event
→ CTIS recalculates SystemRiskIndex
→ If threshold crossed:
• State may transition to AtRisk
Timeline stage unaffected.

===== PAGE 125 =====
10.13 Failure Scenarios
ML Failure
Fallback to rule-based logic.
System remains usable.
Replay Failure
Revert to last stable snapshot.
Admin notified.
Media Processing Failure
Stress not updated.
Log error.
Continue timeline.
WhatsApp Misformatted Input
Reject input.
Send structured instruction.
10.14 Crop Deletion Flow
Farmer deletes crop.
→ Soft delete (is_archived = true)
→ Timeline hidden
→ Action logs preserved
→ Extracted stats retained
No ML data removed.

===== PAGE 126 =====
10.15 Full Lifecycle Summary
Farmer Journey:
1. Create crop
2. Log actions
3. Upload media
4. Adapt timeline
5. Request services
6. Submit yield
7. Influence regional learning
8. Start next crop with adapted insights
System remains:
• Deterministic
• Transparent
• Replay-safe
• ML-bounded
• Governed
SECTION 10 — LIFECYCLE
WALKTHROUGH ENHANCEMENTS
ENHANCEMENT 1 — Extreme Offline
Scenario Walkthrough
Add After: 10.3 Offline Logging Flow
Add subsection:
10.3.1 Extended Offline Scenario (Long Duration)
⚙ Define Scenario:
Farmer:
• Offline for 30–45 days

===== PAGE 127 =====
• Logs 20 actions
• Uploads 5 images (Edge AI processed)
• Edits sowing date
• Submits yield immediately after reconnect
Flow must explicitly describe:
1. SyncBatch → insert actions
2. Validate timestamps
3. Clamp anomalies
4. Insert media features
5. Backend re-evaluates media
6. ReplayTriggered
7. Snapshot usage
8. YieldSubmitted
9. YieldVerificationScore applied
10. RegionalClusterUpdated (if eligible)
Explicitly state:
No data loss
No stage skipping
No cross-instance corruption
Why:
This is your hardest real-world scenario.
ENHANCEMENT 2 — Conflicting Signal
Resolution Scenario
Add After: 10.4 Media Upload Flow
Add subsection:
10.4.1 Multi-Signal Conflict Scenario
Scenario:
• Edge AI stress low
• Backend ML stress moderate
• Weather risk high
• Deviation pattern high
Walk through:

===== PAGE 128 =====
• Arbitration hierarchy
• Weighted stress integration
• Risk_index update
• AtRisk state trigger (if applicable)
• No stage skip
Explicitly reference authority order.
Why:
Proves ML boundedness.
ENHANCEMENT 3 — Admin Rule Change
Mid-Season
Strengthen: 10.10 Admin Rule Modification Flow
Add:
Scenario:
Admin modifies irrigation window mid-season.
System must:
• Create new rule version
• Existing CropInstances retain old rule_version_applied
• No automatic replay
• Admin may manually trigger replay
• All logged in audit log
Why:
This is a common viva trap question.

===== PAGE 129 =====
ENHANCEMENT 4 — Replay Failure
Recovery Scenario
Add After: 10.13 Failure Scenarios
Add subsection:
10.13.1 Replay Failure Handling Walkthrough
Scenario:
Replay fails due to unexpected data corruption.
System must:
1. Revert to last snapshot
2. Mark instance as RecoveryRequired
3. Lock further mutations
4. Log failure
5. Notify admin
6. Allow admin replay reattempt
Explicitly state:
System never leaves undefined state.
Why:
Shows resilience.
ENHANCEMENT 5 — Regional Cluster
Update Safety Walkthrough
Strengthen: 10.8 Regional Learning Update Flow
Add:
Scenario:
Cluster updated for Wheat in Haryana.
Ensure:
• Only new CropInstances use updated cluster

===== PAGE 130 =====
• No retroactive mutation
• No ReplayTriggered for old instances
• Cluster drift validated before commit
Why:
Prevents infinite replay loop concerns.
ENHANCEMENT 6 — Simultaneous Multi-
Crop Isolation Scenario
Add After: 10.15 Full Lifecycle Summary
Add subsection:
10.16 Multi-Crop Parallel Execution Scenario
Scenario:
Farmer manages:
• Wheat (AtRisk)
• Mustard (Healthy)
Simultaneously:
• Wheat triggers replay
• Mustard logs action
Explain:
• Partition isolation
• Independent event queues
• No cross-instance contamination
Why:
Validates partitioned dispatcher.

===== PAGE 131 =====
ENHANCEMENT 7 — Service Orchestration
Interaction Scenario
Strengthen: 10.9 Service Orchestration Flow
Add:
Scenario:
CTIS recommends irrigation.
Farmer requests equipment.
Provider fails to show.
Complaint filed.
TrustScoreUpdated event triggered.
Ensure:
• No CTIS mutation
• Timeline unchanged
• SOE state isolated
Why:
Proves pillar independence.
ENHANCEMENT 8 — What-If Simulation
Safety Scenario
Strengthen: 10.6 What-If Simulation Flow
Add:
Scenario:
Farmer runs What-If for 3-day irrigation delay.
Ensure:
• Deep clone created
• No event emitted
• No DB mutation
• Snapshot untouched

===== PAGE 132 =====
• Simulation disposed after response
Why:
Prevents hidden mutation bugs.
ENHANCEMENT 9 — ML Disabled Scenario
Add Under Failure Scenarios
Add:
If ML kill switch activated:
• RiskPrediction disabled
• Clustering disabled
• Edge AI ignored
• CTIS rule engine only active
System remains functional.
Why:
Shows bounded intelligence.
ENHANCEMENT 10 — Data Retention
Lifecycle Scenario
Add After Media Upload Flow
Add:
3-month media deletion walkthrough:
• scheduled_deletion_at reached
• Cron triggers deletion
• Raw file removed
• extracted_features retained
• Audit entry created
• Replay unaffected

===== PAGE 133 =====
Why:
Ensures deletion does not break determinism.
SECTION 11 — SECURITY & THREAT
MODEL
(Full System Security Architecture & Risk Analysis)
11.1 Purpose
This section defines:
• Security boundaries of CultivaX
• Threat identification
• Attack surfaces
• Mitigation strategies
• Data protection measures
• Role isolation
• ML data protection
• Infrastructure safeguards
The goal is to ensure:
• No timeline corruption
• No unauthorized access
• No trust score manipulation
• No data poisoning
• No replay abuse
• No admin misuse without trace
11.2 Security Philosophy
CultivaX follows:
1. Zero direct state mutation.

===== PAGE 134 =====
2. Event-driven controlled mutation.
3. Role-based access control.
4. Least privilege principle.
5. Soft-delete preservation.
6. Audit-first governance.
7. Deterministic replay safety.
11.3 Threat Surface Overview
Primary attack surfaces:
1. Authentication endpoints
2. Action logging API
3. Media upload API
4. WhatsApp webhook
5. Admin endpoints
6. Event dispatcher
7. ML training pipeline
8. Trust score system
9. Offline timestamp manipulation
Each must be secured.
11.4 Authentication & Authorization
11.4.1 Authentication
• JWT-based authentication.
• Short-lived access tokens.
• Refresh token mechanism.
• Password hashing using bcrypt or Argon2.
JWT must include:
• user_id
• role
• expiration
• issued_at

===== PAGE 135 =====
11.4.2 Role-Based Access Control (RBAC)
Roles:
• Farmer
• Service Provider
• Admin
Access Rules:
Endpoint Farmer Provider Admin
Crop APIs ✔ ✖ ✔
SOE APIs ✔ ✔ ✔
Admin APIs ✖ ✖ ✔
Role validated server-side only.
11.5 Timeline Integrity Threats
Threat 1: Offline Timestamp Manipulation
Attack:
User logs actions with future timestamps.
Mitigation:
• Compare action_effective_date with server_time.
• If impossible future date:
o Clamp to server_time.
o Log adjustment.
• Store device_timestamp for audit.
Replay always uses effective_date.
Threat 2: Action Replay Attack
Attack:
User re-sends same action multiple times.
Mitigation:

===== PAGE 136 =====
• Idempotency key required.
• action_id uniqueness enforced.
• Duplicate action rejected.
Threat 3: Direct DB Manipulation
Mitigation:
• No public DB exposure.
• No direct mutation endpoint.
• All changes through EventDispatcher.
11.6 ML Data Poisoning Threats
Threat 4: False Yield Inflation
Mitigation:
• Dual yield storage.
• ML uses capped ml_yield_value.
• Outliers excluded from clustering.
Farmer truth preserved.
ML protected.
Threat 5: Coordinated Regional Manipulation
Attack:
Multiple fake accounts submit inflated yields.
Mitigation:
• Minimum sample threshold.
• Verified yield confirmation logic.
• Admin monitoring of anomaly clusters.
• Optional anomaly detection on yield submissions.

===== PAGE 137 =====
11.7 Media Threats
Threat 6: Malicious File Upload
Mitigation:
• File type validation.
• File size limit.
• Virus scan (if feasible).
• Store in private bucket.
• Signed URL access only.
Threat 7: Media URL Exposure
Mitigation:
• No direct public URLs.
• Authenticated retrieval only.
• Time-limited signed links.
11.8 WhatsApp Threats
Threat 8: Webhook Spoofing
Mitigation:
• Validate Meta signature.
• Reject unsigned requests.
• Rate limit webhook.
Threat 9: Free-text Injection
Mitigation:
• Only structured responses accepted.

===== PAGE 138 =====
• Reject arbitrary text.
• Require numeric options.
11.9 Trust Score Manipulation
Threat 10: Self-rating
Mitigation:
• Must have valid ServiceRequest.
• Provider cannot review self.
• Only one review per request.
Threat 11: Review Deletion
Mitigation:
• Reviews immutable.
• Only admin can flag.
• Deletion logged.
11.10 Admin Threat Model
Admin has high authority.
Threat:
Admin misuse.
Mitigation:
• All admin actions logged in AdminAuditLog.
• Before & after values stored.
• No silent rule modification.
• Optional dual-admin approval (future enhancement).

===== PAGE 139 =====
11.11 Event Dispatcher Threats
Threat 12: Event Order Corruption
Mitigation:
• Partitioned FIFO queue.
• Lock per crop_instance.
• Idempotent event execution.
• Persisted event_log.
Threat 13: Infinite Replay Loop
Mitigation:
• Regional updates prospective only.
• Replay not triggered by cluster update for past instances.
11.12 Data Retention & Privacy Threats
Threat 14: Personal Data Leakage
Mitigation:
• No PII used in ML training.
• Raw media deleted after 3 months.
• Extracted features anonymized.
Threat 15: Unauthorized Media Access
Mitigation:
• Private storage.
• Signed URLs.
• Auth validation before access.

===== PAGE 140 =====
11.13 API Security
• HTTPS mandatory.
• Input validation (Pydantic or equivalent).
• SQL injection protection via ORM.
• Rate limiting on sensitive endpoints.
• CSRF protection for web forms.
11.14 Infrastructure Security
• Cloud SQL private IP.
• Service account isolation.
• Separate staging & production environments.
• Secrets stored in secret manager.
• No credentials hardcoded.
11.15 Logging & Monitoring Security
• Log failed login attempts.
• Log abnormal yield submissions.
• Log repeated failed media uploads.
• Alert on abnormal event failure patterns.
11.16 Denial of Service Mitigation
• Rate limit API endpoints.
• Limit media upload size.
• Limit video duration.
• Limit action logging frequency.
11.17 Disaster Recovery
• Daily DB backups.

===== PAGE 141 =====
• Media storage redundancy.
• Event log persistence.
11.18 Security Guarantees Summary
CultivaX ensures:
✔ Chronological integrity protection
✔ ML poisoning resistance
✔ Trust manipulation prevention
✔ Media security
✔ Admin accountability
✔ Event ordering safety
✔ Offline tampering protection
✔ Privacy compliance
✔ Role-based isolation
SECTION 11 — SECURITY & THREAT
MODEL ENHANCEMENTS
ENHANCEMENT 1 — Formal Threat
Categorization Framework
Add After: 11.3 Threat Surface Overview
Add subsection:
11.3.1 Threat Classification Matrix
Define threats by category:
Category Example
Data Integrity Timestamp manipulation
Event Integrity Replay duplication
ML Integrity Yield poisoning
Authorization Role escalation

===== PAGE 142 =====
Infrastructure Container compromise
Marketplace Abuse Fake reviews
Offline Abuse Backdated mass logging
This formalizes your threat model academically.
ENHANCEMENT 2 — Replay Tamper
Detection
Strengthen: 11.12 Data Retention & Privacy Threats
Add subsection:
11.12.1 Replay Integrity Verification
Define:
• Event chain hash verification on replay
• Snapshot integrity checksum
• If mismatch:
o Lock crop_instance
o Flag admin
o Prevent further mutation
Why:
Detects silent event corruption.
ENHANCEMENT 3 — Privilege Escalation
Guard
Strengthen: 11.4 Role-Based Access Control
Add:
All admin endpoints must:
• Verify role server-side
• Require short-lived token
• Optionally require step-up verification (OTP for destructive actions)

===== PAGE 143 =====
Add:
Admin deletion of crop/service must log:
• previous_state
• deleted_by
• reason
Why:
Prevents accidental or malicious misuse.
ENHANCEMENT 4 — ML Poisoning Rate
Limiter
Strengthen: 11.6 ML Data Poisoning Threats
Add:
Define:
If user submits:
• Excessive yield edits
• Extreme outlier values
• Frequent deletions and re-submissions
Then:
• Flag anomaly_score
• Temporarily exclude user from ML training pool
• Log event
Why:
Protects ML training integrity.

===== PAGE 144 =====
ENHANCEMENT 5 — Offline Abuse
Detection
Strengthen: 11.5 Timeline Integrity Threats
Add subsection:
11.5.3 Offline Pattern Abuse Detection
Define detection for:
• Large backdated batch uploads
• Non-monotonic local counter resets
• Excessive timestamp clustering
Flag if:
pattern_anomaly_score > threshold.
Why:
Offline flexibility must not enable silent manipulation.
ENHANCEMENT 6 — Media Abuse
Safeguard
Strengthen: 11.7 Media Threats
Add:
• Max image count per day per crop
• Max video length limit
• Reject executable content disguised as media
• Virus scanning (if feasible)
Why:
Prevents media overload or exploit vectors.

===== PAGE 145 =====
ENHANCEMENT 7 — API Abuse Mitigation
Layer
Strengthen: 11.13 API Security
Add:
• IP throttling for repeated failed logins
• CAPTCHA for repeated suspicious attempts (optional)
• Reject malformed JSON early
Why:
Prevents automated abuse.
ENHANCEMENT 8 — Service Marketplace
Fraud Safeguard
Strengthen: 11.9 Trust Score Manipulation
Add:
Fraud detection if:
• Multiple reviews from related accounts
• Rapid trust score spike
• Self-review detection
• IP clustering detection
Flag for admin review.
Why:
Marketplace integrity critical.

===== PAGE 146 =====
ENHANCEMENT 9 — Infrastructure
Compromise Containment
Strengthen: 11.14 Infrastructure Security
Add:
If backend container compromised:
• No local file persistence
• Stateless design ensures minimal data exposure
• DB credentials stored in Secret Manager only
• Service accounts least privilege
Why:
Limits blast radius.
ENHANCEMENT 10 — Audit Log
Immutability Rule
Strengthen: 11.15 Logging & Monitoring Security
Define:
• AdminAuditLog must be append-only
• No hard deletion allowed
• Only soft flag for correction
• Hash-based log verification (optional advanced)
Why:
Prevents silent tampering.

===== PAGE 147 =====
ENHANCEMENT 11 — Data Minimization
Enforcement
Strengthen: 11.12 Privacy Threats
Add:
Explicitly state:
• No biometric data
• No exact geolocation storage beyond required precision
• No PII used in ML training
• Extracted features anonymized
Why:
Aligns with ethical stance.
ENHANCEMENT 12 — ML Kill Switch
Security Binding
Strengthen: 11.16 Denial of Service Mitigation
Add:
If anomaly spike detected:
Admin may:
• Disable ML risk influence
• Disable clustering
• Disable media-triggered stress
Without affecting deterministic CTIS.
Why:
Operational containment.

===== PAGE 148 =====
SECTION 12 — TESTING & VALIDATION
STRATEGY
(Full System Validation Framework)
12.1 Purpose
This section defines:
• How each subsystem is tested.
• How determinism is verified.
• How replay integrity is validated.
• How ML boundaries are enforced.
• How security threats are simulated.
• How accessibility consistency is verified.
Testing must ensure:
• No timeline corruption.
• No cross-crop contamination.
• No race conditions.
• No silent state mutation.
• No ML authority overreach.
12.2 Testing Layers Overview
Testing is divided into:
1. Unit Testing
2. Integration Testing
3. Replay Determinism Testing
4. Event Ordering Testing
5. ML Validation Testing
6. Media Processing Testing
7. Security Testing
8. Accessibility Testing
9. Performance Testing
10. User Acceptance Testing (Prototype Phase)

===== PAGE 149 =====
12.3 Unit Testing
Scope:
• CTIS stage transitions
• Stress smoothing formula
• Risk index computation
• Trust score calculation
• Yield capping logic
• Seasonal window classification
Example:
Test:
Given stress_score=0.4 and new_signal=0.8
Ensure smoothing produces correct new value.
Test:
Yield > biological_limit
Ensure ml_yield_value capped but reported_yield unchanged.
Unit tests must cover:
• All state transitions.
• All mathematical models.
• All threshold conditions.
12.4 Integration Testing
Test interactions between modules:
• ActionLogged → ReplayTriggered
• MediaAnalyzed → StressUpdated
• YieldSubmitted → RegionalClusterUpdated
• RuleModified → New version applied
Ensure:
• No direct DB mutation bypassing event dispatcher.
• Proper event chaining.
• No circular event loop.

===== PAGE 150 =====
12.5 Replay Determinism Testing
Critical test category.
Procedure:
1. Create crop instance.
2. Log sequence of actions.
3. Save final state.
4. Trigger replay manually.
5. Compare final state.
Expected:
States must match exactly.
Also test:
• Offline sync scenario.
• Sowing date modification.
• Rule version update.
Replay must be deterministic.
12.6 Event Ordering Testing
Test scenario:
• Multiple actions logged rapidly.
• Ensure FIFO per crop_instance.
Simulate:
Concurrent events across different crops.
Expected:
• Correct isolation.
• No cross-instance interference.
Test replay lock:
Replay must block subsequent events until completion.

===== PAGE 151 =====
12.7 ML Validation Testing
Test cases:
• High stress scenario → risk threshold crossed.
• Edge AI vs backend conflict → backend overrides.
• Yield inflation attempt → capped for ML.
• Regional cluster update threshold not met → no update.
Ensure:
ML does not:
• Skip stage.
• Directly modify sowing date.
• Delete logs.
12.8 Media Processing Testing
Test image upload:
• Valid image.
• Invalid format.
• Oversized file.
• Corrupted image.
Test video upload:
• Short video.
• Max-length video.
• Frame extraction success.
• Frame extraction failure.
Ensure:
• Async processing works.
• No timeline corruption if processing fails.
• Raw media deletion after 3 months.
12.9 Security Testing
Simulated attacks:

===== PAGE 152 =====
1. Timestamp manipulation.
2. Duplicate action submission.
3. Invalid JWT.
4. Admin endpoint access without role.
5. Webhook spoofing.
6. SQL injection attempt.
7. Excessive action logging spam.
Expected:
All rejected.
No state mutation.
12.10 Accessibility Testing
Test modes:
• High contrast
• Lite mode
• Low literacy mode
• WhatsApp mode
• Voice assist
Ensure:
• All actions still validated.
• No bypass of confirmation.
• Timeline integrity maintained.
12.11 Performance Testing
Replay stress test:
Simulate:
500 actions per crop.
Measure:
Replay duration < 2 seconds.
Concurrent crop test:
Simulate:
10 crops replaying simultaneously.

===== PAGE 153 =====
Ensure:
No deadlock.
No starvation.
Media processing test:
Simulate:
Multiple video uploads.
Ensure:
Main API latency unaffected.
12.12 Database Integrity Testing
Test:
• Foreign key enforcement.
• Soft delete behavior.
• Orphan prevention.
• Index effectiveness.
Run:
• EXPLAIN queries.
• Load tests.
12.13 Failure Injection Testing
Simulate:
• ML service down.
• Weather API down.
• Market API down.
• Event processing crash.
• DB restart.
Expected:
System degrades gracefully.
No timeline corruption.
Events retried safely.

===== PAGE 154 =====
12.14 End-to-End Scenario Testing
Simulate full lifecycle:
1. Create crop.
2. Log offline actions.
3. Upload media.
4. Delay irrigation.
5. Submit yield.
6. Influence regional learning.
7. Create next crop.
Verify:
System stable.
No retroactive mutation.
Regional update prospective only.
12.15 Regression Testing Strategy
Every time:
• CropRuleTemplate modified
• ML model updated
• Event dispatcher updated
Must rerun:
• Replay determinism tests.
• Stage transition tests.
• Risk threshold tests.
12.16 Logging Validation
Ensure:
• All admin actions logged.
• All replay triggers logged.
• All yield submissions logged.
• All failure events logged.
No silent failure allowed.

===== PAGE 155 =====
12.17 Acceptance Criteria for Prototype
Prototype considered valid when:
✔ Replay deterministic
✔ Offline sync stable
✔ Media pipeline functional
✔ Trust score consistent
✔ Admin governance logged
✔ No direct state mutation
✔ ML bounded and safe
✔ Accessibility modes working
✔ Performance within limits
SECTION 12 — TESTING & VALIDATION
ENHANCEMENTS
ENHANCEMENT 1 — Property-Based
Replay Testing
Add After: 12.5 Replay Determinism Testing
Add subsection:
12.5.1 Property-Based Replay Verification
⚙ Define:
Instead of testing only fixed scenarios:
Generate:
• Random valid action sequences
• Random offline batches

===== PAGE 156 =====
• Random action ordering within constraints
For each sequence:
1. Execute normally
2. Save final state
3. Trigger replay
4. Assert equality of:
a. stage
b. stress_score
c. risk_index
d. deviation_profile
Invariant:
Replay(original_sequence) == CurrentState
Why:
This mathematically proves determinism.
ENHANCEMENT 2 — Invariant Testing
Framework
Add After: 12.3 Unit Testing
Add subsection:
12.3.1 Chronological Invariant Tests
Test:
• effective_date ≥ sowing_date
• No stage skipping
• No negative stress
• No invalid state transitions
• No mutation without event
If any invariant fails → test fails.
Why:
Validates Section 1 invariants explicitly.

===== PAGE 157 =====
ENHANCEMENT 3 — ML Boundary
Condition Testing
Strengthen: 12.7 ML Validation Testing
Add:
Test scenarios:
• Extremely high stress image
• Extremely low stress image
• Extreme yield value
• Minimal dataset cluster
• Conflicting weather + image
Verify:
• No stage mutation
• No baseline rewrite
• No negative values
• ML confidence influences risk proportionally
Why:
Ensures ML boundedness is enforced in practice.
ENHANCEMENT 4 — Chaos Testing
Scenarios
Add After: 12.13 Failure Injection Testing
Add subsection:
12.13.1 Controlled Chaos Simulation
Simulate:
• DB restart during replay
• Event worker crash mid-processing
• Media worker crash mid-video
• ML service timeout
• Network latency spike

===== PAGE 158 =====
Expected:
• No silent corruption
• Retry mechanism works
• Dead letter events isolated
• Replay can resume safely
Why:
Proves resilience.
ENHANCEMENT 5 — Multi-Crop
Concurrency Stress Test
Add After: 12.11 Performance Testing
Simulate:
• 20 crops
• Concurrent action logging
• Simultaneous replay
• Video processing in background
Verify:
• No cross-crop contamination
• Partition isolation intact
• No deadlocks
Why:
Validates partitioned dispatcher.
ENHANCEMENT 6 — Accessibility Matrix
Testing
Strengthen: 12.10 Accessibility Testing
Define matrix:

===== PAGE 159 =====
Each feature tested under:
• Online
• Offline
• Lite Mode
• High Contrast
• WhatsApp
• Voice
Test:
• Action logging integrity
• Replay stability
• No duplicate submissions
• No state corruption
Why:
Accessibility must not break determinism.
ENHANCEMENT 7 — Offline Abuse
Simulation
Add Under: Security Testing
Simulate:
• 100 backdated actions
• Non-monotonic counters
• Repeated batch uploads
• Duplicate Idempotency-Key
Expected:
• Anomaly detection triggered
• Duplicate rejected
• Replay safe
Why:
Validates abuse resistance.

===== PAGE 160 =====
ENHANCEMENT 8 — Statistical Stability
Testing
Add Under: 12.7 ML Validation Testing
Simulate:
• Small cluster size
• Cluster variance spike
• Outlier yield submissions
Verify:
• Cluster update blocked if unstable
• No retroactive mutation
• Yield capping applied
Why:
Prevents unstable regional learning.
ENHANCEMENT 9 — Media Quality
Degradation Test
Add Under: 12.8 Media Processing Testing
Test:
• Blurry image
• Dark image
• Overexposed image
• Duplicate frame video
Verify:
• Confidence_score lowered
• Stress update weighted appropriately
• No stage shift

===== PAGE 161 =====
Why:
Validates noise filtering logic.
ENHANCEMENT 10 — Snapshot Integrity
Test
Add Under: 12.4 Database Integrity Testing
Test:
• Replay from snapshot
• Corrupt snapshot simulation
• Snapshot rollback
Verify:
• Recovery to last valid snapshot
• No partial state application
Why:
Snapshot strategy must be validated.
ENHANCEMENT 11 — API Contract
Stability Testing
Add After: 12.1 Overview
Define:
• Backward compatibility test suite
• Deprecated endpoint response validation
• Schema version mismatch test
Why:
Prevents contract breakage.

===== PAGE 162 =====
ENHANCEMENT 12 — Load Testing Replay
Scaling Curve
Add Under: 12.11 Performance Testing
Simulate:
• Increasing action count per crop:
10 → 50 → 200 → 500
Measure:
• Replay duration
• Memory usage
• CPU spikes
Plot scaling curve.
Why:
Proves replay remains linear.
SECTION 13 — DEVELOPMENT
ROADMAP & IMPLEMENTATION
MILESTONES
(Execution Blueprint for CultivaX Prototype)
13.1 Development Philosophy
CultivaX must be built in controlled layers, not feature chaos.
Order of development priority:
1. Chronological integrity (CTIS core)

===== PAGE 163 =====
2. Event dispatcher
3. Database stability
4. Basic UI + action logging
5. Replay engine
6. Stress & rule engine
7. Media pipeline (basic)
8. SOE basic matching
9. Accessibility layer
10. ML statistical refinement
11. Regional clustering
12. Advanced media & edge AI
We do not start with ML.
We start with timeline determinism.
13.2 Phase 0 — Architecture Finalization (Pre-
Development)
Deliverables:
• Final system document (this document).
• Database schema finalized.
• Event model finalized.
• Role definitions finalized.
• Deployment stack chosen (GCP + FastAPI + PostgreSQL).
No coding until this is stable.
13.3 Phase 1 — Core CTIS Foundation (Critical)
Goal:
Deterministic crop lifecycle engine working.
Build:
• Farmer authentication
• Crop instance creation
• CropRuleTemplate versioning
• Action logging
• Replay engine
• Stage transitions

===== PAGE 164 =====
• Stress smoothing formula
• Seasonal risk logic
• Yield submission (without ML clustering)
Must Validate:
• Replay determinism.
• Offline simulation (manual test).
• No direct DB mutation.
• State machine correctness.
At end of Phase 1:
System must handle:
Full crop lifecycle without media, ML, or services.
13.4 Phase 2 — Event Dispatcher & Offline Sync
Goal:
Stabilize event-driven mutation.
Build:
• Partitioned event queue.
• Idempotent event handling.
• Offline action sync.
• Replay lock mechanism.
• Event log persistence.
• Crash recovery logic.
Must Validate:
• FIFO per crop_instance.
• Replay atomicity.
• No cross-crop interference.
At end of Phase 2:
System fully deterministic under offline conditions.

===== PAGE 165 =====
13.5 Phase 3 — Media Pipeline (Basic Backend
Only)
Goal:
Enable image-based stress update.
Build:
• Media upload endpoint.
• Basic image classifier (backend).
• Structured feature storage.
• MediaAnalyzed event integration.
• Stress update logic.
Do NOT build yet:
• Video processing.
• Edge AI.
At end of Phase 3:
Stress updates through image upload must function.
13.6 Phase 4 — Service Orchestration (Basic)
Goal:
Marketplace-style resource discovery.
Build:
• Provider registration.
• Provider verification (admin).
• Equipment listing.
• Service request creation.
• Review submission.
• Trust score calculation.
Must Validate:
• Trust score correct.
• One review per request.
• Admin override logged.

===== PAGE 166 =====
At end of Phase 4:
Basic matching ecosystem functional.
13.7 Phase 5 — Accessibility Layer
Goal:
Make system rural-usable.
Implement:
• High contrast mode.
• Lite mode.
• Multi-language support.
• Structured WhatsApp integration.
• Typography scaling.
Validate:
• No logic corruption.
• All modes pass replay & action tests.
13.8 Phase 6 — ML Risk Model (Statistical Layer)
Goal:
Introduce bounded statistical prediction.
Implement:
• Logistic regression model.
• Risk probability integration.
• Yield verification score.
• Dual yield storage.
Validate:
• ML cannot skip stages.
• ML cannot modify sowing date.
• ML cannot delete logs.
• Outlier yield capped for ML only.

===== PAGE 167 =====
13.9 Phase 7 — Regional Clustering
Goal:
Adaptive regional intelligence.
Implement:
• Cluster key definition.
• Sample threshold enforcement.
• Prospective-only application.
• Batch update job.
Validate:
• No retroactive timeline mutation.
• No infinite replay loop.
• Sample threshold respected.
13.10 Phase 8 — Advanced Media & Edge AI
(Optional Advanced Feature)
Goal:
Demonstrate advanced architecture.
Implement:
• Edge AI model (MobileNet-based).
• Offline image inference.
• Video frame processing worker.
• Aggregation logic.
• Async pipeline stability.
This phase strengthens defense of project but is not required for MVP.
13.11 Phase 9 — Hardening & Security Testing
• Rate limiting.
• JWT validation.
• Webhook signature validation.
• Media retention enforcement.

===== PAGE 168 =====
• Admin audit validation.
• Replay stress testing.
13.12 Phase 10 — Deployment & Final Validation
Deploy:
• Frontend (PWA)
• Backend API
• Database
• Media storage
• Worker service
Conduct:
• Full end-to-end testing.
• Performance validation.
• Failure injection testing.
13.13 Estimated Implementation Order Timeline
(Academic Prototype)
Phase Complexity Priority
Phase 1 High Critical
Phase 2 High Critical
Phase 3 Medium High
Phase 4 Medium High
Phase 5 Medium Medium
Phase 6 High Medium
Phase 7 High Advanced
Phase 8 Very High Advanced
For lab submission:
Phases 1–5 are sufficient.
Phases 6–8 elevate to research-grade prototype.

===== PAGE 169 =====
13.14 MVP Definition (Minimum Viable
Prototype)
MVP must include:
✔ CTIS deterministic
✔ Replay engine
✔ Offline logging
✔ Basic media stress
✔ Basic service matching
✔ Accessibility baseline
✔ Admin governance
Edge AI, clustering, advanced ML optional for MVP.
13.15 Risk Control Strategy During Development
Never:
• Build clustering before replay stable.
• Build edge AI before backend pipeline stable.
• Modify rule templates without versioning.
• Add ML without yield cap logic.
Chronology first. Intelligence second.
SECTION 13 — DEVELOPMENT
ROADMAP ENHANCEMENTS

===== PAGE 170 =====
ENHANCEMENT 1 — Phase Gate Criteria
(Formal Exit Conditions)
Add After: Each Phase Description
Add subsection:
Phase Exit Validation Criteria
For each phase define:
Phase 1 Exit (CTIS Core):
• Replay deterministic
• Chronological invariants enforced
• No state mutation outside event dispatcher
• Snapshot creation working
Phase 2 Exit (Dispatcher):
• Idempotency enforced
• Partition isolation verified
• Dead letter queue tested
• Replay atomic guard validated
… and so on.
Why:
Prevents moving forward with unstable base.
ENHANCEMENT 2 — Research Feature
Isolation Layer
Add After: Phase 8
Add subsection:
Advanced Research Layer (Optional for MVP)
Define research-only components:
• Edge AI
• Long-horizon forecasting

===== PAGE 171 =====
• Regional clustering adaptation
• Confidence propagation dashboard
• What-if sandbox visualization
Clearly mark:
Not required for lab MVP
Required for research-paper grade
Why:
Prevents evaluation confusion.
ENHANCEMENT 3 — Technical Debt
Tracking Policy
Add After: 13.15 Risk Control Strategy
Add subsection:
Technical Debt Register
Define:
Every shortcut must:
• Be documented
• Have severity level
• Have mitigation plan
• Be linked to phase
Example:
Shortcut:
No external message broker (using in-process dispatcher)
Mitigation:
Replace with Pub/Sub if scale > 500 users
Why:
Shows maturity and long-term awareness.

===== PAGE 172 =====
ENHANCEMENT 4 — Data Acquisition Plan
for ML
Add After: Phase 6
Add subsection:
Data Collection & Validation Plan
Define:
Sources:
• Public agricultural datasets
• Farmer-uploaded media (filtered)
• Manual agronomic baseline data
• Historical mandi price APIs
Before ML deployment:
• Data quality threshold
• Dataset audit
• Bias evaluation
Why:
ML without dataset strategy is weak academically.
ENHANCEMENT 5 — Risk Mitigation per
Phase
Add After Each Phase
Define per-phase primary risk and mitigation.
Example:
Phase 3 (Media):
Risk:
Poor image accuracy.
Mitigation:
Confidence-weighted stress + backend authority.

===== PAGE 173 =====
Phase 6 (ML):
Risk:
Overfitting.
Mitigation:
Cross-validation + drift detection.
Why:
Shows execution awareness.
ENHANCEMENT 6 — Feature Flag
Integration Plan
Add Under: Phase 9 Hardening
Define:
Each major feature behind feature flag:
• ML risk
• Clustering
• Edge AI
• SOE integration
Allows progressive rollout.
Why:
Supports safe deployment.
ENHANCEMENT 7 — Scalability Migration
Plan
Add After: 13.13 Timeline Table
Add subsection:
Post-Prototype Scalability Migration Plan
Define future steps:

===== PAGE 174 =====
1. Replace in-memory dispatcher → Pub/Sub or Kafka
2. Extract CTIS into dedicated service
3. Extract Media worker into auto-scaling cluster
4. Introduce read replica DB
5. Add Redis caching for read-heavy endpoints
Why:
Shows architectural foresight.
ENHANCEMENT 8 — Risk Containment
Policy
Add Under: 13.15 Risk Control Strategy
Define:
If a major subsystem unstable:
Disable subsystem via feature flag
Continue deterministic CTIS core
System must degrade gracefully.
Why:
Protects core functionality.
ENHANCEMENT 9 — UX Freeze Policy
Add After Accessibility Phase
Define:
After Phase 5:
• Freeze major UI redesign
• Only bug fixes allowed
• Prevent scope creep

===== PAGE 175 =====
Why:
UI churn delays backend stabilization.
ENHANCEMENT 10 — Milestone-Based
Validation Demo Plan
Add Final Subsection
Define:
Internal milestone demos:
Demo 1:
Replay determinism
Demo 2:
Offline batch sync
Demo 3:
Media stress update
Demo 4:
Service orchestration
Demo 5:
ML bounded risk
This becomes your evaluation defense plan.
SECTION 14 — RISK & LIMITATION
ANALYSIS
(Architectural Trade-offs, Constraints & Known Boundaries)

===== PAGE 176 =====
14.1 Purpose
This section documents:
• Known technical risks
• System design trade-offs
• Prototype limitations
• ML limitations
• Operational constraints
• Ethical considerations
• Future scalability risks
This ensures transparency and architectural maturity.
14.2 Architectural Trade-Offs
14.2.1 Modular Monolith vs Microservices
Current design:
Modular Monolith with logical separation.
Trade-off:
✔ Easier development
✔ Deterministic event control
✔ Lower operational complexity
But:
✖ Not horizontally scalable at massive scale
✖ Single service bottleneck at high load
Mitigation:
Future microservice extraction possible:
• Separate CTIS service
• Separate SOE service
• Separate ML worker cluster

===== PAGE 177 =====
14.2.2 In-Memory Event Queue (Prototype)
Risk:
If using in-memory queue without Redis/Kafka:
• Event loss possible during crash (if not persisted first)
Mitigation:
Event must be written to event_log table before processing.
Long-term:
External message broker recommended.
14.3 Replay Engine Risks
Replay complexity increases with:
• High action volume
• Frequent rule updates
• Long crop lifecycles
Risk:
Replay latency may increase under scale.
Mitigation:
• Enforce O(n) replay algorithm.
• Limit excessive action logging.
• Optimize indexing.
• Consider incremental replay strategy in future.
14.4 ML Limitations
14.4.1 Logistic Regression Simplicity
Current model:
Logistic Regression for risk prediction.

===== PAGE 178 =====
Limitation:
• May not capture complex nonlinear crop stress patterns.
• May underperform under diverse environmental variation.
Mitigation:
Future upgrade path:
• Gradient Boosted Trees
• Neural Networks
• Hybrid agronomic-ML model
14.4.2 Regional Clustering Sensitivity
Risk:
Small sample sizes may produce unstable clusters.
Mitigation:
• Minimum threshold enforcement.
• Prospective-only application.
• Manual admin override.
14.4.3 Edge AI Accuracy
Edge AI uses lightweight model.
Limitation:
• Lower accuracy compared to backend CNN.
• May misclassify complex disease patterns.
Mitigation:
Backend re-evaluation always authoritative.
14.5 Media Processing Limitations

===== PAGE 179 =====
14.5.1 Video Complexity
Limitations:
• Cannot guarantee precise disease identification.
• Lighting conditions affect accuracy.
• Multiple plant overlap may reduce detection reliability.
System must clearly state:
“Stress probability” not “disease diagnosis”.
14.6 Data Limitations
Prototype scale:
100–200 users.
Limitations:
• Regional clustering accuracy limited by sample size.
• Statistical patterns may be weak initially.
Mitigation:
• Gradual expansion.
• Controlled dataset augmentation.
• Public dataset integration for model pretraining.
14.7 Offline Mode Risks
Risk:
User may manipulate timestamps.
Mitigation:
• Clamp extreme values.
• Audit device_timestamp vs server_timestamp.
• Maintain monotonic counter.
Still:
Complete prevention of offline falsification impossible.

===== PAGE 180 =====
14.8 Trust Score System Limitations
Limitations:
• Cannot guarantee service quality.
• May be biased if review sample small.
• Subject to coordinated false reviews (mitigated but not impossible).
Future improvement:
• Weight reviews by verified transaction quality.
• Introduce anomaly detection.
14.9 Accessibility Trade-Offs
Lite mode reduces:
• Visual richness
• Data visualization depth
Trade-off:
Performance over visual detail.
Acceptable for rural deployment.
14.10 Infrastructure Limitations
Current deployment assumes:
• Stable cloud environment.
• Moderate traffic.
• Reliable storage.
Risks:
• Single-region hosting.
• No active-active redundancy.
Future solution:

===== PAGE 181 =====
• Multi-region replication.
• Managed distributed queue.
14.11 Security Limitations
Even with safeguards:
• Social engineering attacks possible.
• Shared device usage risks.
• Phone-based authentication may be compromised.
Mitigation:
• OTP-based login.
• Optional 2FA for admin.
• Session expiration enforcement.
14.12 Ethical Considerations
CultivaX must not:
• Present predictions as guarantees.
• Replace farmer judgment.
• Bias recommendations toward specific providers.
• Store identifiable training data.
• Exploit user data.
Clear disclaimers required:
• ML predictions probabilistic.
• Market forecast indicative only.
• Stress detection not medical-grade diagnosis.
14.13 Legal & Compliance Considerations
(Prototype Phase)
• No payment handling reduces liability.
• No financial scoring (Zero-Knowledge Ledger moved to future).
• No biometric storage.

===== PAGE 182 =====
• No medical claims.
14.14 Scalability Risk Summary
If user base increases to 10,000+:
Risks:
• Replay latency increases.
• DB performance pressure.
• Media storage scaling.
• Event queue congestion.
Future required:
• Event broker (Kafka/PubSub).
• Horizontal service separation.
• Read replica DB.
• Distributed worker pool.
14.15 Overall System Risk Profile
Low Risk:
• Timeline corruption (due to replay design)
• Trust manipulation (controlled)
• ML authority overreach (bounded)
Moderate Risk:
• Small-sample clustering instability
• Offline timestamp manipulation attempts
Higher Risk (Long-Term Scale):
• Infrastructure bottlenecks
• Event queue scaling
• Video processing load
All acceptable for prototype scale.

===== PAGE 183 =====
SECTION 14 — RISK & LIMITATION
ENHANCEMENTS
ENHANCEMENT 1 — Formal Risk
Classification Model
Add After: 14.1 Purpose
Add subsection:
14.1.1 Risk Classification Framework
Classify risks into:
Category Description
Architectural Risk Design limitations
Operational Risk Deployment/runtime limitations
Statistical Risk ML uncertainty
Human Behavioral Risk User misuse
Ethical Risk Over-dependence or misinterpretation
Scalability Risk Future growth constraints
This gives structure to the entire section.
ENHANCEMENT 2 — Determinism Trade-
Off Acknowledgment
Add Under: 14.2 Architectural Trade-Offs
Add subsection:
14.2.3 Deterministic Replay vs Real-Time Responsiveness
Explain:
Because system prioritizes replay determinism:
• Slight delay may occur during replay
• Atomic locking may block concurrent updates
• Real-time streaming feedback limited

===== PAGE 184 =====
Mitigation:
Replay optimized with snapshot strategy.
Why:
Shows conscious architectural trade-off.
ENHANCEMENT 3 — ML Overconfidence
Risk
Add Under: 14.4 ML Limitations
Add subsection:
14.4.4 Model Confidence Misinterpretation Risk
Risk:
Farmer may interpret probability as certainty.
Mitigation:
• Confidence score display
• Advisory labeling
• No deterministic claims
Why:
Aligns with your advisory philosophy.
ENHANCEMENT 4 — Small Dataset
Regional Bias
Strengthen: 14.6 Data Limitations
Add:
Regional clustering may:
• Overfit small geography
• Reflect sampling bias

===== PAGE 185 =====
• Reinforce incorrect patterns
Mitigation:
• Minimum sample threshold
• Drift validation
• Prospective-only application
ENHANCEMENT 5 — Edge AI
Environmental Bias
Add Under: 14.5 Media Processing Limitations
Add:
Edge AI may:
• Underperform in extreme sunlight
• Misinterpret soil color as leaf stress
• Be biased by camera quality
Mitigation:
Backend re-evaluation authoritative.
ENHANCEMENT 6 — Marketplace Network
Effect Risk
Add Under: 14.8 Trust Score System Limitations
Risk:
Early providers may dominate permanently.
Mitigation:
• Trust decay
• Exposure fairness model
• Randomized exposure factor

===== PAGE 186 =====
ENHANCEMENT 7 — Offline Autonomy
Risk
Add Under: 14.7 Offline Mode Risks
Add:
Offline freedom may:
• Delay stress detection
• Reduce timely recommendations
• Allow pattern manipulation
Mitigation:
• Sync anomaly detection
• Timestamp clamping
• Risk recalculation post-sync
ENHANCEMENT 8 — Infrastructure Single-
Region Limitation
Strengthen: 14.10 Infrastructure Limitations
Add:
Single-region deployment may:
• Suffer regional outage
• Increase latency for distant users
Future mitigation:
• Multi-region deployment
• Read replicas

===== PAGE 187 =====
ENHANCEMENT 9 — Replay Scalability
Ceiling
Add Under: 14.14 Scalability Risk Summary
Define:
At very large scale (>10k users):
• Replay per crop could become expensive
• Event queue contention may increase
Mitigation:
• Incremental replay
• Distributed message broker
• Horizontal scaling
ENHANCEMENT 10 — Governance
Centralization Risk
Add Under: 14.11 Security Limitations
Risk:
Admin holds strong authority.
Potential issues:
• Rule bias
• Misuse
• Accidental override
Mitigation:
• Admin audit log
• Immutable logs
• Role restrictions

===== PAGE 188 =====
ENHANCEMENT 11 — Ethical Over-
Dependence Risk
Strengthen: 14.12 Ethical Considerations
Add:
Risk:
Farmer may follow system blindly.
Mitigation:
• “Farmer Final Authority” messaging
• Explainable recommendations
• No automatic irreversible action
ENHANCEMENT 12 — AI Liability
Disclaimer
Add Under: 14.13 Legal & Compliance
Explicitly state:
CultivaX:
• Provides advisory insights
• Does not guarantee yield
• Does not replace agronomic expertise
• Not liable for weather unpredictability
\end{verbatim}
\section{Technical Design Document (TDD)}
\begin{verbatim}
# FILE: CultivaX TDD.pdf
# PAGES: 107



===== PAGE 1 =====
Technical Design Document (TDD)
SECTION 1 — SYSTEM IMPLEMENTATION
ARCHITECTURE
1.1 Purpose of This Document
This document specifies:
• Exact technology stack
• Concrete implementation details
• Service boundaries
• Database schema structure
• API contract format
• Event dispatcher mechanics
• ML integration plan
• Deployment architecture
• Operational constraints
This document is intended for:
• Developers
• DevOps engineers
• ML engineers
• Reviewers validating implementation feasibility
This is not conceptual.
This is executable design.
1.2 Finalized Technology Stack
Backend
• Language: Python 3.11
• Framework: FastAPI
• ASGI Server: Uvicorn (with Gunicorn for production)
• ORM: SQLAlchemy 2.0

===== PAGE 2 =====
• Migration Tool: Alembic
• Data Validation: Pydantic v2
• Auth: JWT (python-jose)
Database
• Engine: PostgreSQL 15
• Hosting: GCP Cloud SQL
• Extensions:
o uuid-ossp
o pgcrypto
o pg_stat_statements (monitoring)
• Data types used:
o UUID
o JSONB
o TIMESTAMP WITH TIME ZONE
Event System
• In-process dispatcher
• DB-backed event persistence
• Partition key: crop_instance_id
• Row-level locking for isolation
• Idempotency enforced via unique event_hash
Future-ready abstraction for Pub/Sub migration.
ML Stack
• Risk Model: Scikit-learn (LogisticRegression)
• CNN (backend): PyTorch
• Edge Model: TensorFlow Lite
• Feature Extraction: OpenCV + NumPy
• Model serialization: joblib / torch.save
Media
• Storage: GCP Cloud Storage (Private bucket)

===== PAGE 3 =====
• Access: Signed URLs (15 min expiry)
• Processing: Separate Cloud Run worker service
• Video Frame Extraction: OpenCV
Frontend
• Framework: Next.js (React 18)
• Styling: TailwindCSS
• State Management: React Context + minimal Redux (if needed)
• PWA: Service Worker enabled
• Auth storage: HttpOnly cookie (recommended)
Deployment
• Backend API: Cloud Run service
• Media Worker: Cloud Run service
• Video Worker: Cloud Run service (optional)
• Database: Cloud SQL
• Storage: Cloud Storage
• Cron Jobs: Cloud Scheduler
• Logging: Cloud Logging
• Monitoring: Cloud Monitoring
1.3 High-Level Service Topology
[ Next.js Frontend ]
↓ HTTPS
[ Cloud Run - Backend API ]
↓
[ Cloud SQL - PostgreSQL ]
[ Cloud Run - Media Worker ]
↓
[ Cloud Storage ]
[ Cloud Scheduler ]
↓
[ Backend API / Worker endpoints ]
No service shares in-memory state.

===== PAGE 4 =====
All state stored in DB.
1.4 Service Boundaries
Backend API Service Responsibilities
• Authentication
• Crop lifecycle operations
• Action logging
• Replay engine
• Risk computation
• SOE management
• Admin governance
• Event dispatch processing
• ML inference (lightweight)
• Media upload orchestration
Must NOT:
• Perform heavy video frame processing
• Store files locally
• Hold long-lived in-memory state
Media Worker Responsibilities
• Image inference
• Video frame extraction
• CNN inference
• Feature aggregation
• Emit MediaAnalyzed event
Must NOT:
• Modify CTIS state directly
• Bypass event dispatcher

===== PAGE 5 =====
Database Responsibility
Single source of truth for:
• CropInstances
• Action logs
• Event logs
• Trust scores
• ML audit
• Snapshots
• Deviation profiles
1.5 Directory Structure (Backend)
cultivax/
│
├── app/
│ ├── main.py
│ ├── config.py
│ ├── database.py
│ ├── models/
│ ├── schemas/
│ ├── api/
│ ├── services/
│ │ ├── ctis/
│ │ ├── soe/
│ │ ├── event_dispatcher/
│ │ ├── ml/
│ │ ├── media/
│ │ ├── replay/
│ ├── workers/
│ ├── utils/
│ └── security/
│
├── alembic/
├── Dockerfile
├── requirements.txt
└── .env
Clear modular boundaries.

===== PAGE 6 =====
1.6 Core Runtime Flow (Implementation-Level)
Example: Log Action
1. API endpoint receives request
2. Validate via Pydantic schema
3. Verify JWT
4. Insert ActionLog (pending)
5. Create EventLog entry (ActionLogged)
6. Commit transaction
7. Dispatcher processes event
8. ReplayTriggered event created
9. Replay executes
10. Snapshot updated
11. Response returned
All DB writes wrapped in transaction.
1.7 Concurrency & Locking Strategy
For replay:
• Acquire SELECT FOR UPDATE on crop_instance row
• Prevent concurrent replay
• Queue other events
• Release lock after commit
Prevents race conditions.
1.8 Migration Readiness Constraints
To allow migration to Pub/Sub later:
• Dispatcher interface abstracted
• Event payload serializable
• No business logic inside dispatcher
• Events immutable
This ensures scaling to 500–600 users requires no rewrite.

===== PAGE 7 =====
1.9 Non-Functional Guarantees
• Replay O(n) with snapshot optimization
• Media inference async
• API latency target < 300ms (excluding media)
• Event processing idempotent
• DB ACID consistency enforced
• No mutation without event
SECTION 2 — DATABASE SCHEMA DESIGN
2.1 Database Design Principles
All mutable state must be event-driven.
No direct state mutation without event_log entry.
Replay must be reconstructable from action_logs.
All destructive operations are soft-delete.
Idempotency enforced at DB level.
All tables must support migration without rewrite.
No in-memory-only data.
2.2 Core Identity Tables
2.2.1 users
CREATE TABLE users (
user_id UUID PRIMARY KEY DEFAULT gen_random_uuid(),
full_name TEXT NOT NULL,
phone_number TEXT UNIQUE NOT NULL,
role TEXT NOT NULL CHECK (role IN ('farmer','provider','admin')),
password_hash TEXT NOT NULL,
is_active BOOLEAN DEFAULT TRUE,
created_at TIMESTAMPTZ DEFAULT NOW(),
deleted_at TIMESTAMPTZ
);
Indexes:
CREATE INDEX idx_users_role ON users(role);
2.3 CTIS Core Tables

===== PAGE 8 =====
2.3.1 crop_instances
CREATE TABLE crop_instances (
crop_instance_id UUID PRIMARY KEY DEFAULT gen_random_uuid(),
user_id UUID REFERENCES users(user_id),
crop_type TEXT NOT NULL,
region TEXT NOT NULL,
sowing_date DATE NOT NULL,
rule_version_applied INT NOT NULL,
baseline_stage TEXT NOT NULL,
current_stage TEXT NOT NULL,
stress_score NUMERIC(5,4) DEFAULT 0.0,
risk_index NUMERIC(5,4) DEFAULT 0.0,
state TEXT CHECK (state IN ('Active','AtRisk','Harvested','RecoveryRequired')),
event_chain_hash TEXT,
created_at TIMESTAMPTZ DEFAULT NOW(),
updated_at TIMESTAMPTZ DEFAULT NOW(),
deleted_at TIMESTAMPTZ
);
Indexes:
CREATE INDEX idx_crop_user ON crop_instances(user_id);
CREATE INDEX idx_crop_state ON crop_instances(state);
2.3.2 action_logs
CREATE TABLE action_logs (
action_id UUID PRIMARY KEY DEFAULT gen_random_uuid(),
crop_instance_id UUID REFERENCES crop_instances(crop_instance_id),
action_type TEXT NOT NULL,
action_impact_type TEXT NOT NULL,
action_effective_date DATE NOT NULL,
local_monotonic_counter BIGINT NOT NULL,
metadata JSONB,
created_at TIMESTAMPTZ DEFAULT NOW(),
deleted_at TIMESTAMPTZ,

===== PAGE 9 =====
UNIQUE (crop_instance_id, action_effective_date, local_monotonic_counter)
);
Replay index:
CREATE INDEX idx_action_replay
ON action_logs(crop_instance_id, action_effective_date DESC, local_monotonic_counter DESC);
2.3.3 crop_instance_snapshots
CREATE TABLE crop_instance_snapshots (
snapshot_id UUID PRIMARY KEY DEFAULT gen_random_uuid(),
crop_instance_id UUID REFERENCES crop_instances(crop_instance_id),
snapshot_index INT NOT NULL,
stage TEXT NOT NULL,
baseline_stage TEXT NOT NULL,
stress_score NUMERIC(5,4),
risk_index NUMERIC(5,4),
deviation_profile JSONB,
rule_version_applied INT,
created_at TIMESTAMPTZ DEFAULT NOW()
);
Index:
CREATE INDEX idx_snapshot_lookup
ON crop_instance_snapshots(crop_instance_id, snapshot_index DESC);
2.3.4 deviation_profiles
CREATE TABLE deviation_profiles (
profile_id UUID PRIMARY KEY DEFAULT gen_random_uuid(),
crop_instance_id UUID REFERENCES crop_instances(crop_instance_id),
consecutive_deviation_count INT DEFAULT 0,
deviation_trend_slope NUMERIC(6,4),
recurring_pattern_flag BOOLEAN DEFAULT FALSE,
last_updated TIMESTAMPTZ DEFAULT NOW()
);
2.3.5 yield_records
CREATE TABLE yield_records (

===== PAGE 10 =====
yield_id UUID PRIMARY KEY DEFAULT gen_random_uuid(),
crop_instance_id UUID REFERENCES crop_instances(crop_instance_id),
reported_yield NUMERIC(10,2) NOT NULL,
ml_yield_value NUMERIC(10,2),
yield_verification_score NUMERIC(5,4),
submitted_at TIMESTAMPTZ DEFAULT NOW()
);
2.4 Event System Tables
2.4.1 event_log
CREATE TABLE event_log (
event_id UUID PRIMARY KEY DEFAULT gen_random_uuid(),
crop_instance_id UUID REFERENCES crop_instances(crop_instance_id),
event_type TEXT NOT NULL,
event_payload JSONB NOT NULL,
event_hash TEXT UNIQUE NOT NULL,
chain_hash TEXT,
status TEXT CHECK (status IN ('Created','Processing','Processed','Failed','DeadLetter')),
retry_count INT DEFAULT 0,
failure_reason TEXT,
priority INT DEFAULT 5,
created_at TIMESTAMPTZ DEFAULT NOW(),
processed_at TIMESTAMPTZ
);
Indexes:
CREATE INDEX idx_event_partition
ON event_log(crop_instance_id, status, priority, created_at);
2.5 SOE Tables
2.5.1 service_providers
CREATE TABLE service_providers (
provider_id UUID PRIMARY KEY DEFAULT gen_random_uuid(),
user_id UUID REFERENCES users(user_id),
equipment_type TEXT,

===== PAGE 11 =====
trust_score NUMERIC(5,4) DEFAULT 0.5,
completion_rate NUMERIC(5,4),
complaint_ratio NUMERIC(5,4),
is_verified BOOLEAN DEFAULT FALSE,
created_at TIMESTAMPTZ DEFAULT NOW(),
deleted_at TIMESTAMPTZ
);
2.5.2 service_requests
CREATE TABLE service_requests (
request_id UUID PRIMARY KEY DEFAULT gen_random_uuid(),
crop_instance_id UUID REFERENCES crop_instances(crop_instance_id),
provider_id UUID REFERENCES service_providers(provider_id),
status TEXT CHECK (status IN ('Pending','Accepted','Completed','Cancelled','Escalated')),
created_at TIMESTAMPTZ DEFAULT NOW(),
updated_at TIMESTAMPTZ DEFAULT NOW()
);
2.5.3 service_reviews
CREATE TABLE service_reviews (
review_id UUID PRIMARY KEY DEFAULT gen_random_uuid(),
request_id UUID REFERENCES service_requests(request_id),
rating INT CHECK (rating BETWEEN 1 AND 5),
complaint_category TEXT,
severity_level INT,
resolution_status TEXT,
created_at TIMESTAMPTZ DEFAULT NOW()
);
2.6 ML Tables
2.6.1 ml_models
CREATE TABLE ml_models (
model_id UUID PRIMARY KEY DEFAULT gen_random_uuid(),
model_name TEXT NOT NULL,
model_version TEXT NOT NULL,

===== PAGE 12 =====
model_type TEXT,
metrics JSONB,
created_at TIMESTAMPTZ DEFAULT NOW()
);
2.6.2 ml_training_audit
CREATE TABLE ml_training_audit (
training_id UUID PRIMARY KEY DEFAULT gen_random_uuid(),
model_version TEXT,
dataset_reference TEXT,
feature_schema_version TEXT,
evaluation_metrics JSONB,
training_start TIMESTAMPTZ,
training_end TIMESTAMPTZ,
created_at TIMESTAMPTZ DEFAULT NOW()
);
2.7 Media Tables
2.7.1 media_files
CREATE TABLE media_files (
media_id UUID PRIMARY KEY DEFAULT gen_random_uuid(),
crop_instance_id UUID REFERENCES crop_instances(crop_instance_id),
file_type TEXT CHECK (file_type IN ('image','video')),
storage_path TEXT NOT NULL,
confidence_score NUMERIC(5,4),
image_quality_score NUMERIC(5,4),
extracted_features JSONB,
source TEXT CHECK (source IN ('backend','edge')),
created_at TIMESTAMPTZ DEFAULT NOW(),
scheduled_deletion_at TIMESTAMPTZ
);
Index:
CREATE INDEX idx_media_crop
ON media_files(crop_instance_id);

===== PAGE 13 =====
2.8 Admin & Audit
2.8.1 admin_audit_log
CREATE TABLE admin_audit_log (
audit_id UUID PRIMARY KEY DEFAULT gen_random_uuid(),
admin_id UUID REFERENCES users(user_id),
action_type TEXT,
target_id UUID,
previous_state JSONB,
new_state JSONB,
reason TEXT,
created_at TIMESTAMPTZ DEFAULT NOW()
);
Immutable rule:
No DELETE allowed.
Only insert.
2.9 Retention Policy Implementation
Scheduled job:
Delete media_files where scheduled_deletion_at < NOW()
Keep extracted_features
Log deletion event
2.10 Migration Readiness
All tables use UUID PK
JSONB flexible for schema evolution
Index strategy supports scaling
No hard constraints blocking future partitioning
Can partition action_logs by crop_instance_id later
SECTION 3 — EVENT DISPATCHER
IMPLEMENTATION DESIGN

===== PAGE 14 =====
3.1 Design Goals
The Event Dispatcher must:
1. Persist every event before execution.
2. Enforce idempotency at DB level.
3. Process events in FIFO order per crop_instance.
4. Isolate partitions by crop_instance_id.
5. Guarantee atomic replay.
6. Support retry and dead-letter handling.
7. Be easily replaceable with Pub/Sub in future.
8. Avoid global in-memory state.
3.2 Dispatcher Architecture Overview
Conservative implementation:
• In-process dispatcher inside FastAPI app
• Uses DB event_log table as source of truth
• Uses SELECT FOR UPDATE SKIP LOCKED
• Processes events in small batches
• Partitioned logically by crop_instance_id
There is NO in-memory queue that stores unpersisted state.
3.3 Core Abstraction Layer (Migration-Ready)
We define a dispatcher interface.
class EventDispatcherInterface:
async def publish(self, event_type: str, payload: dict, crop_instance_id: UUID):
pass
async def process_pending_events(self):
pass
Current implementation:
DBEventDispatcher
Future implementation:
PubSubEventDispatcher
Business logic never directly depends on concrete implementation.

===== PAGE 15 =====
3.4 Event Lifecycle Model (Concrete)
Event states in DB:
• Created
• Processing
• Processed
• Failed
• DeadLetter
Lifecycle:
1. publish()
a. Insert event_log row
b. status = 'Created'
c. event_hash generated
2. process_pending_events()
a. Fetch events WHERE status='Created'
b. Order by priority, created_at
c. Lock rows
d. status='Processing'
e. Execute handler
f. On success → status='Processed'
g. On failure → increment retry_count
i. If retry_count < threshold → status='Created'
ii. Else → status='DeadLetter'
3.5 Event Hash & Idempotency Enforcement
When publishing:
event_hash = sha256(
f"{event_type}:{crop_instance_id}:{json.dumps(payload)}"
)
DB has:
UNIQUE(event_hash)
If duplicate publish attempt occurs:
• Insert fails
• Existing event returned

===== PAGE 16 =====
• No duplicate processing
This prevents:
• Double action logging
• Network retry duplication
• Offline sync duplication
3.6 Partitioned Processing Strategy
Partition key:
crop_instance_id
Processing query:
SELECT *
FROM event_log
WHERE status = 'Created'
ORDER BY priority DESC, created_at ASC
LIMIT 20
FOR UPDATE SKIP LOCKED;
Then:
Before processing each event:
SELECT * FROM crop_instances
WHERE crop_instance_id = ?
FOR UPDATE;
This ensures:
• Only one replay per crop at a time
• No cross-crop blocking
• No deadlock cascade
3.7 Replay Trigger Mechanism
Certain events automatically emit ReplayTriggered:
Events that require replay:
• ActionLogged
• YieldSubmitted

===== PAGE 17 =====
• SowingDateModified
• MediaAnalyzed (if stress updated)
Replay event:
event_type = 'ReplayTriggered'
Replay logic:
• Acquire crop lock
• Load snapshot (if exists)
• Load action_logs
• Recompute state
• Update crop_instances row
• Insert snapshot if threshold met
• Commit transaction
Replay is atomic.
If replay fails:
• Rollback transaction
• Mark event Failed
• Do not partially update crop_instance
3.8 Dead Letter Queue Handling
If retry_count exceeds threshold (e.g., 5):
• status='DeadLetter'
• Insert admin alert
• Mark crop_instance state='RecoveryRequired'
DeadLetter events do NOT block other events.
Admin endpoint:
GET /api/v1/admin/dead-letters
Allows inspection and manual retry.
3.9 Backpressure Strategy
If event backlog grows:

===== PAGE 18 =====
• process_pending_events() runs in background loop
• Max concurrent replays limited (e.g., 5)
• Use asyncio.Semaphore to control concurrency
If queue length > threshold:
• Raise monitoring alert
• Reduce intake rate (soft throttle)
Prevents CPU exhaustion.
3.10 Debounce Replay Optimization
If multiple replay-triggering events arrive quickly:
• Merge them logically
• Only one ReplayTriggered inserted per crop per short interval (e.g., 300ms window)
Prevents replay storm during bulk sync.
3.11 Crash Recovery Behavior
On service restart:
• Any event with status='Processing'
→ Reset to 'Created'
→ retry_count += 1
Ensures:
• No event permanently stuck
• No partial state applied
Because state changes occur only after successful commit.
3.12 Event Dependency Map (Hard-Coded Rules)
Allowed chains:
ActionLogged → ReplayTriggered
YieldSubmitted → ReplayTriggered

===== PAGE 19 =====
MediaAnalyzed → ReplayTriggered
RuleModified → (No automatic replay)
Forbidden:
RegionalClusterUpdated → ReplayTriggered (retroactively)
This prevents infinite loops.
3.13 Concurrency Model
Event loop:
• Background task started in FastAPI startup event
• Async loop polling DB every X ms
• Uses SKIP LOCKED to avoid contention
• Each Cloud Run instance runs its own loop
Row-level locking ensures no double execution across instances.
3.14 Performance Envelope
Expected load:
• 200 users
• ~20 concurrent sessions
• <100 events/minute
This design handles that easily.
For 600 users:
• Increase Cloud Run max instances
• Increase DB pool size
• Possibly reduce batch size
• Still no rewrite required
3.15 Migration Plan to Pub/Sub
Because:

===== PAGE 20 =====
• EventDispatcherInterface exists
• Event payload serializable
• No business logic inside dispatcher
Migration steps:
1. Replace DBEventDispatcher.publish()
2. Publish to Pub/Sub topic
3. Subscriber service consumes events
4. Same handler functions reused
Zero business logic rewrite.
3.16 Security Considerations
• No direct table mutation outside event processing
• All event handlers wrapped in transaction
• Replay cannot be triggered without event entry
• Event hash prevents tampering
• Chain hash verification can be added
SECTION 4 — CTIS IMPLEMENTATION
DESIGN
4.1 CTIS Core Responsibilities
CTIS must:
1. Maintain chronological integrity.
2. Recompute crop state deterministically.
3. Enforce invariants.
4. Apply drift constraints.
5. Integrate multi-signal stress safely.
6. Separate Farmer Truth from ML Truth.
7. Never mutate without event trigger.
8. Support snapshot optimization.

===== PAGE 21 =====
4.2 Core Data Inputs for Replay
Replay uses:
• crop_instances
• action_logs (ordered)
• rule_version_applied
• deviation_profiles
• yield_records
• latest stress signals
• weather risk (current snapshot)
Replay does NOT:
• Depend on external API live calls
• Use ephemeral in-memory values
All needed data must be stored or versioned.
4.3 Chronological Invariants (Code-Level
Enforcement)
Before replay begins:
def validate_invariants(crop, actions):
for action in actions:
assert action.action_effective_date >= crop.sowing_date
assert action.local_monotonic_counter >= 0
Additionally:
• No action date beyond current_date + tolerance_limit
• No stage skip allowed
• No negative stress
• No direct state mutation without event
If invariant fails:
• Raise ReplayIntegrityException
• Mark crop_instance as RecoveryRequired

===== PAGE 22 =====
4.4 Replay Algorithm (Precise Pseudocode)
def replay_crop_instance(crop_instance_id):
acquire_row_lock(crop_instance_id)
crop = load_crop_instance()
actions = load_ordered_actions()
snapshot = load_latest_snapshot()
if snapshot exists:
state = snapshot_state
start_index = snapshot.snapshot_index
else:
state = initialize_from_baseline()
start_index = 0
for action in actions[start_index:]:
validate_action_invariant(action)
state = apply_action(state, action)
state = apply_stress_integration(state)
state = compute_risk_index(state)
state = enforce_drift_constraints(state)
update_crop_instance(state)
maybe_create_snapshot()
release_lock()
Atomic transaction enforced.
4.5 Stage Transition Logic
Each crop has:
CropRuleTemplate:
{
"stages": [
{ "name": "Germination", "duration_days": 7 },
{ "name": "Tillering", "duration_days": 21 },
...

===== PAGE 23 =====
]
}
Baseline stage computed as:
days_since_sowing = today - sowing_date
baseline_stage = stage_lookup(days_since_sowing)
Actual stage adjusted by:
• deviation_offset
• stress_modifier
But must not exceed:
• max_stage_drift (e.g., ±14 days)
No direct stage jump allowed.
4.6 Drift Constraint Enforcement
After state computation:
if abs(state.stage_offset) > crop.max_allowed_drift:
state.stage_offset = clamp(state.stage_offset)
Prevents:
• Infinite delay
• Infinite acceleration
Drift bounded per crop type.
4.7 Stress Integration Formula
Multi-signal arbitration:
stress_signal = (
w_backend * backend_ml_stress +
w_weather * weather_risk +
w_deviation * deviation_penalty +
w_edge * edge_signal
)
new_stress = alpha * (stress_signal * confidence) + (1 - alpha) * previous_stress

===== PAGE 24 =====
Where:
w_backend > w_weather > w_deviation > w_edge
Constraints:
• stress in [0,1]
• max daily increase limit
• max single-event jump limit
4.8 Risk Index Computation
System Risk Index:
risk_index = (weather_risk * 0.7) + (farmer_input_risk * 0.3)
Adjusted by:
• stress_score
• stage sensitivity factor
Must remain probabilistic (0–1).
4.9 Yield Truth Enforcement
When YieldSubmitted event processed:
reported_yield = input_value
if reported_yield > biological_limit:
ml_yield_value = biological_limit
else:
ml_yield_value = reported_yield
Store both.
UI always shows reported_yield.
ML training uses ml_yield_value.
Never modify reported_yield.

===== PAGE 25 =====
4.10 Snapshot Strategy
Snapshot created when:
• action_count % 50 == 0
OR
• major structural action occurs
Snapshot stores:
• stage
• stress_score
• risk_index
• deviation_profile
Reduces replay cost.
4.11 Structural Action Handling
If action_impact_type == 'Structural':
• Invalidate snapshot
• Full replay required
• Recompute baseline
Examples:
• SowingDateModified
4.12 Behavioral Adaptation Enforcement
Deviation logic:
• Track consecutive deviation count
• If recurring pattern detected:
o Apply small bounded adaptation offset
o Must not exceed max_adaptation_days
Reset at season end.

===== PAGE 26 =====
4.13 What-If Simulation Implementation
For What-If:
• Clone crop_instance state in memory
• Apply hypothetical action
• Run replay logic
• Do NOT persist
• Return projected state
Never emit event.
4.14 State Transition Safeguards
Valid transitions:
Active → AtRisk
Active → Harvested
AtRisk → Active
Active → RecoveryRequired
Invalid transitions blocked by:
if not is_valid_transition(old_state, new_state):
raise InvalidStateTransition()
4.15 Concurrency Safety
Replay executed inside:
• DB transaction
• Row-level lock
No two replays on same crop simultaneously.
4.16 Performance Envelope
With snapshots:
Replay complexity:

===== PAGE 27 =====
O(k) where k = actions since last snapshot.
For 200 users:
Negligible overhead.
For 600 users:
Still safe with tuning.
4.17 ML Kill Switch Integration
If feature_flag.ml_enabled == False:
• Skip ML signal in stress integration
• Use rule-only computation
No code branching outside controlled function.
4.18 Logging & Observability
Each replay logs:
• replay_duration
• actions_processed
• snapshot_used
• drift_applied
• stress_delta
Helps debugging and performance tuning.
SECTION 5 — SERVICE
ORCHESTRATION ENGINE (SOE)
IMPLEMENTATION DESIGN

===== PAGE 28 =====
5.1 SOE Responsibilities
SOE handles:
1. Service provider registration
2. Equipment availability
3. Service request matchmaking
4. Complaint handling
5. Trust score computation
6. Exposure fairness
7. Review audit
8. Escalation workflow
SOE must NOT:
• Mutate crop stage
• Trigger replay
• Influence CTIS baseline
• Override farmer authority
CTIS and SOE are logically separate systems.
5.2 Provider Registration Flow
When provider registers:
1. Create user (role='provider')
2. Insert into service_providers
3. Default trust_score = 0.5
4. is_verified = False
Admin may later verify provider.
No trust score advantage given for early signup.
5.3 Equipment & Availability Model
Add availability table:
CREATE TABLE provider_availability (
availability_id UUID PRIMARY KEY DEFAULT gen_random_uuid(),
provider_id UUID REFERENCES service_providers(provider_id),
available_from DATE,

===== PAGE 29 =====
available_to DATE,
equipment_status TEXT CHECK (equipment_status IN ('Available','Maintenance','Unavailable'))
);
Match requests only if:
current_date ∈ available window
equipment_status = 'Available'
5.4 Service Request Flow (Implementation)
Service request lifecycle:
1. Farmer creates request
2. Matching providers retrieved
3. Sorted by exposure-adjusted trust
4. Provider accepts or declines
5. Completion status updated
6. Review submitted
7. Trust recalculated
All mutations wrapped in transaction.
5.5 Trust Score Model (Mathematically Defined)
Trust score components:
• Completion Rate (CR)
• Complaint Ratio (CPR)
• Average Rating (AR)
• Verified Bonus (VB)
• Escalation Penalty (EP)
Formula:
trust_score = (
w1 * CR +
w2 * (1 - CPR) +
w3 * normalized_rating +
w4 * VB -
w5 * EP
)
Weights:

===== PAGE 30 =====
w1 > w2 > w3
w4 small bonus
w5 moderate penalty
All values normalized 0–1.
Clamp result to [0,1].
5.6 Trust Decay Mechanism
To prevent permanent dominance:
Apply time decay:
decay_factor = exp(-lambda * days_since_last_completion)
trust_score *= decay_factor
Ensures inactive providers gradually lose ranking advantage.
5.7 Exposure Fairness Algorithm
When listing providers:
Sort by:
ranking_score = trust_score - exposure_penalty
Exposure penalty computed as:
exposure_penalty = alpha * recent_impressions_count
This prevents:
• Top providers monopolizing visibility
• New providers being invisible
Ensures marketplace fairness.
5.8 Complaint Handling Logic
Complaint submission:

===== PAGE 31 =====
• Insert service_review
• Update complaint_ratio
• If severity > threshold:
o Escalate to admin
o status = 'Escalated'
Admin resolution required for:
• Repeated high-severity complaints
• Fraud suspicion
Escalation does NOT block other provider operations automatically.
5.9 Fraud Detection Rules
Flag if:
• Same IP multiple reviews
• Review from linked account
• Sudden rating spike
• High complaint but high rating mismatch
Flagged providers:
• Mark under review
• Reduce exposure temporarily
No automatic ban without admin review.
5.10 CTIS Isolation Guarantee
SOE must never:
• Call replay
• Modify crop_instance
• Influence stress_score
• Influence risk_index
Only shared link:
service_requests references crop_instance_id.
That is informational only.

===== PAGE 32 =====
5.11 Concurrency Model
For trust recalculation:
• Lock provider row FOR UPDATE
• Recalculate
• Update atomically
Prevents concurrent inconsistent trust updates.
5.12 Service Review Constraints
Rules:
• Only one review per request
• Cannot review before completion
• Cannot self-review
DB-level constraints:
ALTER TABLE service_reviews
ADD CONSTRAINT unique_review_per_request UNIQUE(request_id);
5.13 Escalation Workflow
Escalation path:
Complaint → Severity Check → Escalate → Admin Review → Resolution
Admin options:
• Dismiss
• Warn provider
• Reduce trust score
• Suspend provider (soft delete)
• Permanently deactivate
All admin actions logged in admin_audit_log.

===== PAGE 33 =====
5.14 Performance Envelope
Expected load:
• <50 service requests per day
• Trust recalculation minimal CPU
• Exposure sorting lightweight
Safe up to 600 users.
Beyond 1000:
• Add caching layer
• Precompute ranking score
No redesign needed.
5.15 Migration Readiness
If SOE needs to scale independently:
• Extract into separate microservice
• Keep DB shared initially
• Later move to separate DB
Because:
• Logic is self-contained
• No CTIS dependency
5.16 Observability
Log:
• Trust recalculation
• Escalation events
• Exposure adjustments
• Fraud flags
Helps audit marketplace fairness.

===== PAGE 34 =====
SECTION 6 — ML IMPLEMENTATION
DESIGN
6.1 ML Design Principles
1. 2. 3. 4. 5. 6. 7. ML must never override deterministic CTIS rules.
ML outputs must be bounded and probabilistic.
All ML outputs must carry confidence.
All training runs must be auditable.
Small dataset bias must be controlled.
No PII used in training.
All model versions stored and traceable.
6.2 ML System Architecture
ML Components:
1. Risk Prediction Model (Logistic Regression)
2. Stress Detection CNN (PyTorch)
3. Regional Cluster Learning (Statistical)
4. Yield Verification Model
5. Confidence Propagation Layer
ML runs inside backend container
Heavy training runs as background job.
6.3 Feature Engineering Pipeline
6.3.1 Risk Model Features
Input features:
• stress_score
• deviation_count
• deviation_trend_slope

===== PAGE 35 =====
• weather_risk
• days_since_sowing
• action_density_last_14_days
• historical_average_yield_region
• crop_stage_numeric
All features normalized to [0,1] or standardized.
Feature schema version tracked:
feature_schema_version = "v1.0"
Stored in ml_training_audit.
6.4 Risk Model (Logistic Regression)
Purpose:
Predict probability of AtRisk state.
Implementation:
model = LogisticRegression()
model.fit(X_train, y_train)
Output:
risk_probability = model.predict_proba(X)[1]
Constraints:
• Output ∈ [0,1]
• Must not exceed CTIS stage drift authority
• Only influences risk_index, not stage directly
If ML disabled:
Use rule-based fallback.
6.5 Stress Detection CNN (Backend)
Model type:
Lightweight CNN (e.g., ResNet18 variant)

===== PAGE 36 =====
Input:
• 224x224 image
• RGB normalized
Output:
• stress_probability ∈ [0,1]
• confidence_score ∈ [0,1]
Inference pipeline:
image → preprocessing → model.forward() → sigmoid
Post-processing:
stress_weighted = stress_probability * confidence_score
Bounded daily stress increase enforced.
6.6 Multi-Angle Media Aggregation
When multiple images uploaded:
Aggregate:
aggregate_stress = weighted_mean(
stress_i,
weight = confidence_i
)
If variance high:
Reduce final confidence.
Prevents single bad frame corruption.
6.7 Video Frame Pipeline
Video worker:
1. 2. Extract 1 frame per N seconds
Run CNN inference per frame
3. Remove low-quality frames
4. Aggregate via trimmed mean

===== PAGE 37 =====
5. Compute final stress + confidence
Heavy processing isolated in worker.
6.8 Regional Cluster Learning
Purpose:
Detect regional adaptation patterns.
Method:
• Group by region + crop_type
• Collect verified yield + deviation patterns
• Use K-Means clustering OR statistical mean shift
Constraint:
Minimum sample size threshold (e.g., 30 instances)
If cluster_size < threshold:
Do NOT apply adaptation.
Cluster update rule:
Prospective only
No retroactive mutation.
6.9 Yield Verification Logic
When farmer submits yield:
Compute:
z_score = (reported_yield - regional_mean) / regional_std
If |z_score| > threshold:
• reduce yield_verification_score
• clamp ML yield value
Never modify reported_yield.
Store:

===== PAGE 38 =====
reported_yield
ml_yield_value
yield_verification_score
6.10 Confidence Propagation Layer
Every ML output must include:
confidence_score
Stress integration uses:
final_stress_delta = raw_stress * confidence_score
Low confidence reduces impact.
Confidence sources:
• Image quality score
• Frame consistency
• Model certainty
• Cluster stability
• Sample size adequacy
6.11 Small Dataset Safeguards
Because initial dataset small:
Rules:
• Freeze ML if dataset < minimum threshold
• Use rule-based risk only
• Apply stronger smoothing factor
• Reduce stress weight
This prevents overfitting early.
6.12 Model Versioning Strategy
Each model:

===== PAGE 39 =====
• Stored with model_version
• Version saved in crop_instance
• Inference logs include model_version
• Training runs logged in ml_training_audit
Allows rollback if needed.
6.13 Training Lifecycle
Training triggered:
• Manually by admin
• Or scheduled batch (Cloud Scheduler)
Training pipeline:
1. Extract dataset
2. Filter invalid entries
3. Remove outliers
4. Split train/test
5. Evaluate metrics
6. Save model
7. Log metrics
8. Activate model
No automatic activation without admin approval.
6.14 ML Kill Switch
Feature flag:
ml_enabled = True/False
If False:
• Skip CNN stress
• Skip logistic regression
• Use rule-only risk
System still functional.

===== PAGE 40 =====
6.15 Drift Detection
Track:
• Prediction confidence trend
• Regional cluster variance
• Increase in outlier frequency
If anomaly detected:
• Alert admin
• Pause cluster adaptation
6.16 Performance Envelope
Expected load:
• 50–100 images/day
• 10–20 videos/day
• Risk inference per action
Python handles this easily.
For 600 users:
• Scale worker instances
• Possibly move training to separate service
No redesign required.
6.17 Data Privacy Enforcement
Before training:
• Remove user_id
• Remove phone_number
• Remove exact coordinates
• Only keep region-level data
ML training dataset must be anonymized.

===== PAGE 41 =====
6.18 Research Extension Hooks
Later:
• Replace logistic regression with XGBoost
• Add time-series LSTM
• Add advanced crop health segmentation
• Add anomaly detection autoencoder
Because architecture modular, easy to extend.
SECTION 7 — MEDIA PROCESSING
IMPLEMENTATION
7.1 Media Architecture Overview
Media pipeline consists of:
1. Frontend upload request
2. Signed URL generation
3. Upload to Cloud Storage
4. MediaUploaded event creation
5. Worker service processes media
6. Extracted features stored in DB
7. MediaAnalyzed event emitted
8. ReplayTriggered (if stress changed)
No synchronous inference in API path.
7.2 Upload Flow (Implementation-Level)
Step 1 — Client requests upload token
API:
POST /api/v1/media/request-upload
Backend:
• Validate JWT

===== PAGE 42 =====
• Generate signed URL (15 min expiry)
• Insert media_files row (status=Pending)
Step 2 — Client uploads directly to Cloud Storage
No file flows through backend memory.
Step 3 — Client confirms upload
POST /api/v1/media/confirm-upload
Backend:
• Verify object exists in bucket
• Publish MediaUploaded event
7.3 Cloud Storage Configuration
Bucket settings:
• Private bucket
• Uniform bucket-level access enabled
• No public ACL
• Lifecycle rule: delete after 90 days
• Versioning optional (off for prototype)
• Object path structure:
/crop_instance_id/media_id/file.ext
Signed URL validity:
15 minutes
7.4 Media Worker Service
Deployed as separate Cloud Run container.
Responsibilities:
• Listen for MediaUploaded events
• Download file from bucket
• Perform inference
• Update media_files.extracted_features
• Emit MediaAnalyzed event

===== PAGE 43 =====
Worker must NOT:
• Modify crop_instances directly
• Trigger replay directly
• Bypass dispatcher
7.5 Image Processing Pipeline
Flow:
1. Load image via OpenCV
2. Compute image_quality_score
a. brightness check
b. blur detection (Laplacian variance)
c. contrast check
3. Reject if quality below threshold
4. Resize to 224x224
5. Normalize pixel values
6. Run PyTorch model inference
7. Compute stress_probability
8. Compute confidence_score
Confidence calculation includes:
• Image quality
• Model output certainty
• Variance across augmentation passes (optional)
Final result stored as:
{
"stress_probability": 0.42,
"confidence": 0.78,
"leaf_curl_detected": true,
"yellowing_score": 0.35
}
7.6 Video Processing Pipeline
Video worker steps:
1. Download video
2. 3. Extract frame every N seconds (configurable)
Limit max frames (e.g., 30 frames)

===== PAGE 44 =====
4. Run image pipeline on each frame
5. Remove low-quality frames
6. Compute trimmed mean stress
7. Compute variance across frames
8. Reduce confidence if variance high
Final stored:
{
"aggregate_stress": 0.51,
"confidence": 0.72,
"frame_count": 18,
"variance": 0.11
}
7.7 Multi-Media Aggregation Logic
If multiple images/videos exist for same day:
Aggregate using:
weighted_mean(stress_i, weight=confidence_i)
If disagreement high:
Reduce final confidence.
No direct stage mutation.
Only stress_score updated.
7.8 MediaAnalyzed Event Handling
Worker publishes:
event_type = "MediaAnalyzed"
Payload includes:
• crop_instance_id
• stress_delta
• confidence_score
• media_id
Dispatcher:

===== PAGE 45 =====
• Inserts event
• ReplayTriggered event created
• CTIS integrates stress in bounded way
7.9 Edge AI Integration (Optional Advanced
Mode)
If device uses Edge AI:
• Device computes preliminary stress
• Sends extracted summary
• Backend stores as source='edge'
• Backend still runs verification if full media later uploaded
If low data mode:
Only extracted summary stored.
Confidence lower for edge-only data.
7.10 Retention & Deletion Pipeline
Cloud Scheduler runs daily:
1. Query media_files where scheduled_deletion_at < NOW()
2. Delete object from bucket
3. Keep extracted_features
4. Log deletion in admin_audit_log
Extracted features retained permanently.
Replay unaffected.
7.11 Failure Handling
If worker crashes during processing:
• Event remains Processing
• On restart → reset to Created
• retry_count incremented

===== PAGE 46 =====
• After 5 retries → DeadLetter
DeadLetter visible in admin panel.
7.12 Rate Limiting & Abuse Control
Limits:
• Max images per crop per day (configurable, e.g., 10)
• Max video length (e.g., 60 sec)
• Max video size (e.g., 50MB)
Reject excessive uploads.
Prevents storage abuse.
7.13 Performance Envelope
Expected load:
100–200 users:
• <100 images/day
• <20 videos/day
600 users:
• ~300 images/day
• ~60 videos/day
Scaling strategy:
• Increase worker instances
• Limit frame extraction
• Increase CPU allocation
No redesign required.
7.14 Security Enforcement
Worker:

===== PAGE 47 =====
• Uses service account with bucket access only
• No DB root credentials
• No local file persistence
• Files processed in memory or temp storage
Media never public.
Signed URLs required for access.
7.15 Observability
Worker logs:
• processing_time
• model_version
• confidence_score
• frame_count
• failure_reason
Alerts triggered if:
• processing_time > threshold
• failure_rate spike
• backlog grows
SECTION 8 — API CONTRACT
SPECIFICATION
8.1 API Design Principles
1. All endpoints versioned (/api/v1/)
2. 4. 5. 6. 7. All state-changing endpoints require Idempotency-Key header
3. Uniform response envelope
JWT required for protected endpoints
No direct state mutation without event emission
Bulk sync supported for offline mode
Long-running tasks tracked via operation_id

===== PAGE 48 =====
8.2 Base URL Structure
/api/v1/
Future breaking changes → /api/v2/
8.3 Standard Response Envelope
Success:
{
"success": true,
"data": {},
"meta": {
"request_id": "uuid",
"timestamp": "ISO8601",
"version": "v1"
}
}
Error:
{
"success": false,
"error": {
"code": "ERROR_CODE",
"message": "Human readable message",
"details": {}
}
}
8.4 Authentication Endpoints
POST /api/v1/auth/login
Request:
{
"phone_number": "string",
"password": "string"
}

===== PAGE 49 =====
Response:
{
"access_token": "jwt",
"role": "farmer"
}
JWT expiry: 1 hour
Refresh token optional (future)
8.5 Crop Lifecycle APIs (CTIS)
POST /api/v1/crops
Create crop instance.
Headers:
Idempotency-Key: UUID
Request:
{
"crop_type": "Wheat",
"region": "Haryana",
"sowing_date": "YYYY-MM-DD"
}
Triggers:
• CropCreated event
GET /api/v1/crops
Returns active crops.
Optional:
?include_archived=true

===== PAGE 50 =====
GET /api/v1/crops/{crop_id}
Returns:
• current_stage
• stress_score
• risk_index
• state
• deviation_profile
POST /api/v1/crops/{crop_id}/actions
Headers:
Idempotency-Key: UUID
Request:
{
"action_type": "Irrigation",
"action_effective_date": "YYYY-MM-DD",
"local_monotonic_counter": 12,
"metadata": {}
}
Triggers:
• ActionLogged
• ReplayTriggered
POST /api/v1/crops/{crop_id}/sync-batch
Used for offline bulk sync.
Headers:
Idempotency-Key: UUID
Request:
{
"actions": [
{
"action_type": "Irrigation",
"action_effective_date": "YYYY-MM-DD",
"local_monotonic_counter": 1
}

===== PAGE 51 =====
],
"device_session_id": "UUID"
}
Server:
• Inserts all valid actions
• Triggers single ReplayTriggered
Response:
{
"accepted_count": 5,
"rejected": []
}
POST /api/v1/crops/{crop_id}/yield
Headers:
Idempotency-Key: UUID
Request:
{
"reported_yield": 4800
}
Triggers:
• YieldSubmitted
• ReplayTriggered
8.6 Media APIs
POST /api/v1/media/request-upload
Request:
{
"crop_instance_id": "UUID",
"file_type": "image"
}

===== PAGE 52 =====
Response:
{
"upload_url": "signed_url",
"media_id": "UUID"
}
POST /api/v1/media/confirm-upload
Request:
{
"media_id": "UUID"
}
Triggers:
• MediaUploaded
GET /api/v1/media/{media_id}
Returns signed download URL.
8.7 SOE APIs
GET /api/v1/providers
Optional filters:
• equipment_type
• region
Sorted by exposure-adjusted trust.
POST /api/v1/service-requests
Headers:
Idempotency-Key

===== PAGE 53 =====
Request:
{
"crop_instance_id": "UUID",
"provider_id": "UUID"
}
POST /api/v1/service-requests/{id}/complete
POST /api/v1/service-requests/{id}/review
Request:
{
"rating": 4,
"complaint_category": "Delay",
"severity_level": 2
}
Triggers:
• TrustScoreUpdated
8.8 Admin APIs
Role: admin only
GET /api/v1/admin/dead-letters
POST /api/v1/admin/providers/{id}/verify
POST /api/v1/admin/crops/{id}/force-replay

===== PAGE 54 =====
DELETE /api/v1/admin/users/{id}
Soft delete only.
All actions logged in admin_audit_log.
8.9 What-If Simulation
POST /api/v1/crops/{id}/simulate
Request:
{
"hypothetical_action": {
"action_type": "DelayIrrigation",
"days": 3
}
}
Response:
Projected stage, stress, risk.
No event emitted.
8.10 Long-Running Operation Tracking
GET /api/v1/operations/{operation_id}
Returns:
{
"status": "processing",
"progress": 60
}
Used for:
• Video processing
• Replay heavy batches

===== PAGE 55 =====
8.11 Idempotency Rules
All state-changing endpoints require:
Idempotency-Key
Server behavior:
• Store key in idempotency table
• Reject duplicate within window
• Return previous response
8.12 Rate Limiting
Per user:
• 60 requests/min
• Media uploads limited separately
• Admin endpoints stricter logging
8.13 Error Codes (Standardized)
Examples:
• INVALID_INPUT
• UNAUTHORIZED
• FORBIDDEN
• REPLAY_IN_PROGRESS
• DEAD_LETTER_EVENT
• ML_DISABLED
• MEDIA_PROCESSING_FAILED
• RATE_LIMIT_EXCEEDED
8.14 Role Enforcement Matrix
Endpoint Farmer Provider Admin
Create Crop ✔ ✖ ✔
Log Action ✔ ✖ ✔

===== PAGE 56 =====
Create Service ✔ ✖ ✔
Review Service ✔ ✖ ✔
Dead Letters ✖ ✖ ✔
Enforced in middleware.
8.15 Health Endpoint
GET /api/v1/health
Returns:
• DB connectivity
• Event queue status
• Worker status
• ML model loaded
Used by Cloud Run health checks.
SECTION 9 — AUTHENTICATION &
AUTHORIZATION IMPLEMENTATION
DESIGN
This section must align with:
• Phase 3 (Construction) expectations
Lab_Manual
• Your finalized tech stack (JWT-based auth)
CultivaX Tech stack
• Replay integrity (CTIS must be protected)
• Admin audit traceability
• Migration readiness (OAuth/Firebase later)
We will design this cleanly and conservatively.

===== PAGE 57 =====
9.1 Security Design Principles
1. 2. 3. 4. 5. 7. 8. All protected endpoints require authentication.
All state mutations require valid JWT + Idempotency-Key.
Role-based access control (RBAC) enforced at middleware.
Admin actions must be audit-logged.
No trust in frontend role flags.
6. JWT must be short-lived.
Passwords never stored in plaintext.
Soft-delete users instead of hard deletion.
9.2 Authentication Architecture
Authentication Model: Stateless JWT
Flow:
User → Login → Backend validates credentials →
JWT issued →
Frontend sends JWT in Authorization header →
Backend validates token →
Role enforced →
Request processed
No server-side session storage required.
Scalable across Cloud Run instances.
9.3 JWT Structure
Header:
{
"alg": "HS256",
"typ": "JWT"
}
Payload:
{
"sub": "user_id",
"role": "farmer",
"exp": timestamp,

===== PAGE 58 =====
"iat": timestamp,
"iss": "cultivax"
}
Important Claims:
• sub → UUID of user
• role → farmer | provider | admin
• exp → 1 hour expiry
• iss → cultivax
Signed using strong secret key (stored in Secret Manager).
9.4 Password Security
Password storage:
• Hashing algorithm: bcrypt
• Salt automatically handled by bcrypt
• Minimum length enforcement
• Optional: password strength validator
Database stores only:
password_hash
Never store raw password.
9.5 Login Flow (Implementation-Level)
1. Receive phone + password
2. Fetch user row
3. Verify bcrypt hash
4. If valid:
a. Generate JWT
b. Return token
5. If invalid:
a. Return UNAUTHORIZED
Rate limit login attempts to prevent brute force.

===== PAGE 59 =====
9.6 Authorization Model (RBAC)
Role matrix:
Endpoint Type Farmer Provider Admin
Crop CRUD ✔ ✖ ✔
Service listing ✔ ✔ ✔
Trust update ✖ ✖ ✔
Dead letter ✖ ✖ ✔
Force replay ✖ ✖ ✔
Enforced via FastAPI dependency:
def require_role(required_roles):
...
No business logic inside controller checks.
9.7 Middleware Enforcement Flow
Each protected request:
1. Extract Authorization header
2. Verify JWT signature
3. Validate expiration
4. Extract user_id + role
5. Attach user to request context
6. Apply role guard
7. Continue
If token expired:
Return 401.
9.8 Admin Protection Model
Admin endpoints require:
• role == admin
• additional logging
Before executing admin action:
Insert admin_audit_log entry:

===== PAGE 60 =====
{
admin_id,
action_type,
target_id,
previous_state,
new_state,
reason
}
No admin action without audit trail.
9.9 Soft Delete Enforcement
When admin deletes user:
UPDATE users SET is_active = false, deleted_at = NOW();
User cannot login afterward.
Historical data preserved.
9.10 Rate Limiting
Per user:
• 60 requests/minute
• 10 media uploads/day
• 5 login attempts/minute
Implemented via:
• Middleware counter (prototype)
• Later: Redis-based limiter
Protects against API abuse.
9.11 Token Refresh Strategy (Phase 1
Conservative)
For prototype:

===== PAGE 61 =====
• Access token: 1 hour
• No refresh token initially
Future:
• Add refresh endpoint
• Store refresh token in DB
• Enable multi-device session control
No redesign needed.
9.12 Device Session Tracking (Offline Mode
Support)
For bulk sync:
Include device_session_id in sync requests.
Store:
device_sessions:
- session_id
- user_id
- last_seen
- device_info
Helps detect:
• Duplicate sync attempts
• Abuse patterns
9.13 Secure Storage on Frontend
Preferred:
• Store JWT in HttpOnly cookie (safer)
Alternative (if needed):
• Store in memory (avoid localStorage)
Never store token in plain localStorage without caution.

===== PAGE 62 =====
9.14 CSRF Protection
If using cookies:
• Enable CSRF token mechanism
If using Authorization header:
• CSRF risk minimal
9.15 HTTPS Enforcement
All endpoints must be served over HTTPS.
Cloud Run enforces TLS by default.
No HTTP allowed.
9.16 Account Lockout Policy
If > 5 failed logins in 10 minutes:
• Temporarily lock account for 15 minutes
• Log event
Prevents brute force.
9.17 Audit Trail Enforcement
All state mutations log:
• user_id
• endpoint
• timestamp
• IP address
Stored in structured logs (Cloud Logging).

===== PAGE 63 =====
9.18 Scaling to 600 Users
JWT-based stateless auth:
• No session storage bottleneck
• No sticky sessions required
• Cloud Run instances can scale freely
No redesign required for 600 users.
9.19 Future Authentication Migration
Because auth logic abstracted:
Later can integrate:
• Firebase Auth
• OAuth
• Aadhaar-like government auth (if required)
Only token verification layer changes.
Business logic unaffected.
9.20 Security Alignment with Lab Manual
Phase 1 requires security requirements in SRS
Lab_Manual
Phase 2 requires architecture + security protocols
Lab_Manual
This section fulfills both:
✔ Encryption method (bcrypt)
✔ Role-based control
✔ Token security
✔ Audit logging
✔ Rate limiting

===== PAGE 64 =====
Fully compliant.
SECTION 10 — FRONTEND
INTEGRATION DESIGN
This section must align with:
• Your Canonical Architecture (Time-centric system)
CultivaX
• Lab Phase 2 requirements (System Design + Modeling)
Lab_Manual
• Finalized Tech Stack (Next.js PWA)
CultivaX Tech stack
• Event-driven backend + replay model
• Accessibility-first philosophy
• Offline + sync behavior
This is not UI decoration. This is frontend engineering design.
10.1 Frontend Architecture Overview
Frontend Stack (Locked):
• Next.js (React 18)
• TypeScript (recommended for clarity)
• TailwindCSS
• PWA enabled
• Service Worker for offline caching
• JWT auth via secure cookie
Architecture Pattern:
UI Components
↓
State Layer (Context)
↓
API Client Layer

===== PAGE 65 =====
↓
Backend (FastAPI)
No business logic stored permanently in frontend.
Frontend is presentation + local queue manager only.
10.2 Core Frontend Modules
The frontend is divided into modules aligned with canonical architecture
CultivaX
:
1. Crop Timeline View
2. Action Logging Panel
3. Media Upload Panel
4. What-If Simulator UI
5. Service Marketplace View
6. Dashboard & Insights
7. Admin Panel
8. Accessibility Control Panel
10.3 Crop Timeline UI (Core Experience)
This is the most important screen.
Layout Structure:
Top Section:
• 2D Crop Visualization Canvas
• Weather Indicator
• Risk Indicator
Middle Section:
• Current Day
• Growth Stage
• Stress Indicator Bar
• Drift Indicator
Bottom Section:

===== PAGE 66 =====
• “Recommended Action”
• Primary Button: Mark as Done
• Secondary Button: Override
• “Simulate” Option
Key Rule:
Timeline shifts must reflect backend response, not local assumption.
No optimistic stage mutation without confirmation.
10.4 2D Dynamic Crop Visualization
Implementation Strategy:
• SVG-based dynamic plant rendering
• Growth height mapped to stage progression
• Color hue adjusted via stress_score
• Animated wind effect optional
Mapping Example:
height = stage_progress * max_height
color = healthy_green → yellow → brown
Rendering must be deterministic based on backend data.
No fake visual states.
10.5 Offline Mode Architecture
Frontend must support:
• Action logging while offline
• Media queueing
• Sync on reconnect
Local Storage Strategy:
Use IndexedDB to store:
• Pending actions
• Pending media confirmations
• Device session id
• local_monotonic_counter

===== PAGE 67 =====
Sync Flow:
When connection restored:
1. Batch send /sync-batch
2. Send in chronological order
3. Receive accepted_count
4. Clear confirmed local entries
Important:
Frontend never assumes replay success until server response received.
10.6 Local Monotonic Counter Handling
Each device session:
• Maintains incremental counter
• Attached to each action
This preserves deterministic replay ordering.
If conflict detected:
Server rejects duplicate.
Frontend must handle gracefully.
10.7 Media Upload UI Flow
Steps:
1. User selects image/video
2. Call request-upload endpoint
3. Upload directly to signed URL
4. Call confirm-upload
5. Show “Processing” indicator
6. Poll operation status (if needed)
UI states:
• Uploading
• Processing
• Completed
• Failed (retry available)

===== PAGE 68 =====
Never block UI during inference.
10.8 What-If Simulation UI
Simulation panel:
• Select hypothetical action
• Adjust delay days (slider)
• Click “Simulate”
Frontend calls:
POST /simulate
Display:
• Stage shift
• Stress change
• Risk change
No local calculations allowed.
10.9 Service Marketplace UI
List providers sorted by backend ranking.
Display:
• Trust score
• Completion rate
• Complaint ratio
• Verified badge
Request Flow:
• Select provider
• Confirm request
• Status updates shown in timeline
No direct trust score manipulation on frontend.

===== PAGE 69 =====
10.10 Dashboard & Insights UI
Dashboard displays:
• Crop history
• Deviation trends
• Stress timeline chart
• Yield reports
• Service history
Charts generated using:
• Recharts or Chart.js
All data pulled from backend.
10.11 Accessibility Architecture
Accessibility must be global.
Accessibility Panel Options:
• High Contrast Mode
• Adjustable Text Size
• Reduced Motion Mode
• Brightness Simulation Mode
• Voice Input (optional advanced)
• Language Selection
Implementation:
Global Context:
<AccessibilityProvider>
<App/>
</AccessibilityProvider>
CSS variables used to toggle themes.
Must apply across all pages.

===== PAGE 70 =====
10.12 Language Support Strategy
Initial language:
• English
• Hindi (phase 1 optional)
Use i18n framework (next-i18next).
All text externalized.
No hardcoded strings.
10.13 Error Handling UX
Error states:
• API failure
• Replay in progress
• Dead letter state
• Rate limit exceeded
Display friendly messages:
Example:
“Timeline update in progress. Please wait.”
No technical stack traces shown.
10.14 Admin Panel UI
Separate route:
/admin
Features:
• Dead letter viewer
• Force replay
• Provider verification
• User suspension
• ML model toggle
• Event backlog monitor

===== PAGE 71 =====
Admin-only access enforced by role guard.
10.15 PWA Configuration
Enable:
• Service Worker
• Cache static assets
• Cache last known crop state
• Background sync capability
Manifest includes:
• App name
• Icons
• Theme color
Works as installable app.
10.16 Performance Optimization
Techniques:
• Lazy load heavy pages
• Dynamic import for charts
• Compress images before upload (optional)
• Limit re-renders
• Use React.memo for timeline components
10.17 State Management Strategy
Use:
• React Context for global state
• Local component state for UI
• Avoid Redux unless complexity increases
Keep frontend logic simple.

===== PAGE 72 =====
10.18 Security Considerations
• JWT stored in HttpOnly cookie
• No sensitive data in localStorage
• No direct API keys exposed
• Signed URL used for media
• Strict CORS config
10.19 Scalability to 600 Users
Frontend horizontally scalable:
• Static deployment via Vercel or Cloud Run
• No session server required
• Stateless requests
No redesign required.
10.20 Alignment with Lab Manual Phase 2
Lab Phase 2 requires:
• System architecture
• Database schema
• Security protocols
Lab_Manual
This section fulfills UI integration modeling requirement.
SECTION 11 — DEPLOYMENT &
INFRASTRUCTURE CONFIGURATION
11.1 Deployment Philosophy
CultivaX deployment must:

===== PAGE 73 =====
1. Be containerized.
2. Be stateless at application layer.
3. Use managed services for persistence.
4. Support horizontal scaling.
5. Separate compute roles (API vs Media Worker).
6. Be environment-separated (Dev / Prod).
7. Be CI/CD ready.
11.2 Infrastructure Overview (GCP)
Final Infrastructure Layout:
Internet
↓
Cloud Load Balancer
↓
Cloud Run (Backend API)
Cloud Run (Media Worker)
Cloud Run (Video Worker - optional)
↓
Cloud SQL (PostgreSQL)
Cloud Storage (Media Bucket)
↓
Cloud Scheduler
Cloud Logging
Cloud Monitoring
Secret Manager
This aligns directly with your locked tech stack
CultivaX Tech stack
.
11.3 Service Breakdown
Service 1 — Backend API
Purpose:
• Authentication
• CTIS
• SOE

===== PAGE 74 =====
• Event dispatcher
• ML inference (lightweight)
• Admin APIs
Cloud Run Configuration:
• CPU: 1 vCPU
• Memory: 512MB (start)
• Max instances: 5 (prototype)
• Concurrency: 80
• Min instances: 0 (cost saving)
Scaling to 600 users:
Increase max instances to 15–20.
Service 2 — Media Worker
Purpose:
• Image processing
• CNN inference
• Feature extraction
Cloud Run Configuration:
• CPU: 1–2 vCPU
• Memory: 1GB
• Concurrency: 1 (important)
• Max instances: 3 (prototype)
Heavy jobs isolated here.
Service 3 — Video Worker (Optional)
Purpose:
• Frame extraction
• Multi-frame aggregation
Higher memory allocation recommended (1.5–2GB).

===== PAGE 75 =====
11.4 Database Deployment
Service: Cloud SQL (PostgreSQL 15)
Configuration:
• Tier: db-f1-micro (prototype)
• Storage: 10GB SSD
• Automatic backups enabled
• Binary logging off (prototype)
Important Settings:
• max_connections tuned
• Connection pooling via SQLAlchemy
Scaling to 600 users:
• Upgrade instance tier
• Add read replica
• Increase storage
No schema redesign required.
11.5 Media Storage Configuration
Service: Cloud Storage
Bucket Settings:
• Private
• Uniform bucket-level access
• Lifecycle rule: delete after 90 days
• Object retention enforced via scheduler
Access:
• Signed URLs (15 min expiry)
No public read access.

===== PAGE 76 =====
11.6 Secrets Management
Use GCP Secret Manager.
Store:
• JWT secret
• DB credentials
• ML model version flag
• Feature flags
Cloud Run service account given access only to required secrets.
Never hardcode secrets.
11.7 Environment Separation
Two environments minimum:
• Development
• Production
Each has:
• Separate Cloud Run services
• Separate Cloud SQL instance
• Separate storage bucket
Environment variables define:
• DATABASE_URL
• ENVIRONMENT
• FEATURE_FLAGS
Prevents accidental data corruption.
11.8 CI/CD Pipeline
Lab Manual Phase 5 requires deployment + optional CI/CD
Lab_Manual
.

===== PAGE 77 =====
Recommended Pipeline:
GitHub → GitHub Actions → Build Docker → Push to Artifact Registry → Deploy to Cloud Run
Pipeline Steps:
1. On push to main:
a. Run lint
b. Run tests
c. Build Docker image
d. Push image
e. Deploy to Cloud Run
Frontend pipeline:
• Vercel auto-deploy OR
• Separate Cloud Run service
11.9 Version Control Requirements
All code must be in GitHub (Lab requirement)
Lab_Manual
.
Repository structure:
/backend
/frontend
/infrastructure
/docs
Branches:
• main
• dev
• feature/*
Commit discipline required.
11.10 Logging & Monitoring Integration
Cloud Logging:

===== PAGE 78 =====
Structured JSON logs:
• replay_duration
• event_id
• stress_delta
• trust_score_update
Cloud Monitoring:
Alerts:
• Event backlog > threshold
• Replay failure spike
• Media processing failure spike
• CPU > 80%
Helps during peer testing phase.
11.11 Health Check Configuration
Endpoint:
GET /api/v1/health
Checks:
• DB connection
• Dispatcher operational
• ML model loaded
• Worker connectivity
Cloud Run health check enabled.
11.12 Scaling Plan (100 → 600 Users)
Users Required Change
100 Default config
200 Increase max instances
400 Increase DB tier
600 Add read replica
1000+ Move to Pub/Sub
Because dispatcher abstraction exists

===== PAGE 79 =====
CultivaX Tech stack
, migration safe.
11.13 Backup & Disaster Recovery
Cloud SQL:
• Daily automatic backup
• Point-in-time recovery
Cloud Storage:
• 90-day lifecycle
• Optional backup bucket (future)
Code:
• GitHub repository backup
11.14 Security at Infrastructure Level
• HTTPS only
• IAM-based access
• Principle of least privilege
• Service account isolation
• No public DB access
• VPC connector if required
11.15 Cost Control Strategy
Prototype phase:
• Min instances = 0
• Low-tier DB
• Auto scale
• Log retention limited
Avoid over-provisioning.

===== PAGE 80 =====
11.16 Alignment with Lab Manual
Phase 5 requires:
• Deployment
• Hosting
• CI/CD
SECTION 12 — MONITORING, LOGGING
& OBSERVABILITY
12.1 Observability Philosophy
CultivaX must be:
• Deterministically debuggable
• Replay traceable
• ML behavior explainable
• Infrastructure observable
• Admin-auditable
Every important state change must be traceable.
No silent failures allowed.
12.2 Structured Logging Strategy
All backend logs must be JSON structured.
Example:
{
"timestamp": "ISO8601",
"service": "backend-api",
"event_type": "ReplayTriggered",
"crop_instance_id": "uuid",
"actions_processed": 12,
"replay_duration_ms": 48,
"stress_delta": 0.04,
"risk_index": 0.31,

===== PAGE 81 =====
"model_version": "risk_v1.0"
}
Fields must include:
• request_id
• user_id
• service_name
• severity_level
This enables Cloud Logging filtering.
12.3 Event-Level Traceability
Each event_log entry must include:
• event_id
• event_hash
• status
• retry_count
• processed_at
Replay logs must include:
• snapshot_used
• drift_applied
• stage_transition
• state_before
• state_after
This ensures replay corruption can be audited.
12.4 Request-Level Tracing
Each incoming API request:
• Generate request_id (UUID)
• Attach to context
• Include in all logs
Helps trace:
User action → Event → Replay → State change

===== PAGE 82 =====
Critical for peer testing phase
Lab_Manual
.
12.5 ML Observability
Every ML inference must log:
{
"model_name": "stress_cnn",
"model_version": "cnn_v1.0",
"confidence": 0.78,
"input_quality_score": 0.84,
"stress_probability": 0.42
}
Additionally track:
• Distribution of stress scores
• Drift in prediction averages
• Sudden confidence drops
If distribution shifts unexpectedly → alert.
12.6 Regional Learning Monitoring
Track:
• Number of samples per region
• Cluster size
• Adaptation offset applied
• Variance within cluster
If cluster size < threshold:
Log warning:
"Cluster size insufficient for adaptation."
Prevents silent overfitting.

===== PAGE 83 =====
12.7 Dispatcher Monitoring
Metrics to monitor:
• Pending event count
• DeadLetter count
• Average replay time
• Retry rate
• Processing latency
Alert thresholds:
• 50 pending events
• 5 dead letters in 1 hour
• Replay duration spike > 3x baseline
12.8 Media Worker Monitoring
Track:
• Media processing time
• Failure rate
• Frame extraction count
• Average confidence score
Alert if:
• Processing time > threshold
• Memory spike
• Worker crash loops
12.9 Infrastructure Health Monitoring
Use Cloud Monitoring dashboards for:
• CPU utilization
• Memory utilization
• DB connections
• Cloud Run instance count
• Request latency
• Error rate

===== PAGE 84 =====
Set alerts:
• CPU > 80% sustained
• Error rate > 5%
• DB connection exhaustion
12.10 Replay Debug Mode
Admin-only endpoint:
GET /admin/replay-debug/{crop_id}
Returns:
• Action list
• Snapshot used
• Stage transitions
• Drift adjustments
• Stress timeline
This is extremely powerful during viva.
12.11 Admin Observability Dashboard
Admin UI shows:
• Event backlog
• Dead letters
• ML drift indicator
• Regional cluster stats
• Worker health
• Replay performance graph
All pulled from logs + metrics.
12.12 Log Retention Policy
• Application logs: 30 days
• Error logs: 90 days
• ML training logs: permanent (small)

===== PAGE 85 =====
• Admin audit log: permanent (DB)
Avoid excessive storage costs.
12.13 Incident Response Model
If alert triggered:
1. Admin notified
2. Identify affected crop_instance
3. Inspect event chain
4. Force replay if safe
5. Log resolution
All incidents logged in admin_audit_log.
12.14 ML Drift Detection Strategy
Compute:
Rolling mean of stress_probability
If deviation > threshold:
• Mark ML model unstable
• Disable adaptation
• Alert admin
Never auto-disable silently.
12.15 Abuse Detection Logging
Monitor:
• Excessive override patterns
• Rapid action submissions
• Upload spam
• Repeated login failures
Log with severity="warning".

===== PAGE 86 =====
12.16 Performance Metrics Targets
Target metrics:
• Replay time < 100ms (normal)
• API latency < 300ms
• Media processing < 5 seconds/image
• Video processing < 30 seconds
These targets support scaling to 600 users
CultivaX Tech stack
.
12.17 Alignment with Lab Manual
Phase 4: Peer Testing
Lab_Manual
This section enables structured bug reporting.
Phase 5: Deployment & DevOps
Lab_Manual
Monitoring required.
This section fulfills both.
SECTION 13 — SECURITY
IMPLEMENTATION DETAILS

===== PAGE 87 =====
13.1 Security Design Philosophy
CultivaX security model is built around:
1. Integrity of crop timeline (CTIS protection)
2. Identity verification (JWT-based)
3. Event immutability
4. Least-privilege access
5. Media isolation
6. Abuse detection
7. Data minimization
8. Replay protection
Primary threat model:
Not nation-state actors.
But:
• Malicious user manipulation
• Timeline tampering
• Duplicate event injection
• Media abuse
• Trust score fraud
• ML poisoning
13.2 Threat Model Overview
Threat Categories:
1. Authentication attacks
2. Authorization bypass
3. Replay attack (duplicate action injection)
4. Data tampering
5. Media upload abuse
6. Trust score manipulation
7. ML dataset poisoning
8. Infrastructure misconfiguration
Each has mitigation defined below.

===== PAGE 88 =====
13.3 Authentication Security
Mechanisms:
• bcrypt password hashing
• Short-lived JWT (1 hour)
• JWT signature validation
• Account lockout after repeated failures
• Rate limiting on login
Secret key stored in Secret Manager.
Mitigation:
Prevents credential brute force and token forgery.
13.4 Authorization Security
Role-based access control (RBAC).
Enforced in middleware before controller logic.
Admin endpoints double-protected by:
• role check
• audit log requirement
Prevents horizontal privilege escalation.
13.5 Event Replay Protection
Critical area.
Each event:
• Has event_hash (SHA256)
• Unique constraint at DB level
• Idempotency-Key required in API
Mitigation:
• Prevents duplicate action injection
• Prevents network retry duplication

===== PAGE 89 =====
• Prevents offline replay duplication
Additionally:
Chain hash can be added for stronger immutability.
13.6 Input Validation Strategy
All API inputs validated via:
• Pydantic schema
• Type validation
• Range checks
• Enum enforcement
• String length limits
Reject:
• Invalid dates
• Stage skip attempts
• Stress > 1
• Negative land size
Prevents injection and state corruption.
13.7 SQL Injection Protection
• SQLAlchemy ORM
• Parameterized queries
• No raw string interpolation
Prevents SQL injection.
13.8 Media Upload Security
Threats:
• Malware upload
• Oversized file upload
• Storage abuse

===== PAGE 90 =====
• Direct bucket access
Mitigations:
• Signed URLs (15 min)
• Private bucket
• File type validation
• Size limits enforced
• Rate limiting
• No public bucket access
Optional future:
• Virus scanning
13.9 Offline Abuse Detection
Offline mode risk:
• Fabricated chronological manipulation
• Artificial action backdating
Mitigation:
• local_monotonic_counter required
• No backdating beyond tolerance window
• Server-side validation of action_effective_date
• Drift cap enforced
Prevents timeline fraud.
13.10 Trust Score Manipulation Prevention
Threats:
• Fake reviews
• Self-review
• Collusion
• Rating spikes
Mitigations:
• Unique review per request
• Cannot review before completion

===== PAGE 91 =====
• IP anomaly detection
• Sudden rating spike monitoring
• Exposure fairness algorithm
Admin review required for suspicious patterns.
13.11 ML Data Poisoning Mitigation
Threat:
User submits false yield data to manipulate regional learning.
Mitigation:
• Yield_verification_score
• Z-score validation
• Clamp ML_yield_value
• Minimum cluster size threshold
• Manual admin approval for model update
Prevents small-sample corruption.
13.12 Data Anonymization Policy
Before ML training:
Remove:
• user_id
• phone_number
• exact location coordinates
• device identifiers
Keep only:
• region
• crop_type
• anonymized patterns
No personally identifiable data used in training.

===== PAGE 92 =====
13.13 Infrastructure Security
• HTTPS enforced
• No public DB access
• IAM least privilege
• Separate service accounts
• Secret Manager usage
• Environment separation (Dev/Prod)
Prevents configuration-level attacks.
13.14 CSRF & CORS Protection
If JWT via cookie:
• CSRF token required
If JWT via header:
• Strict CORS policy
• Only frontend domain allowed
13.15 Rate Limiting
Limits:
• Login attempts
• Action submissions
• Media uploads
• Simulation requests
Prevents denial-of-service via user abuse.
13.16 Replay Integrity Check
Before replay execution:
• Validate action chronology
• Validate stage consistency

===== PAGE 93 =====
• Validate drift limits
If violation:
• Reject event
• Mark crop_instance RecoveryRequired
• Log incident
Prevents internal corruption.
13.17 Logging Security
Logs must:
• Avoid storing passwords
• Avoid storing raw JWT
• Avoid storing full media content
• Mask sensitive fields
Prevents accidental data leakage.
13.18 Disaster Recovery Security
Backups:
• Cloud SQL automated
• No direct DB shell access
• Controlled restore procedure
Ensures resilience.
13.19 Feature Flags & Kill Switches
Admin can:
• Disable ML
• Disable regional learning
• Disable media processing
Prevents runaway AI behavior.

===== PAGE 94 =====
13.20 Alignment with Lab Manual
Lab Phase 1 requires security in SRS
Lab_Manual
Lab Phase 2 requires security protocols in design
Lab_Manual
This section fully satisfies both.
SECTION 14 — PERFORMANCE &
SCALING PLAN
14.1 Performance Philosophy
CultivaX is:
• Logic-heavy
• State-sensitive
• ML-bounded
• Event-driven
The bottleneck is state integrity, not request throughput
CultivaX Tech stack
.
We optimize:
1. Replay efficiency
2. DB locking discipline
3. Worker isolation
4. Event backlog control

===== PAGE 95 =====
14.2 User Load Assumptions
Phase 1 (Prototype):
• 100–200 users
• Peak concurrent sessions: 20–30
• Average actions per user per day: 5
• Media uploads per day: 100 images
• Videos per day: 10–20
Phase 2 (Expanded):
• 500–600 users
• Peak concurrent sessions: 80–120
• Media uploads per day: ~300
• Videos per day: ~60
These are realistic agricultural app numbers.
14.3 Replay Complexity Analysis
Replay without snapshot:
O(n) where n = number of actions in crop_instance.
Worst case:
n = 200 actions per crop
With snapshot every 50 actions:
Replay complexity becomes:
O(k), where k ≤ 50
This guarantees bounded replay time.
Target replay time:
< 100ms (normal)
< 300ms (worst case)

===== PAGE 96 =====
14.4 API Performance Targets
Target metrics:
• Median latency: < 200ms
• 95th percentile latency: < 400ms
• Error rate: < 1%
Non-media endpoints must remain fast.
Heavy tasks offloaded to workers.
14.5 Media Worker Throughput
Image processing:
• Avg time per image: 2–5 seconds
• 100 images/day = trivial load
Video processing:
• Avg time per video: 15–30 seconds
• 20/day manageable with 2 worker instances
Scaling to 600 users:
• Increase worker max instances to 5
• Maintain concurrency=1 per instance
No redesign required.
14.6 Database Load Analysis
Estimated DB operations per action:
• Insert action_log
• Insert event_log
• Replay updates crop_instance
• Possibly snapshot insert
~5–10 DB queries per action
At 600 users:

===== PAGE 97 =====
600 × 5 actions/day = 3000 actions/day
3000 × 10 queries = 30,000 queries/day
This is extremely safe for PostgreSQL Cloud SQL.
Even micro-tier handles this comfortably.
14.7 Lock Contention Analysis
Row-level locking on:
crop_instance_id
Since each farmer operates on separate crop_instance:
No global lock contention.
Worst case:
Same user spams actions → sequential lock.
Safe.
14.8 Event Backlog Thresholds
Safe zone:
Pending events < 20
Warning:
50 events
Critical:
100 events
If backlog > 100:
• Increase Cloud Run max instances
• Reduce replay batch size
• Investigate DB bottleneck

===== PAGE 98 =====
14.9 Scaling Path (100 → 600 Users)
Stage Required Action
100 users Default config
200 users Increase max instances
400 users Increase DB tier
600 users Add DB read replica
1000+ users Migrate dispatcher to Pub/Sub
Dispatcher abstraction ensures migration safe
CultivaX Tech stack
.
14.10 Pub/Sub Migration Trigger Conditions
Move to distributed event bus if:
• Event backlog consistently > 500
• Replay latency > 1 second
• CPU > 80% sustained
• 1000 users
Migration steps:
1. Replace DBEventDispatcher
2. Publish to Pub/Sub
3. Dedicated subscriber service
4. Keep same replay logic
Zero business logic rewrite.
14.11 Horizontal Scaling Strategy
Cloud Run:
• Stateless containers
• Auto-scale on request load
• Independent scaling for:
o API
o Media Worker
o Video Worker

===== PAGE 99 =====
No sticky sessions required.
JWT stateless.
14.12 ML Scaling Strategy
Inference:
Lightweight — stays inside backend.
Training:
Run scheduled batch job.
If dataset grows large:
• Move training to separate container
• Possibly use Cloud Run job
Still no architectural redesign.
14.13 Memory & CPU Planning
Backend API:
• 512MB safe for 200 users
• 1GB for 600 users
Media worker:
• 1GB for image
• 2GB for video
Replay memory footprint small.
No in-memory queue accumulation.
14.14 Cost-Performance Tradeoff
Prototype:
• Min instances = 0

===== PAGE 100 =====
• Small DB tier
• Auto scale
Scaling:
Increase compute only when needed.
Avoid over-provisioning.
14.15 Failure Scenarios & Recovery
Scenario 1: Worker crash
• Event resets to Created
• Retry safe
Scenario 2: API crash
• Stateless → auto restart
• No state loss
Scenario 3: DB overload
• Increase tier
• Add read replica
Scenario 4: ML instability
• Kill switch enabled
System resilient by design.
14.16 Alignment with Lab Manual
Phase 5 requires deployment & hosting
Lab_Manual
.
This section demonstrates performance planning and scalability.
You can defend this mathematically in viva.

===== PAGE 101 =====
SECTION 15 — FUTURE EXPANSION &
RESEARCH ROADMAP
15.1 Architectural Evolution Path
Current architecture:
• Monolithic backend
• In-process event dispatcher
• Worker-based media isolation
• Snapshot-optimized replay
• ML embedded inside backend
Future Evolution Path:
Stage 1 → Extract Dispatcher
Stage 2 → Distributed Event Bus (Pub/Sub)
Stage 3 → Microservice separation (CTIS / SOE / ML)
Stage 4 → Multi-region deployment
Because dispatcher abstraction already exists
CultivaX Tech stack
, migration requires config change, not redesign.
15.2 Advanced ML Expansion
Current ML:
• Logistic regression (risk)
• CNN (stress)
• Cluster-based regional learning
Future Enhancements:
1. XGBoost for risk modeling
2. Time-series LSTM for growth forecasting
3. Segmentation-based leaf disease detection
4. Autoencoder-based anomaly detection
5. Bayesian confidence calibration

===== PAGE 102 =====
All possible without breaking CTIS because:
ML influence remains bounded and confidence-weighted.
15.3 Yield Modeling Research Extension
Current:
• Yield stored
• Z-score verification
• ML_yield_value separated from reported_yield
Future:
• Predictive yield modeling
• Early yield forecast at sowing
• Climate-adjusted yield prediction
• Scenario-based probabilistic yield intervals
Important:
Never promise exact yield.
Always probabilistic.
15.4 Regional Intelligence Deepening
Current:
• Regional cluster adaptation
• Minimum sample threshold
Future:
• Sub-district adaptation layers
• Microclimate modeling
• Geo-spatial crop similarity mapping
• Spatial clustering using coordinates (anonymized)
Global model remains stable.
Regional model adaptive.
This preserves research defensibility.

===== PAGE 103 =====
15.5 Edge AI Expansion
Current:
• Backend CNN
• Edge optional (TFLite ready)
Future:
• Fully offline stress detection
• On-device model updates
• Low-bandwidth rural deployment
• Federated learning integration
Federated learning possible because:
• ML separated from CTIS core
• Feature schema versioned
15.6 IoT Integration Roadmap
Future expansion:
• Soil moisture sensors
• Temperature sensors
• NDVI satellite integration
• Automated irrigation triggers
Integration method:
IoT device → API → Event → ReplayTriggered
No structural redesign required.
15.7 Marketplace Expansion
Current SOE:
• Equipment
• Labor
• Inputs
Future:

===== PAGE 104 =====
• Credit scoring integration
• Insurance integration
• Micro-financing
• Cold storage partnerships
• Logistics network
Trust system extensible to:
• Financial reliability scoring
• Service guarantee scoring
15.8 Multi-Tenant SaaS Model
Future:
• District-level deployments
• NGO partnerships
• Government collaboration
• Institutional dashboards
Because system is stateless and cloud-native
CultivaX Tech stack
, multi-tenant extension possible.
15.9 Advanced Visualization Roadmap
Current:
• 2D dynamic SVG crop
Future:
• 3D WebGL-based visualization
• Realistic crop growth modeling
• Augmented Reality field overlay
• Satellite imagery overlay
Must remain optional.
Core logic unaffected.

===== PAGE 105 =====
15.10 Blockchain / Ledger (Deferred Feature)
You intentionally deferred:
Zero-Knowledge Ledger
Future:
• Cryptographic crop history proof
• Immutable yield history
• Bank-readable farm performance proof
Can be implemented as:
Event chain hash → anchored externally.
No core redesign required.
15.11 Distributed Architecture Roadmap
When user base > 1000:
• Replace dispatcher with Pub/Sub
• Separate CTIS service
• Separate SOE service
• Dedicated ML service
• Introduce Redis cache
• Use CDN for media
Because abstraction boundaries already exist
CultivaX Tech stack
, migration safe.
15.12 Research Publication Opportunities
Potential research themes:
1. Time-centric agricultural modeling
2. Human-AI co-pilot systems
3. Drift-bounded adaptive crop modeling
4. Replay-based deterministic agricultural systems

===== PAGE 106 =====
5. Multi-signal stress arbitration
6. Regional adaptive agronomy intelligence
Your architecture supports publishable work.
15.13 Startup Readiness Roadmap
Phase 1: Academic prototype
Phase 2: Pilot in one district
Phase 3: NGO partnership
Phase 4: Marketplace monetization
Phase 5: SaaS model
Revenue models possible:
• Commission on services
• Premium analytics
• Subscription advisory
• Institutional dashboards
Architecture supports monetization.
15.14 Compliance & Policy Expansion
Future:
• Data governance compliance
• Government integration
• Agricultural data privacy compliance
• Digital agriculture mission alignment
Because system logs structured and anonymizes ML data, compliance feasible.
15.15 Final Architectural Maturity Summary
CultivaX now:
✔ Deterministic replay core
✔ Event-driven architecture
✔ ML bounded and versioned

===== PAGE 107 =====
✔ Worker-isolated media
✔ Marketplace fairness
✔ Cloud-native deployment
✔ Migration-ready abstraction
✔ Offline-safe design
✔ Accessibility-ready frontend
✔ Lab-compliant lifecycle alignment
\end{verbatim}
\section{UML}
\begin{verbatim}
# FILE: CultivaX UML.pdf
# PAGES: 123



===== PAGE 1 =====
CultivaX UML
SECTION 1 — Modeling Foundations
1.1 Purpose of UML Modeling
This UML document provides a structural and behavioral representation of CultivaX, covering:
• Functional system interactions (Use Case view)
• Domain object relationships (Class view)
• Runtime interactions (Sequence view)
• Lifecycle behavior (State Machines)
• Workflow logic (Activity view)
• System composition (Component view)
• Physical deployment (Deployment view)
The modeling approach respects:
• Event-driven state mutation
• Deterministic replay architecture
• Partition isolation per CropInstance
• ML assistive (non-authoritative) behavior
• Service Orchestration isolation
1.2 Modeling Philosophy
CultivaX is not a CRUD system.
It is:
• A time-aware state machine
• Event-driven
• Replay-deterministic
• ML-augmented but rule-governed
• Partition-isolated
Therefore:
• All state changes originate from Events.
• No object mutates itself silently.
• ReplayEngine is modeled explicitly.
• Dispatcher is a first-class architectural element.

===== PAGE 2 =====
• SOE and CTIS are isolated subsystems.
SECTION 2 — System-Level Use Case
Modeling
We begin with the Primary System Use Case Diagram.
This diagram answers:
Who interacts with the system, and what major capabilities exist?
🖼 System-Level Use Case Diagram
Now here is the logical structure we will model for CultivaX:
Actors
Primary Actors
• Farmer
• Service Provider
• Admin
Secondary / External Actors
• Weather API
• ML Engine
• Edge AI Module
Major Farmer Use Cases
• Register / Login
• Create Crop Instance
• View Crop Timeline
• Log Action
• Upload Media

===== PAGE 3 =====
• Run What-If Simulation
• Submit Yield
• Request Service
• Submit Review
• View Dashboard
System Internal Use Cases (Triggered)
• Trigger Replay
• Recalculate Stress
• Compute Risk Index
• Create Snapshot
• Update Regional Cluster
• Calculate Trust Score
• Handle Offline Sync
These are system-driven but shown for architectural clarity.
Admin Use Cases
• Modify Crop Rules
• View Dead Letter Events
• Suspend Provider
• Adjust Trust Weights
• Trigger Manual Replay
• View System Logs
Provider Use Cases
• Register as Provider
• Manage Equipment
• Accept Service Request
• Update Availability
• Respond to Escalation

===== PAGE 4 =====
Important Modeling Decisions
We will:
• Use <<include>> for:
o Log Action → Trigger Replay
o Submit Yield → Yield Verification
o Upload Media → Media Analysis
o Service Completion → Trust Recalculation
• Use <<extend>> for:
o AtRisk Warning extends View Timeline
o Escalation extends Service Completion
• Show system boundary labeled:
CultivaX Platform
• Show ML Engine as secondary actor to:
o Media Analysis
o Risk Evaluation
SECTION 2.1 — Master System Overview Use
Case Diagram
This diagram answers:
Who interacts with CultivaX at a high level?
It will not show internal replay or ML mechanics yet — only major capability groups.
Actors (High-Level)
Primary:
• Farmer
• Service Provider
• Admin
Secondary:
• Weather API
• ML Engine

===== PAGE 5 =====
High-Level Capability Groups
For clarity, we group use cases into 4 capability clusters:
1. Crop Lifecycle Management
2. Media & Intelligence Processing
3. Service Orchestration
4. Governance & System Control
🖼 Master Overview – Visual Style Reference
Logical Structure of the Master Diagram
System Boundary:
“CultivaX Platform”
Inside boundary (grouped):
Crop Lifecycle Management
• Create Crop Instance
• Manage Timeline
• Log Action
• Submit Yield
• Run What-If Simulation
• View Dashboard
Actor:
Farmer
Media & Intelligence
• Upload Image/Video
• Analyze Crop Stress
• Compute Risk Index
• Update Regional Learning
Actors:
Farmer

===== PAGE 6 =====
ML Engine
Weather API
Service Orchestration
• Request Service
• Manage Equipment
• Accept Service Request
• Submit Review
• Calculate Trust Score
Actors:
Farmer
Service Provider
🏛 Governance & Control
• Modify Crop Rules
• Suspend Provider
• Adjust Trust Weights
• View System Logs
• Handle Dead Letters
Actor:
Admin
Architectural Clarity Decisions
We will:
• Keep internal “Replay Engine” hidden at master level.
• Show ML Engine as assisting “Analyze Crop Stress”.
• Show Weather API assisting “Compute Risk Index”.
• Show Service Provider only interacting with SOE-related use cases.
• Not mix governance with farmer flows.

===== PAGE 7 =====
SECTION 2.2 — Farmer Subsystem Use Case
Diagram (Detailed)
This diagram focuses only on:
Actor: Farmer
It will model:
• Crop lifecycle control
• Timeline mutation
• Media interaction
• Yield submission
• Service initiation
• What-if simulation
• Dashboard analytics
And it will explicitly show:
• <<include>> for mandatory internal triggers
• <<extend>> for conditional flows
• Event-driven replay logic
🖼 Farmer Subsystem – Use Case Diagram
CultivaX – Farmer Subsystem Logical
Structure
Primary Actor
Farmer

===== PAGE 8 =====
Core Crop Lifecycle Use Cases
Create Crop Instance
• <<include>> Initialize Timeline
• <<include>> Assign Rule Version
Log Action
• <<include>> Persist Action Log
• <<include>> Trigger Replay
• <<include>> Update Snapshot (conditional)
• <<extend>> Raise AtRisk Alert
View Crop Timeline
• <<extend>> Display Risk Warning
• <<extend>> Display Drift Notification
Upload Media
• <<include>> Analyze Media
• <<include>> Update Stress Score
• <<include>> Trigger Replay
Secondary actor:
• ML Engine
Submit Yield
• <<include>> Verify Yield
• <<include>> Update Regional Cluster
• <<include>> Trigger Replay
• <<extend>> Cap ML Yield Value

===== PAGE 9 =====
Run What-If Simulation
• <<include>> Clone Crop State
• <<include>> Execute Replay (Isolated)
• <<include>> Generate Projected Outcome
No state mutation allowed.
Request Service
• Connects to SOE subsystem
• <<include>> Fetch Ranked Providers
View Dashboard
• <<include>> Aggregate Crop Metrics
• <<include>> Load Historical Insights
Architectural Notes
Why Trigger Replay is included under multiple use cases?
Because in CultivaX:
Any structural or stage-affecting action must result in deterministic recomputation.
This enforces:
• No silent mutation
• Event-driven integrity
• Snapshot optimization
Why Yield Verification is included?
Because of:
• Farmer Truth vs ML Truth separation
• Biological cap enforcement
• Regional cluster stability constraints

===== PAGE 10 =====
Why What-If is isolated?
Because:
• It must not emit events
• It must not persist mutation
• It must deep-clone instance state
Implicit Invariants Modeled
The diagram reflects:
• Chronological ordering
• Partition isolation
• ML assistive (not authoritative)
• Replay atomicity
SECTION 2.3 — Service Orchestration (SOE)
Subsystem Use Case Diagram
This diagram focuses on:
Actors:
• Farmer
• Service Provider
• Admin
It explicitly models:
• Provider lifecycle
• Trust computation
• Exposure fairness
• Trust decay
• Complaint escalation
• Fraud detection
• Isolation from CTIS

===== PAGE 11 =====
SOE Subsystem – Structural Overview
System Boundary:
“CultivaX – Service Orchestration Ecosystem”
Actors
Farmer
• Request Service
• Submit Review
• File Complaint
• View Ranked Providers
Service Provider
• Register as Provider
• Manage Equipment
• Update Availability
• Accept / Decline Request
• Respond to Escalation
Admin
• Verify Provider
• Suspend Provider
• Override Trust Score
• Resolve Escalation
• Moderate Reviews
Core SOE Use Cases
Register as Provider
• <<include>> Create Provider Profile
• <<include>> Set Initial Trust Score (0.5)
• <<include>> Await Admin Verification

===== PAGE 12 =====
Manage Equipment
• Add Equipment
• Update Availability
• Flag Maintenance Status
Request Service (Farmer)
• <<include>> Fetch Eligible Providers
• <<include>> Apply Exposure Fairness Algorithm
• <<include>> Rank by Trust Score
• <<extend>> Highlight Emergency Services (if AtRisk)
Note:
This interaction does NOT modify crop state.
Accept Service Request (Provider)
• <<include>> Lock Service Slot
• <<include>> Update Request Status
Complete Service Request
• <<include>> Enable Review Submission
• <<include>> Trigger Trust Recalculation
Submit Review (Farmer)
• <<include>> Validate Request Completion
• <<include>> Update Complaint Ratio
• <<include>> Trigger Trust Recalculation
• <<extend>> Escalate If Severity High

===== PAGE 13 =====
Trust & Fairness Use Cases (Explicit)
Calculate Trust Score
• <<include>> Compute Completion Rate
• <<include>> Compute Complaint Ratio
• <<include>> Compute Rating Stability
• <<include>> Apply Escalation Penalty
• <<include>> Apply Verified Bonus
• <<include>> Apply Trust Decay Mechanism
• <<include>> Clamp to [0,1]
Apply Exposure Fairness
• <<include>> Calculate Exposure Penalty
• <<include>> Adjust Ranking Score
• <<include>> Enforce Rotation Constraint
• <<include>> Prevent Monopoly Bias
This is algorithmically separate from trust_score.
Fraud Detection
• <<include>> Detect Review Pattern Anomaly
• <<include>> Detect Sudden Rating Spike
• <<include>> Detect IP Correlation
• <<extend>> Flag Provider for Admin Review
Escalation Workflow
Escalate Complaint
Triggered when:

===== PAGE 14 =====
• Complaint severity threshold exceeded
• Complaint ratio threshold exceeded
Flow:
• <<include>> Auto-Flag Provider
• <<include>> Notify Admin
• <<extend>> Temporarily Reduce Exposure
🏛 Admin Governance Use Cases
Verify Provider
• Approve / Reject
• Change verification_status
Suspend Provider
• Soft delete provider
• Hide from listings
Override Trust Score
• Manually adjust
• Log in AdminAuditLog
Isolation Guarantee (Important Modeling
Element)
We explicitly annotate:
⚠ SOE must NOT:
• Trigger Replay

===== PAGE 15 =====
• Modify CropInstance
• Modify Stress Score
• Modify Risk Index
This will be shown in the diagram via constraint note:
{SOE cannot mutate CTIS state}
SECTION 2.4 — Admin & Governance
Subsystem Use Case Diagram
System Boundary:
“CultivaX – Governance & Control Layer”
Actor: Admin
This diagram will include:
• Rule template lifecycle
• ML model governance
• Replay intervention
• Dead letter handling
• Trust override
• System kill switches
• Audit logging
ADMIN USE CASE STRUCTURE
Crop Rule Governance
Create Rule Template Draft
• <<include>> Define Stage Definitions
• <<include>> Define Drift Constraints
• <<include>> Define Risk Parameters

===== PAGE 16 =====
Validate Rule Template
• <<include>> Run Simulation Replay
• <<include>> Check Stage Continuity
• <<include>> Check Drift Bounds
• <<include>> Detect Infinite Loop Risk
Activate Rule Template
• <<include>> Set is_active = TRUE
• <<include>> Log Governance Event
• <<extend>> Notify Future Instances
Constraint:
⚠ Must NOT retroactively mutate existing CropInstances.
ML Model Governance
Register New ML Model Version
• <<include>> Store Model Metadata
• <<include>> Store Evaluation Metrics
• <<include>> Store Feature Schema Version
Activate ML Model Version
• <<include>> Mark Model Active
• <<include>> Log Model Activation Event
• <<extend>> Disable Previous Version
Constraint:
Only new CropInstances use new model unless manual replay triggered.

===== PAGE 17 =====
Monitor Model Drift
• <<include>> Compare Performance Metrics
• <<include>> Detect Precision Drop
• <<include>> Detect False Positive Increase
• <<extend>> Trigger Model Rollback Suggestion
Rollback ML Model
• <<include>> Restore Previous Active Version
• <<include>> Log Rollback Event
Replay & Recovery Governance
View Dead Letter Events
• <<include>> Inspect Failure Reason
• <<include>> View Event Payload
Retry Dead Letter Event
• <<include>> Reset Event Status
• <<include>> Requeue Event
Lock CropInstance
• <<include>> Set State = RecoveryRequired
• <<include>> Prevent Further Actions
Trigger Manual Replay
• <<include>> Enqueue ReplayTriggered Event
• <<include>> Acquire Partition Lock

===== PAGE 18 =====
⚙ System Kill Switch Controls
These are critical architectural safeguards.
Disable ML Influence
• <<include>> Toggle Feature Flag
• <<include>> Log System Configuration Change
Disable Regional Learning
• <<include>> Prevent Cluster Updates
Disable SOE Recommendations
• <<include>> Suppress Service Suggestions
Constraint:
Core CTIS must remain operational.
🛡 Trust & Marketplace Governance
Suspend Provider
• <<include>> Set is_suspended = TRUE
• <<include>> Remove from Exposure Ranking
• <<include>> Log Audit Entry
Override Trust Score
• <<include>> Manually Set Trust Value
• <<include>> Log Override Reason

===== PAGE 19 =====
Adjust Trust Weights
• <<include>> Modify Weight Parameters
• <<include>> Trigger Trust Recalculation
Audit & Compliance
View Audit Logs
• <<include>> Filter by Entity
• <<include>> Filter by Date Range
Export Governance Report
• <<include>> Compile Rule Changes
• <<include>> Compile Model Changes
• <<include>> Compile Suspensions
Constraint:
AdminAuditLog is immutable.
Architectural Integrity Constraints (Explicit
Notes)
We will annotate diagram with constraints:
{No admin action may bypass Event Dispatcher}
{Rule changes are prospective only}
{ML activation must preserve reproducibility}
{DeadLetter events cannot auto-resolve}

===== PAGE 20 =====
SECTION 2.5 — Internal System Processes
Use Case Diagram
System Boundary:
“CultivaX – Core Processing Engine”
This diagram models:
• Event Dispatcher
• Replay Engine
• Snapshot Manager
• Stress Integration
• Risk Computation
• Event Lifecycle
• Offline Sync
• Partition Locking
It connects to external triggers from:
• Farmer
• Admin
• SOE
• ML Engine
Actor Continuity
From Farmer Diagram:
• Log Action
• Upload Media
• Submit Yield
• Run What-If
• Offline Sync
From SOE Diagram:
• Service Completed (indirect trigger)
From Admin Diagram:
• Trigger Manual Replay
• Retry Dead Letter

===== PAGE 21 =====
• Activate Rule
• Activate ML Model
Core Internal Use Cases
Publish Event
Triggered by:
• Log Action
• Upload Media
• Submit Yield
• Admin Actions
Includes:
• Generate Event Hash
• Persist Event
• Set Status = Created
Constraint:
⚠ Must be idempotent
Process Event Queue
Includes:
• Select Created Events
• Acquire Partition Lock
• Set Status = Processing
Constraint:
⚠ FIFO per crop_instance_id

===== PAGE 22 =====
Trigger Replay
Triggered by:
• ActionLogged
• MediaAnalyzed
• YieldSubmitted
• ManualReplay
Includes:
• Acquire Row Lock
• Load Snapshot (if exists)
• Load Ordered Actions
• Validate Chronological Invariants
• Apply Actions Sequentially
• Integrate Stress Signals
• Compute Risk Index
• Enforce Drift Constraints
• Update CropInstance
• Create Snapshot (conditional)
Constraint:
⚠ Atomic transaction boundary
Handle Offline Sync
Triggered by:
• OfflineSyncCompleted
Includes:
• Insert Unsynced Actions
• Sort by (effective_date, monotonic_counter)
• Trigger Replay
• Debounce Multiple Replay Events
Constraint:
⚠ No partial state exposure

===== PAGE 23 =====
Compute Stress Integration
Triggered during Replay
Includes:
• Combine Backend ML Signal
• Combine Weather Risk
• Combine Deviation Penalty
• Combine Edge AI Signal
• Apply Confidence Weight
• Apply Exponential Smoothing
Constraint:
stress ∈ [0,1]
Compute Risk Index
Includes:
• Combine Weather Risk
• Combine Farmer Input Risk
• Apply Stage Sensitivity
• Bound Output [0,1]
Constraint:
⚠ Must not directly alter stage
Snapshot Management
Includes:
• Determine Snapshot Threshold
• Persist Snapshot State
• Index by snapshot_index
Constraint:
⚠ Structural Action invalidates snapshot

===== PAGE 24 =====
Dead Letter Handling
Triggered when:
• retry_count > threshold
Includes:
• Move to DeadLetter
• Notify Admin
• Set CropInstance = RecoveryRequired
Constraint:
⚠ Must not block other partitions
Model Activation Propagation
Triggered by:
• Admin Activate Model
Includes:
• Mark Model Active
• Apply to New CropInstances
• Preserve Historical Model Version
Constraint:
⚠ No retroactive mutation
Critical Architectural Constraints (Explicit
Notes)
We annotate:
{No state mutation without event_log entry}
{Replay is atomic per partition}
{No cross-crop event interleaving}
{DeadLetter must not cascade failure}
{What-If must not persist mutation}

===== PAGE 25 =====
SECTION 3 — DOMAIN CLASS
DIAGRAMS (LAYERED)
We begin with:
3.1 Identity & Core Ownership Layer
This layer models:
• User hierarchy
• Role separation
• Ownership boundaries
• Governance linkage
Identity & Ownership — Class Diagram

===== PAGE 26 =====
Now the actual CultivaX structural modeling:
Class Structure (CultivaX Identity Layer)
Abstract Class: User
<<abstract>>
User
-------------------------
user_id : UUID
full_name : String
phone_number : String
role : Enum
is_active : Boolean
created_at : Timestamp
deleted_at : Timestamp
Farmer extends User
Farmer
-------------------------
preferred_language : String
accessibility_settings : JSON
Relationship:
User <|-- Farmer

===== PAGE 27 =====
ServiceProvider extends User
ServiceProvider
-------------------------
trust_score : Float
verification_status : Enum
complaint_count : Int
completion_count : Int
is_suspended : Boolean
Relationship:
User <|-- ServiceProvider
🏛 Admin extends User
Admin
-------------------------
admin_level : Enum
Relationship:
User <|-- Admin
Ownership Relationships
Farmer owns CropInstances
Farmer 1 ──── 0..* CropInstance
Composition:
Farmer ◆──── CropInstance
(If Farmer deleted → CropInstances soft-archived)
ServiceProvider owns Equipment
ServiceProvider 1 ──── 0..* Equipment
Composition:
ServiceProvider ◆──── Equipment

===== PAGE 28 =====
ServiceProvider participates in ServiceRequest
ServiceProvider 1 ──── 0..* ServiceRequest
Farmer 1 ──── 0..* ServiceRequest
Association only (not composition).
Governance Link
Admin interacts with:
• CropRuleTemplate
• MLModel
• ServiceProvider
• EventLog
Association (not ownership).
Why This Layer Matters
This layer defines:
✔ Role-based boundaries
✔ Ownership constraints
✔ Soft-delete cascade logic
✔ Governance authority
✔ Separation of concerns
This matches your DB schema exactly.
3.2 CTIS Core Domain Layer
Now we model the deterministic engine itself.
This is the heart of CultivaX.

===== PAGE 29 =====
CTIS Core — Class Diagram
Core Domain Classes
CropInstance
CropInstance
-------------------------
crop_instance_id : UUID
crop_type : String
region : String
sowing_date : Date
baseline_stage : String
current_stage : String
stress_score : Float
risk_index : Float
state : Enum
rule_version_applied : Int
event_chain_hash : String
ActionLog
ActionLog
-------------------------
action_id : UUID
action_type : String
action_impact_type : Enum
action_effective_date : Date
local_monotonic_counter : Long
metadata : JSON
Composition Relationship
CropInstance ◆──── 0..* ActionLog
Meaning:
ActionLogs cannot exist without CropInstance.

===== PAGE 30 =====
Snapshot
Snapshot
-------------------------
snapshot_id : UUID
snapshot_index : Int
stage : String
baseline_stage : String
stress_score : Float
risk_index : Float
Relationship:
CropInstance ◇──── 0..* Snapshot
Aggregation (Snapshot tied but not identity-defining).
DeviationProfile
DeviationProfile
-------------------------
consecutive_deviation_count : Int
deviation_trend_slope : Float
recurring_pattern_flag : Boolean
Composition:
CropInstance ◆──── 1 DeviationProfile
YieldRecord
YieldRecord
-------------------------
reported_yield : Float
ml_yield_value : Float
yield_verification_score : Float
Association:
CropInstance 1 ──── 0..1 YieldRecord

===== PAGE 31 =====
Engine & Service Classes (Behavioral)
These are service-layer domain components.
ReplayEngine
ReplayEngine
-------------------------
replay(crop_instance_id)
validate_invariants()
apply_action()
compute_stress()
compute_risk()
create_snapshot()
Dependency:
ReplayEngine ..> CropInstance
ReplayEngine ..> ActionLog
ReplayEngine ..> Snapshot
ReplayEngine ..> DeviationProfile
StressEngine
StressEngine
-------------------------
integrate_signals()
apply_smoothing()
bound_stress()
Used by ReplayEngine.
RiskCalculator
RiskCalculator
-------------------------
compute_risk_index()
apply_stage_weight()
bound_probability()

===== PAGE 32 =====
Critical Constraints (Modeled as Notes)
{No CropInstance mutation without Event}
{Replay must be atomic}
{Drift bounded per rule template}
SECTION 3.3 — Event System Layer (Class
Diagram)
This layer sits between:
• Farmer/Admin/SOE triggers
• CTIS replay engine
• ML layer
It ensures safe state mutation.
Core Event Domain Classes
Event
Event
-------------------------
event_id : UUID
event_type : String
crop_instance_id : UUID?
event_payload : JSON
event_hash : String
chain_hash : String
status : Enum
retry_count : Int
failure_reason : String
priority : Int
created_at : Timestamp
processed_at : Timestamp

===== PAGE 33 =====
Associations
CropInstance 1 ──── 0..* Event
Events are associated with CropInstance via partition key.
(Not composition — event persists even if crop archived)
Dispatcher Abstraction
EventDispatcherInterface (Interface)
<<interface>>
EventDispatcherInterface
-------------------------
publish(event)
process_pending_events()
This enables migration to Pub/Sub later.
DBEventDispatcher (Concrete Implementation)
DBEventDispatcher
-------------------------
publish()
process_pending_events()
acquire_partition_lock()
set_processing_status()
handle_retry()
move_to_dead_letter()
Relationship:
EventDispatcherInterface <|.. DBEventDispatcher
Implementation (dashed triangle arrow).

===== PAGE 34 =====
EventHandler
EventHandler
-------------------------
handle(event)
Specializations:
• ReplayEventHandler
• MediaAnalyzedHandler
• YieldSubmittedHandler
• RuleModifiedHandler
• TrustRecalculationHandler
Relationship:
DBEventDispatcher ..> EventHandler
EventHandler ..> ReplayEngine
Dependency arrow.
DeadLetterQueue
DeadLetterQueue
-------------------------
store_failed_event()
retry_event()
Association:
Event 0..1 ──── 1 DeadLetterQueue
(Conceptually moved when status = DeadLetter)
Partition & Locking Model
We model logical partition as constraint:
{Partition Key = crop_instance_id}
{FIFO per partition}
{No cross-partition blocking}

===== PAGE 35 =====
Row-level locking modeled as note attached to DBEventDispatcher:
{SELECT FOR UPDATE SKIP LOCKED}
Event Lifecycle State Modeling (Class-Level
Constraints)
{Created → Processing → Processed}
{Retry < threshold → Created}
{Retry ≥ threshold → DeadLetter}
This will later connect to State Machine diagram.
Chain Hash Integrity
Constraint note attached to Event:
{chain_hash_n = hash(previous_hash + event_hash)}
{Tamper detection enforced during replay}
Idempotency Enforcement
Constraint attached to publish():
{event_hash UNIQUE}
{Duplicate publish rejected}
How This Connects to CTIS
Dependency flow:
EventDispatcher ..> ReplayEngine
ReplayEngine ..> CropInstance
Meaning:
Events are the only gateway to state mutation.

===== PAGE 36 =====
🏗 Structural Summary of Event Layer
Classes:
• Event
• EventDispatcherInterface
• DBEventDispatcher
• EventHandler (abstract)
• Specific Event Handlers
• DeadLetterQueue
Relationships:
• Interface implementation
• Dependency injection
• Partitioned association
• Idempotent uniqueness
• Dead-letter fallback
SECTION 3.4 — SOE Domain Layer (Class
Diagram)
This layer models:
• Provider identity extension
• Equipment management
• Service lifecycle
• Trust score computation
• Exposure fairness control
• Fraud detection
• Escalation policy
• Admin override linkage
Core SOE Domain Entities

===== PAGE 37 =====
ServiceProvider
ServiceProvider
-------------------------
provider_id : UUID
verification_status : Enum
trust_score : Float
completion_rate : Float
complaint_ratio : Float
is_suspended : Boolean
created_at : Timestamp
deleted_at : Timestamp
Inheritance (from Identity Layer):
User <|-- ServiceProvider
Equipment
Equipment
-------------------------
equipment_id : UUID
equipment_type : String
availability_status : Enum
hourly_rate : Float
is_flagged : Boolean
Composition:
ServiceProvider ◆──── 0..* Equipment
Meaning:
Equipment lifecycle bound to provider.
ServiceRequest
ServiceRequest
-------------------------
request_id : UUID
service_type : String
status : Enum
region : String
created_at : Timestamp
updated_at : Timestamp

===== PAGE 38 =====
Associations:
Farmer 1 ──── 0..* ServiceRequest
ServiceProvider 1 ──── 0..* ServiceRequest
CropInstance 0..1 ──── 0..* ServiceRequest
Not composition — requests can persist for audit.
ServiceReview
ServiceReview
-------------------------
review_id : UUID
rating : Int
complaint_category : String
severity_level : Int
resolution_status : Enum
created_at : Timestamp
Association:
ServiceRequest 1 ──── 0..1 ServiceReview
Constraint:
{One review per completed request}
Trust & Marketplace Engines
These are behavioral service-layer classes.
⚖ TrustScoreEngine
TrustScoreEngine
-------------------------
compute_completion_rate()
compute_complaint_ratio()
compute_rating_stability()
apply_verified_bonus()
apply_escalation_penalty()
apply_trust_decay()
clamp_score()

===== PAGE 39 =====
Dependency:
TrustScoreEngine ..> ServiceProvider
TrustScoreEngine ..> ServiceReview
TrustScoreEngine ..> ServiceRequest
Important:
TrustScoreEngine does NOT rank providers.
It only computes trust_score.
ExposureFairnessEngine
ExposureFairnessEngine
-------------------------
calculate_exposure_penalty()
adjust_ranking_score()
enforce_rotation_limit()
prevent_monopoly_bias()
Dependency:
ExposureFairnessEngine ..> ServiceProvider
Critical architectural separation:
{RankingScore = TrustScore - ExposurePenalty}
Trust and exposure are separate concerns.
EscalationPolicy
EscalationPolicy
-------------------------
check_complaint_threshold()
auto_flag_provider()
notify_admin()
reduce_exposure()
Dependency:
EscalationPolicy ..> ServiceProvider
EscalationPolicy ..> ServiceReview

===== PAGE 40 =====
🕵 FraudDetector
FraudDetector
-------------------------
detect_rating_spike()
detect_ip_correlation()
detect_pattern_anomaly()
flag_suspicious_provider()
Dependency:
FraudDetector ..> ServiceReview
FraudDetector ..> ServiceProvider
Fraud detection does NOT automatically suspend — only flags.
🏛 Governance Linkage
Admin interacts with:
Admin ──── ServiceProvider
Admin ──── TrustScoreEngine
Admin ──── EscalationPolicy
Admin may:
• Override trust score
• Suspend provider
• Adjust weight parameters
Constraint note:
{All admin overrides logged in AdminAuditLog}
CTIS Isolation Constraint (Very Important)
We explicitly annotate:
{SOE must NOT modify CropInstance}
{SOE must NOT trigger Replay}
{SOE must NOT mutate stress_score or risk_index}
Association to CropInstance is informational only.

===== PAGE 41 =====
SECTION 3.5 — ML & Statistical
Architecture Layer (Class Diagram)
This layer is assistive, not authoritative.
It integrates with:
• ReplayEngine
• StressEngine
• RiskCalculator
• Admin governance
Core ML Domain Entities
MLModel (Abstract Base)
<<abstract>>
MLModel
-------------------------
model_id : UUID
model_name : String
model_version : String
training_date : Timestamp
metrics : JSON
is_active : Boolean
RiskPredictor (Logistic Regression)
RiskPredictor
-------------------------
predict_risk(features)
return_probability()
return_confidence()
Inheritance:
MLModel <|-- RiskPredictor

===== PAGE 42 =====
Dependency:
RiskPredictor ..> CropInstance
RiskPredictor ..> StressEngine
StressCNN (Backend ML)
StressCNN
-------------------------
analyze_image()
analyze_video_frames()
return_stress_probability()
return_confidence_score()
Inheritance:
MLModel <|-- StressCNN
Dependency:
StressCNN ..> MediaFile
EdgeAIModel
EdgeAIModel
-------------------------
classify_image()
estimate_stress()
return_confidence()
Inheritance:
MLModel <|-- EdgeAIModel
Constraint:
{Edge AI has lowest authority}
Regional Learning & Clustering

===== PAGE 43 =====
ClusterEngine
ClusterEngine
-------------------------
group_instances()
compute_cluster_stats()
validate_sample_size()
update_cluster_parameters()
Dependency:
ClusterEngine ..> YieldRecord
ClusterEngine ..> CropInstance
Constraint:
{Cluster updates apply prospectively only}
{Minimum sample threshold required}
Yield Verification Layer
YieldVerifier
YieldVerifier
-------------------------
compute_verification_score()
detect_outlier()
apply_biological_cap()
Dependency:
YieldVerifier ..> YieldRecord
YieldVerifier ..> RegionalCluster
Critical constraint:
{Farmer Truth preserved}
{ML Truth used for training only}
Confidence Propagation Layer

===== PAGE 44 =====
ConfidencePropagator
ConfidencePropagator
-------------------------
apply_confidence_weight()
reduce_low_confidence_influence()
adjust_recommendation_tone()
Dependency:
ConfidencePropagator ..> RiskPredictor
ConfidencePropagator ..> StressCNN
ConfidencePropagator ..> EdgeAIModel
Constraint:
{Low confidence softens influence}
Model Governance & Training
ModelRegistry
ModelRegistry
-------------------------
register_model()
activate_model()
rollback_model()
track_version_history()
Association:
Admin 1 ──── 0..* ModelRegistry
ModelRegistry 1 ──── 0..* MLModel
MLTrainingAudit
MLTrainingAudit
-------------------------
training_id : UUID
dataset_reference : String
feature_schema_version : String
evaluation_metrics : JSON

===== PAGE 45 =====
training_start : Timestamp
training_end : Timestamp
Association:
MLModel 1 ──── 0..* MLTrainingAudit
Constraint:
{Every model activation must have audit entry}
🏗 Authority Hierarchy (Critical Structural
Constraint)
We explicitly model:
Authority Order:
1. FarmerConfirmedAction
2. CTIS Rule Engine
3. Backend ML (RiskPredictor / StressCNN)
4. Weather Risk Model
5. Behavioral Adaptation
6. Edge AI
Attached as constraint note across ML classes.
Integration With CTIS
Dependency relationships:
ReplayEngine ..> RiskPredictor
ReplayEngine ..> StressCNN
ReplayEngine ..> ConfidencePropagator
ReplayEngine ..> YieldVerifier
Constraint:
{ML must NOT mutate stage directly}
{ML influences risk_score only}

===== PAGE 46 =====
SECTION 4.1 — Sequence Diagram
Log Action → Event → Replay → Snapshot
This is the core deterministic mutation pipeline.
Trigger
Actor: Farmer
Action: Logs an irrigation / fertilizer / structural update.
Participants (Lifelines)
1. Farmer (UI)
2. Backend API (FastAPI Controller)
3. Database (PostgreSQL)
4. DBEventDispatcher
5. EventHandler
6. ReplayEngine
7. StressEngine
8. RiskCalculator
9. SnapshotManager
Sequence Flow (Industrial Detail)
Phase 1 — Action Logging
Farmer → BackendAPI : POST /log-action
BackendAPI → BackendAPI : Validate JWT
BackendAPI → Database : INSERT ActionLog
BackendAPI → Database : INSERT Event (status=Created)
BackendAPI → Database : COMMIT TRANSACTION
BackendAPI → Farmer : 200 OK

===== PAGE 47 =====
Constraint:
{ActionLog + Event must persist atomically}
Phase 2 — Event Processing (Async Background Loop)
DBEventDispatcher → Database :
SELECT event WHERE status='Created'
ORDER BY priority, created_at
FOR UPDATE SKIP LOCKED
DBEventDispatcher → Database :
UPDATE event SET status='Processing'
Partition constraint:
{Partition key = crop_instance_id}
{FIFO within partition}
Phase 3 — Trigger Replay
EventHandler → ReplayEngine : replay(crop_instance_id)
ReplayEngine → Database :
SELECT crop_instance FOR UPDATE
Row-level lock acquired.
Constraint:
{Only one replay per crop_instance at a time}
Phase 4 — Snapshot Check
ReplayEngine → SnapshotManager :
load_latest_snapshot()
IF snapshot exists:
load state from snapshot
ELSE:
initialize baseline

===== PAGE 48 =====
Phase 5 — Ordered Action Application
ReplayEngine → Database :
SELECT ActionLogs
ORDER BY action_effective_date,
local_monotonic_counter
For each action:
ReplayEngine → ReplayEngine :
validate_invariants(action)
ReplayEngine → ReplayEngine :
apply_action()
ReplayEngine → StressEngine :
integrate_signals()
ReplayEngine → RiskCalculator :
compute_risk_index()
Constraints enforced:
{No stage jump}
{Drift bounded}
{Stress ∈ [0,1]}
Phase 6 — Snapshot Creation
If action_count % 50 == 0 OR structural action:
ReplayEngine → SnapshotManager :
create_snapshot()
Phase 7 — State Update
ReplayEngine → Database :
UPDATE crop_instance
SET stage,
stress_score,
risk_index,
updated_at
Commit transaction.
Row lock released.

===== PAGE 49 =====
Phase 8 — Event Finalization
EventHandler → Database :
UPDATE event SET status='Processed'
Failure Branch (Industrial Detail)
If replay fails:
ReplayEngine → Database : ROLLBACK
EventHandler → Database :
UPDATE event SET status='Failed'
retry_count += 1
If retry_count ≥ threshold:
EventHandler → Database :
UPDATE event SET status='DeadLetter'
EventHandler → Database :
UPDATE crop_instance SET state='RecoveryRequired'
Constraint:
{DeadLetter does not block other partitions}
Architectural Guarantees Modeled
This sequence diagram explicitly models:
✔ Atomic state mutation
✔ Partition isolation
✔ Idempotent event processing
✔ Snapshot optimization
✔ Deterministic replay
✔ Drift enforcement
✔ Dead-letter recovery
✔ FIFO order
✔ Row-level locking
This is not textbook-level — this is production-grade modeling.

===== PAGE 50 =====
Critical Invariant Notes Attached
{No state mutation without event}
{Replay must complete fully or rollback}
{Cross-crop events never interleave}
{Action order = (effective_date, monotonic_counter)}
SECTION 4.2 — Sequence Diagram
Media Upload → ML → Stress Update → Replay
Trigger
Actor: Farmer
Action: Upload crop image or video.
Participants (Lifelines)
1. Farmer (Browser / App)
2. Backend API
3. Cloud Storage
4. Media Worker (Cloud Run Worker)
5. StressCNN (Backend ML)
6. ConfidencePropagator
7. Database
8. DBEventDispatcher
9. EventHandler
10. ReplayEngine
11. StressEngine
12. RiskCalculator

===== PAGE 51 =====
Sequence Flow (Industrial Depth)
Phase 1 — Upload Initialization
Farmer → BackendAPI : POST /request-upload
BackendAPI → CloudStorage : Generate Signed URL
BackendAPI → Farmer : Return Signed URL
Security constraint:
{Signed URL expires in 15 minutes}
Phase 2 — Direct Media Upload
Farmer → CloudStorage : Upload media file
CloudStorage → BackendAPI : Notify upload complete (optional webhook)
No file stored locally on API server.
Phase 3 — Emit Media Processing Event
BackendAPI → Database :
INSERT Event(type="MediaUploaded", status="Created")
COMMIT
No crop mutation here.
Phase 4 — Event Dispatcher Picks Up
DBEventDispatcher → Database :
SELECT MediaUploaded WHERE status='Created'
FOR UPDATE SKIP LOCKED
DBEventDispatcher → Database :
UPDATE event SET status='Processing'

===== PAGE 52 =====
Phase 5 — Media Worker Invocation (Async)
EventHandler → MediaWorker :
process_media(media_id)
Isolation constraint:
{MediaWorker must not mutate CropInstance directly}
Phase 6 — ML Inference Pipeline
Inside MediaWorker:
MediaWorker → StressCNN :
analyze_image() OR analyze_video_frames()
StressCNN → MediaWorker :
stress_probability,
confidence_score
If multiple frames:
MediaWorker → MediaWorker :
aggregate_stress = weighted_mean(stress_i, confidence_i)
Phase 7 — Confidence Propagation
MediaWorker → ConfidencePropagator :
apply_confidence_weight()
final_stress_signal =
stress_probability * confidence
Constraint:
{Low confidence reduces influence}
Phase 8 — Emit MediaAnalyzed Event
MediaWorker → Database :
INSERT Event(type="MediaAnalyzed",

===== PAGE 53 =====
payload={stress_signal})
COMMIT
MediaWorker ends here.
No crop mutation performed.
Phase 9 — Event Dispatcher Triggers Replay
DBEventDispatcher → Database :
SELECT MediaAnalyzed WHERE status='Created'
FOR UPDATE SKIP LOCKED
UPDATE event SET status='Processing'
Phase 10 — Replay Execution
EventHandler → ReplayEngine :
replay(crop_instance_id)
Replay flow (similar to Section 4.1):
ReplayEngine → Database :
SELECT crop_instance FOR UPDATE
ReplayEngine :
load snapshot
load ordered actions
integrate stress_signal
compute new_stress
compute risk_index
enforce drift constraints
Constraint:
{ML influences stress_score only}
{ML must NOT directly change stage}

===== PAGE 54 =====
Phase 11 — Persist Updated State
ReplayEngine → Database :
UPDATE crop_instance
SET stress_score,
risk_index,
current_stage (if affected)
COMMIT
Phase 12 — Event Finalization
EventHandler → Database :
UPDATE event SET status='Processed'
Failure Branch
If MediaWorker fails:
EventHandler :
retry_count += 1
IF retry_count >= threshold:
move to DeadLetter
notify Admin
Replay not triggered if inference fails.
Architectural Guarantees Modeled
This sequence ensures:
✔ API never processes heavy ML
✔ Worker never mutates state directly
✔ ML output always event-driven
✔ Confidence reduces influence
✔ Replay enforces atomic update
✔ Stress bounded in integration layer
✔ Partition locking preserved

===== PAGE 55 =====
✔ Dead-letter path modeled
✔ Signed URL security
Explicit Safety Constraints
{No direct DB mutation from MediaWorker}
{All stress updates must pass through ReplayEngine}
{Replay must be atomic}
{ML has no stage authority}
SECTION 4.3 — Sequence Diagram
Yield Submission → Verification → Cluster
Update → Replay
Trigger
Actor: Farmer
Action: Submits crop yield after harvest.
Participants (Lifelines)
1. Farmer
2. Backend API
3. Database
4. YieldVerifier
5. ClusterEngine
6. DBEventDispatcher
7. EventHandler
8. ReplayEngine
9. SnapshotManager

===== PAGE 56 =====
Sequence Flow (Industrial Depth)
Phase 1 — Yield Submission
Farmer → BackendAPI : POST /submit-yield
BackendAPI → BackendAPI : Validate JWT
Phase 2 — Persist Yield Record (Farmer
Truth)
BackendAPI → Database :
INSERT YieldRecord(
reported_yield = input_value,
ml_yield_value = NULL,
yield_verification_score = NULL
)
COMMIT
Constraint:
{reported_yield must NEVER be mutated}
Phase 3 — Emit YieldSubmitted Event
BackendAPI → Database :
INSERT Event(type="YieldSubmitted", status="Created")
COMMIT
Phase 4 — Dispatcher Picks Up Event
DBEventDispatcher → Database :
SELECT YieldSubmitted WHERE status='Created'
FOR UPDATE SKIP LOCKED
UPDATE event SET status='Processing'

===== PAGE 57 =====
Partition isolation preserved.
Phase 5 — Yield Verification
EventHandler → YieldVerifier :
compute_verification_score(reported_yield)
Inside YieldVerifier:
YieldVerifier :
calculate_regional_mean()
calculate_regional_std()
z_score = (reported - mean) / std
If |z_score| > threshold:
YieldVerifier :
mark_as_outlier()
Phase 6 — Biological Cap Enforcement (ML
Truth Creation)
IF reported_yield > biological_limit:
ml_yield_value = biological_limit
ELSE:
ml_yield_value = reported_yield
Constraint:
{FarmerTruth ≠ MLTruth}
{MLTruth used for training only}
Phase 7 — Update Yield Record
YieldVerifier → Database :
UPDATE YieldRecord
SET ml_yield_value,
yield_verification_score
COMMIT

===== PAGE 58 =====
Phase 8 — Regional Cluster Update
(Prospective Only)
EventHandler → ClusterEngine :
validate_sample_size()
IF sample_size ≥ threshold:
update_cluster_parameters()
ELSE:
skip_cluster_update()
Constraint:
{Cluster updates apply to future CropInstances only}
{No retroactive mutation}
ClusterEngine does NOT trigger replay on past instances.
Phase 9 — Trigger Replay
EventHandler → ReplayEngine :
replay(crop_instance_id)
Replay includes:
ReplayEngine :
load snapshot
integrate final stress
compute final risk_index
finalize stage = Harvested
Phase 10 — Persist Final State
ReplayEngine → Database :
UPDATE crop_instance
SET state = 'Harvested',
stress_score,
risk_index
COMMIT

===== PAGE 59 =====
Phase 11 — Event Finalization
EventHandler → Database :
UPDATE event SET status='Processed'
Failure Branch
If YieldVerifier fails:
retry_count += 1
IF retry_count ≥ threshold:
move to DeadLetter
notify Admin
Replay not executed if verification incomplete.
Architectural Guarantees Modeled
This sequence ensures:
✔ Farmer Truth preserved
✔ ML Truth capped and separated
✔ Z-score anomaly detection modeled
✔ Cluster update threshold enforced
✔ No retroactive cluster mutation
✔ Replay atomic
✔ Partition isolation preserved
✔ Dead-letter path modeled
Explicit Safety Constraints
{Yield verification must precede cluster update}
{Cluster update must not trigger retroactive replay}
{MLTruth must not override FarmerTruth}
{Replay must be atomic}

===== PAGE 60 =====
SECTION 4.4 — Sequence Diagram
Service Request → Ranking → Completion →
Trust Recalculation
Trigger
Actor: Farmer
Action: Requests a service (e.g., irrigation equipment).
Participants (Lifelines)
1. Farmer
2. Backend API
3. Database
4. ExposureFairnessEngine
5. TrustScoreEngine
6. ServiceProvider
7. EscalationPolicy
8. FraudDetector
9. Admin (conditional branch)
Sequence Flow (Industrial Detail)
Phase 1 — Service Request Creation
Farmer → BackendAPI : POST /request-service
BackendAPI → Database :
INSERT ServiceRequest(status="Pending")
COMMIT
No CTIS mutation.

===== PAGE 61 =====
Phase 2 — Fetch Eligible Providers
BackendAPI → Database :
SELECT ServiceProviders
WHERE equipment_available
AND is_suspended = FALSE
Phase 3 — Trust Score Retrieval
BackendAPI → TrustScoreEngine :
get_trust_score(provider)
TrustScoreEngine may apply:
apply_trust_decay()
Constraint:
{TrustScore ∈ [0,1]}
Phase 4 — Exposure Fairness Adjustment
BackendAPI → ExposureFairnessEngine :
calculate_exposure_penalty(provider)
Ranking Score Computation:
ranking_score = trust_score - exposure_penalty
Constraint:
{TrustScore and ExposurePenalty computed independently}
Phase 5 — Return Ranked Providers
BackendAPI → Farmer :
Return sorted provider list

===== PAGE 62 =====
Phase 6 — Provider Accepts Request
ServiceProvider → BackendAPI : ACCEPT request
BackendAPI → Database :
UPDATE ServiceRequest(status="Accepted")
COMMIT
Phase 7 — Completion
ServiceProvider → BackendAPI : MARK complete
BackendAPI → Database :
UPDATE ServiceRequest(status="Completed")
COMMIT
Phase 8 — Review Submission
Farmer → BackendAPI : POST /submit-review
BackendAPI → Database :
INSERT ServiceReview
COMMIT
Constraint:
{Only one review per request}
Phase 9 — Trust Recalculation
BackendAPI → TrustScoreEngine :
recompute(provider)
TrustScoreEngine performs:
compute_completion_rate()
compute_complaint_ratio()
compute_rating_stability()
apply_verified_bonus()
apply_escalation_penalty()
apply_trust_decay()
clamp_score()
Then:

===== PAGE 63 =====
TrustScoreEngine → Database :
UPDATE ServiceProvider(trust_score)
COMMIT
Conditional Branch — Escalation
If severity_level > threshold:
BackendAPI → EscalationPolicy :
evaluate(review)
EscalationPolicy may:
flag_provider()
notify_admin()
reduce_exposure()
🕵 Conditional Branch — Fraud Detection
BackendAPI → FraudDetector :
analyze_review_pattern()
If suspicious:
FraudDetector → Database :
flag provider
Admin notified.
🏛 Admin Intervention (Optional)
Admin → BackendAPI : suspend provider
BackendAPI → Database :
UPDATE ServiceProvider(is_suspended=TRUE)
Constraint:
{All admin overrides logged}

===== PAGE 64 =====
CTIS Isolation Guarantee
Explicit constraint note attached:
{SOE flow must NOT trigger Replay}
{SOE must NOT modify CropInstance}
{SOE must NOT modify stress_score or risk_index}
SECTION 4.5 — Sequence Diagram
Admin Rule Activation → Simulation →
Prospective Enforcement
Trigger
Actor: Admin
Action: Activates a new Crop Rule Template.
Participants (Lifelines)
1. Admin
2. Backend API
3. Database
4. RuleEngineValidator
5. ReplayEngine (Simulation Mode)
6. DBEventDispatcher (optional manual replay)
7. EventHandler
8. CropInstance (future instances only)
Sequence Flow (Industrial Depth)

===== PAGE 65 =====
Phase 1 — Rule Draft Creation
Admin → BackendAPI : POST /create-rule-template
BackendAPI → Database :
INSERT CropRuleTemplate(status="Draft")
COMMIT
No crop affected.
Phase 2 — Validation & Simulation
Admin → BackendAPI : POST /validate-rule
BackendAPI → RuleEngineValidator :
validate(rule_template)
Validator checks:
check_stage_continuity()
check_drift_bounds()
check_duration_consistency()
check_loop_risk()
If validation passes:
Phase 3 — Simulation Replay (Isolated)
BackendAPI → ReplayEngine :
simulate_with_rule(sample_crop_instances)
Simulation Mode:
ReplayEngine :
clone crop_instance state
apply rule template
run deterministic replay
generate projected output
Constraint:
{Simulation must NOT persist state}
Results returned to Admin.

===== PAGE 66 =====
Phase 4 — Rule Activation
Admin → BackendAPI : POST /activate-rule
BackendAPI → Database :
UPDATE CropRuleTemplate(status="Active")
COMMIT
Constraint:
{Only one active rule per crop_type}
Phase 5 — Prospective Enforcement
From this point:
New CropInstance → assign rule_version_applied = new_rule_version
Existing CropInstances remain unchanged.
Constraint:
{Rule changes are prospective only}
Optional Branch — Manual Replay
If Admin chooses to reprocess specific crops:
Admin → BackendAPI : POST /manual-replay
BackendAPI → Database :
INSERT Event(type="ManualReplay")
COMMIT
Dispatcher picks up:
DBEventDispatcher → EventHandler :
handle(ManualReplay)
EventHandler → ReplayEngine :
replay(crop_instance_id)
Atomic as usual.

===== PAGE 67 =====
Failure Branch
If validation fails:
RuleEngineValidator → BackendAPI :
return error list
Activation blocked.
Critical Governance Constraints
{No automatic retroactive mutation}
{Manual replay required for old instances}
{Rule activation logged in AdminAuditLog}
{Rule template immutable after activation}
Architectural Guarantees Modeled
This sequence ensures:
✔ Rule validation before activation
✔ Simulation isolation
✔ Prospective enforcement
✔ No silent retroactive mutation
✔ Replay remains event-driven
✔ Admin override logged
✔ Deterministic guarantees preserved
SECTION 5 — State Machine Diagrams
SECTION 5.1 — CropInstance State Machine

===== PAGE 68 =====
States
We model the following states:
1. Initialized
2. Active
3. AtRisk
4. RecoveryRequired
5. Harvested
6. Closed
7. Archived
State Definitions
Initialized
• Created but no meaningful action logged yet.
• Rule version assigned.
• Baseline stage computed.
Transition:
on ActionLogged → Active
Active
Normal operating state.
Crop evolves via replay.
Transitions:
on StressThresholdExceeded → AtRisk
on YieldSubmitted → Harvested
on ReplayFailure → RecoveryRequired
on SeasonEnd → Closed

===== PAGE 69 =====
AtRisk
Stress score beyond safe bound.
Transitions:
on StressReduced → Active
on YieldSubmitted → Harvested
on ReplayFailure → RecoveryRequired
Constraint:
{AtRisk does NOT freeze stage}
RecoveryRequired
Entered when:
• DeadLetter threshold exceeded
• Replay invariant violation
• Critical system failure
Transitions:
on AdminManualReplaySuccess → Active
on AdminLockCrop → Archived
Constraint:
{No farmer actions allowed while RecoveryRequired}
Harvested
Yield submitted and replay finalized.
Transitions:
on ArchiveSeason → Archived
No transition back to Active.
Constraint:
{Harvested is terminal for stage evolution}

===== PAGE 70 =====
Closed
Season ended without yield.
Transitions:
on ArchiveSeason → Archived
Archived
Terminal state.
No outgoing transitions.
Constraint:
{Archived instances are immutable}
Transition Model (Textual UML
Representation)
[Initialized]
| ActionLogged
v
[Active]
| StressThresholdExceeded
v
[AtRisk]
| StressReduced
v
[Active]
[Active] --YieldSubmitted--> [Harvested]
[AtRisk] --YieldSubmitted--> [Harvested]
[Active] --ReplayFailure--> [RecoveryRequired]
[AtRisk] --ReplayFailure--> [RecoveryRequired]
[RecoveryRequired] --AdminManualReplaySuccess--> [Active]
[Harvested] --ArchiveSeason--> [Archived]
[Closed] --ArchiveSeason--> [Archived]

===== PAGE 71 =====
Forbidden Transitions (Explicitly Modeled)
Archived → Active
Harvested → Active
RecoveryRequired → Harvested (without replay)
AtRisk → Archived (without closure)
These will be shown as constraint notes in the diagram.
Event Triggers That Drive State
Event Resulting State Change
ActionLogged May remain Active
MediaAnalyzed May shift Active ↔ AtRisk
YieldSubmitted → Harvested
DeadLetter → RecoveryRequired
ManualReplaySuccess → Active
ArchiveSeason → Archived
Important:
All transitions occur inside ReplayEngine transaction boundary.
Critical Structural Constraints
Attached as notes:
{State transitions must occur inside replay transaction}
{No direct state mutation outside Event system}
{Harvested blocks further ActionLog insertion}
{RecoveryRequired blocks farmer mutation endpoints}
Why This State Machine Is Strong
Because it:
✔ Encodes deterministic lifecycle
✔ Separates operational vs failure state

===== PAGE 72 =====
✔ Models recovery intervention
✔ Prevents illegal resurrection
✔ Ensures immutability after archive
✔ Aligns with event-driven design
This is not a cosmetic diagram — this defines system safety.
SECTION 5.1 — CropInstance State Machine
States
We model the following states:
1. Initialized
2. Active
3. AtRisk
4. RecoveryRequired
5. Harvested
6. Closed
7. Archived
State Definitions
Initialized
• Created but no meaningful action logged yet.
• Rule version assigned.
• Baseline stage computed.
Transition:
on ActionLogged → Active
Active
Normal operating state.

===== PAGE 73 =====
Crop evolves via replay.
Transitions:
on StressThresholdExceeded → AtRisk
on YieldSubmitted → Harvested
on ReplayFailure → RecoveryRequired
on SeasonEnd → Closed
AtRisk
Stress score beyond safe bound.
Transitions:
on StressReduced → Active
on YieldSubmitted → Harvested
on ReplayFailure → RecoveryRequired
Constraint:
{AtRisk does NOT freeze stage}
RecoveryRequired
Entered when:
• DeadLetter threshold exceeded
• Replay invariant violation
• Critical system failure
Transitions:
on AdminManualReplaySuccess → Active
on AdminLockCrop → Archived
Constraint:
{No farmer actions allowed while RecoveryRequired}
Harvested
Yield submitted and replay finalized.
Transitions:

===== PAGE 74 =====
on ArchiveSeason → Archived
No transition back to Active.
Constraint:
{Harvested is terminal for stage evolution}
Closed
Season ended without yield.
Transitions:
on ArchiveSeason → Archived
Archived
Terminal state.
No outgoing transitions.
Constraint:
{Archived instances are immutable}
Transition Model (Textual UML
Representation)
[Initialized]
| ActionLogged
v
[Active]
| StressThresholdExceeded
v
[AtRisk]
| StressReduced
v
[Active]
[Active] --YieldSubmitted--> [Harvested]
[AtRisk] --YieldSubmitted--> [Harvested]

===== PAGE 75 =====
[Active] --ReplayFailure--> [RecoveryRequired]
[AtRisk] --ReplayFailure--> [RecoveryRequired]
[RecoveryRequired] --AdminManualReplaySuccess--> [Active]
[Harvested] --ArchiveSeason--> [Archived]
[Closed] --ArchiveSeason--> [Archived]
Forbidden Transitions (Explicitly Modeled)
Archived → Active
Harvested → Active
RecoveryRequired → Harvested (without replay)
AtRisk → Archived (without closure)
These will be shown as constraint notes in the diagram.
Event Triggers That Drive State
Event Resulting State Change
ActionLogged May remain Active
MediaAnalyzed May shift Active ↔ AtRisk
YieldSubmitted → Harvested
DeadLetter → RecoveryRequired
ManualReplaySuccess → Active
ArchiveSeason → Archived
Important:
All transitions occur inside ReplayEngine transaction boundary.
Critical Structural Constraints
Attached as notes:
{State transitions must occur inside replay transaction}
{No direct state mutation outside Event system}
{Harvested blocks further ActionLog insertion}
{RecoveryRequired blocks farmer mutation endpoints}

===== PAGE 76 =====
Why This State Machine Is Strong
Because it:
✔ Encodes deterministic lifecycle
✔ Separates operational vs failure state
✔ Models recovery intervention
✔ Prevents illegal resurrection
✔ Ensures immutability after archive
✔ Aligns with event-driven design
This is not a cosmetic diagram — this defines system safety.
SECTION 5.2 — Event Lifecycle State
Machine
Event States
We model the following states:
1. Created
2. Processing
3. Processed
4. Failed
5. DeadLetter
State Definitions
Created
Event has been persisted but not yet processed.
Entry Condition:

===== PAGE 77 =====
publish(event)
status = 'Created'
Transition:
on DispatcherSelect → Processing
Constraint:
{Event must be persisted before entering Created}
Processing
Dispatcher has selected event and acquired lock.
Entry Action:
SELECT FOR UPDATE SKIP LOCKED
status = 'Processing'
Transitions:
on HandlerSuccess → Processed
on HandlerException → Failed
on ServiceCrash → Created (with retry_count++)
Crash recovery rule:
{Any event stuck in Processing at restart resets to Created}
Processed
Event handled successfully.
Entry Action:
UPDATE status='Processed'
processed_at = now()
No outgoing transitions.
Constraint:
{Processed events are immutable}

===== PAGE 78 =====
Failed
Temporary failure occurred.
Entry Action:
retry_count += 1
Transitions:
if retry_count < threshold → Created
if retry_count ≥ threshold → DeadLetter
Constraint:
{Retry logic must be idempotent}
DeadLetter
Permanent failure state.
Entry Condition:
retry_count ≥ MAX_RETRY
Entry Actions:
status = 'DeadLetter'
notify_admin()
if crop-related:
set CropInstance.state = RecoveryRequired
Transition:
on AdminManualRetry → Created
Constraint:
{DeadLetter must not block other partitions}
State Transition Model (Textual UML
Representation)
[Created]
v
| DispatcherSelect

===== PAGE 79 =====
[Processing]
| Success
v
[Processed]
[Processing]
| Exception
v
[Failed]
[Failed]
v
| retry_count < threshold
[Created]
[Failed]
v
[DeadLetter]
| retry_count ≥ threshold
[DeadLetter]
| AdminManualRetry
v
[Created]
Partition Isolation Constraint
Attached note:
{Events processed FIFO per crop_instance_id}
{No cross-partition blocking}
{Row-level locking enforced}
Crash Recovery Behavior (Critical Detail)
On service restart:
if status == 'Processing':
status = 'Created'
retry_count += 1
Ensures:

===== PAGE 80 =====
✔ No event permanently stuck
✔ No partial state mutation
✔ Deterministic re-execution
Idempotency Enforcement
Attached constraint:
{event_hash UNIQUE}
{Duplicate publish rejected}
{Handler must be idempotent}
Why This Event State Machine Is Strong
Because it:
✔ Explicitly models retry loop
✔ Models crash recovery
✔ Models DeadLetter isolation
✔ Enforces idempotency structurally
✔ Prevents infinite retry loops
✔ Preserves partition safety
✔ Integrates with CropInstance RecoveryRequired state
This is production-grade reliability modeling.
SECTION 5.3 — ServiceRequest Lifecycle
State Machine
States
We model the following states:
1. Pending
2. Accepted
3. Completed

===== PAGE 81 =====
4. Cancelled
5. Escalated
6. Closed
State Definitions
Pending
Service request created but not yet accepted.
Entry Action:
status = 'Pending'
Transitions:
on ProviderAccept → Accepted
on FarmerCancel → Cancelled
on AutoTimeout → Cancelled
Constraint:
{No review allowed in Pending state}
Accepted
Provider has committed to service.
Entry Action:
lock provider slot
update schedule
Transitions:
on ProviderComplete → Completed
on FarmerCancelBeforeStart → Cancelled
on AdminSuspendProvider → Cancelled
Constraint:
{Provider must be verified and not suspended}

===== PAGE 82 =====
Completed
Service executed successfully.
Entry Action:
enable_review_submission = TRUE
Transitions:
on ReviewSubmitted → Closed
on ReviewSeverityHigh → Escalated
on NoReviewTimeout → Closed
Constraint:
{Only one review per request}
Cancelled
Request aborted before completion.
Entry Conditions:
• Farmer cancellation
• Provider unavailable
• Admin suspension
• Timeout
No outgoing transitions.
Constraint:
{Cancelled requests cannot be reopened}
Escalated
Triggered by:
• High severity complaint
• Fraud detection flag
• Repeated complaints
Entry Action:
notify_admin()
reduce_exposure()

===== PAGE 83 =====
Transitions:
on AdminResolve → Closed
on AdminSuspendProvider → Closed
Constraint:
{Escalation does not automatically suspend provider}
Closed
Terminal state.
Entry Condition:
• Review processed
• Escalation resolved
• No-review timeout
No outgoing transitions.
Constraint:
{Closed requests immutable}
State Transition Model (Textual UML
Representation)
[Pending]
| ProviderAccept
v
[Accepted]
[Pending] --FarmerCancel--> [Cancelled]
[Pending] --Timeout--> [Cancelled]
[Accepted]
| ProviderComplete
v
[Completed]
[Accepted] --AdminSuspend--> [Cancelled]
[Completed]
| ReviewSubmitted

===== PAGE 84 =====
v
[Closed]
[Completed] --HighSeverityComplaint--> [Escalated]
[Escalated]
| AdminResolve
v
[Closed]
Fraud Detection Integration
FraudDetector may trigger:
on FraudFlagDetected → Escalated
But does NOT directly suspend.
Trust Score Integration
Transitions triggering trust recalculation:
Transition TrustScoreEngine Trigger
Completed → Closed Yes
Escalated → Closed Yes
Cancelled No
Trust recalculation does not change state directly.
Critical Constraints
Attached as notes:
{ServiceRequest must not mutate CropInstance}
{Escalation requires admin review}
{Review allowed only after Completed}
{Cancelled and Closed are terminal}

===== PAGE 85 =====
Why This State Machine Is Strong
Because it:
✔ Separates operational vs failure states
✔ Prevents review misuse
✔ Encodes escalation path
✔ Encodes fraud trigger
✔ Encodes admin authority
✔ Prevents illegal reopening
✔ Preserves CTIS isolation
SECTION 6 — Component Diagram
CultivaX Logical Architecture
High-Level Logical Components
We model the following components:
Client Layer
• Web App (Next.js PWA)
API Layer
• Backend API (FastAPI)
• Authentication Module
• Event Dispatcher
• Replay Engine (CTIS)
• SOE Module
• ML Integration Module
Worker Layer
• Media Worker (Cloud Run)
• Optional Video Worker

===== PAGE 86 =====
Data Layer
• PostgreSQL (Cloud SQL)
• Cloud Storage
• Model Registry Storage
External Systems
• Weather API
• Edge AI Device
Component Breakdown (Block
Representation)
🖥 Web Application (Next.js)
Component: WebApp (PWA)
Responsibilities:
- User Interaction
- Dashboard Rendering
- Timeline Visualization
- Service Requests
- What-If Simulation UI
Dependency:
WebApp → BackendAPI (HTTPS REST)
No direct DB access.
⚙ Backend API (FastAPI)
Component: BackendAPI
Submodules:
- AuthController
- CropController
- ServiceController
- AdminController
- EventDispatcher

===== PAGE 87 =====
Dependencies:
BackendAPI → PostgreSQL
BackendAPI → EventDispatcher
BackendAPI → CloudStorage (signed URL generation)
BackendAPI → WeatherAPI
Constraint:
{No heavy ML processing inside API}
Event Dispatcher
Component: DBEventDispatcher
Dependencies:
DBEventDispatcher → PostgreSQL
DBEventDispatcher → ReplayEngine
DBEventDispatcher → SOEHandlers
Constraint:
{All state mutation flows through EventDispatcher}
CTIS Module (Replay Engine)
Component: CTIS
Submodules:
- ReplayEngine
- StressEngine
- RiskCalculator
- SnapshotManager
Dependencies:
CTIS → PostgreSQL
CTIS → MLModule
Constraint:
{No external mutation allowed}

===== PAGE 88 =====
SOE Module
Component: SOE
Submodules:
- TrustScoreEngine
- ExposureFairnessEngine
- EscalationPolicy
- FraudDetector
Dependencies:
SOE → PostgreSQL
Constraint:
{SOE cannot call ReplayEngine}
ML Module
Component: MLModule
Submodules:
- RiskPredictor
- StressCNN
- YieldVerifier
- ClusterEngine
- ConfidencePropagator
- ModelRegistry
Dependencies:
MLModule → PostgreSQL (model metadata)
MLModule → CloudStorage (model files)
Constraint:
{ML does not directly mutate CropInstance}
Media Worker
Component: MediaWorker
Responsibilities:
• Video frame extraction
• Image preprocessing
• CNN inference

===== PAGE 89 =====
• Emit MediaAnalyzed event
Dependencies:
MediaWorker → CloudStorage
MediaWorker → PostgreSQL (event insertion)
MediaWorker → MLModule (StressCNN)
Constraint:
{MediaWorker must not modify CropInstance}
🗄 PostgreSQL (Cloud SQL)
Single source of truth.
Stores:
• Users
• CropInstances
• ActionLogs
• EventLogs
• Snapshots
• ServiceRequests
• Reviews
• Trust data
• ML metadata
Constraint:
{All mutable state persisted here}
☁ Cloud Storage
Stores:
• Media files
• ML model artifacts
Accessed via:
• Signed URLs

===== PAGE 90 =====
🌦 Weather API (External)
Used by:
• RiskCalculator
No persistent mutation.
Component Interaction Model (Textual)
WebApp
↓ HTTPS
BackendAPI
↓
EventDispatcher
↓
ReplayEngine (CTIS)
↓
PostgreSQL
BackendAPI → SOE
BackendAPI → MLModule
MediaWorker → MLModule
MediaWorker → EventDispatcher
All persistent state → PostgreSQL
All media → CloudStorage
Architectural Integrity Constraints
{No shared in-memory state between services}
{Replay isolated per crop_instance}
{Workers containerized}
{ML influence bounded}
{SOE isolated from CTIS}
Why This Component Diagram Is Strong
Because it:

===== PAGE 91 =====
✔ Shows logical module boundaries
✔ Shows worker isolation
✔ Shows event-driven mutation
✔ Shows no shared memory
✔ Shows ML assistive integration
✔ Shows governance boundaries
✔ Shows external dependency isolation
This aligns exactly with your Cloud Run + PostgreSQL stack.
SECTION 7 — Deployment Diagram
CultivaX Cloud Deployment Architecture
Deployment Nodes
We model the following runtime nodes:
1. User Device (Browser / PWA)
2. Google Cloud Platform
a. Cloud Run (Backend API)
b. Cloud Run (Media Worker)
c. Cloud Run (Video Worker - optional)
d. Cloud SQL (PostgreSQL)
e. Cloud Storage
f. Cloud Scheduler
3. External Weather API
4. Edge AI Device (Optional)
Node Breakdown (Block Representation)

===== PAGE 92 =====
🖥 Node 1 — User Device
Node: User Device
Artifact: WebApp (Next.js PWA)
Protocol: HTTPS
Communicates with:
User Device → Backend API (HTTPS REST)
No direct DB or storage access.
☁ Node 2 — Cloud Run (Backend API)
Node: CloudRun_API
Container: FastAPI + Uvicorn
Autoscaling: Enabled
Responsibilities:
• Authentication
• Crop lifecycle endpoints
• SOE endpoints
• Admin endpoints
• Event dispatcher loop
• ML inference (lightweight only)
Connections:
CloudRun_API → CloudSQL (TCP)
CloudRun_API → CloudStorage (Signed URL)
CloudRun_API → WeatherAPI (HTTPS)
Constraint:
{No local file storage}
Node 3 — Cloud Run (Media Worker)
Node: CloudRun_MediaWorker
Container: Python Worker
Autoscaling: Enabled
Responsibilities:

===== PAGE 93 =====
• Media processing
• CNN inference
• Frame extraction
• Emit MediaAnalyzed events
Connections:
MediaWorker → CloudStorage (read media)
MediaWorker → CloudSQL (insert event)
Constraint:
{No direct CropInstance mutation}
Node 4 — Cloud Run (Video Worker - Optional)
Dedicated to heavy video processing.
Isolation prevents API overload.
🗄 Node 5 — Cloud SQL (PostgreSQL)
Node: CloudSQL_PostgreSQL
Stores:
• All domain data
• Events
• Snapshots
• SOE data
• ML metadata
Connections:
CloudRun_API ↔ CloudSQL
MediaWorker ↔ CloudSQL
Constraint:
{Single source of truth}
{ACID enforced}

===== PAGE 94 =====
☁ Node 6 — Cloud Storage
Node: CloudStorage
Stores:
• Images
• Videos
• Model artifacts
Access via:
Signed URL (temporary)
Node 7 — Cloud Scheduler
Node: CloudScheduler
Triggers:
• Media cleanup
• Retention jobs
• Model monitoring
• Dead-letter monitoring
• Periodic health checks
Communicates with:
CloudScheduler → BackendAPI
🌦 External Weather API
Node: WeatherAPI
Protocol: HTTPS
Used by:
BackendAPI → WeatherAPI
No persistent connection.
Edge AI Device (Optional)
Node: EdgeDevice

===== PAGE 95 =====
Runs:
• TensorFlow Lite model
Uploads:
EdgeDevice → BackendAPI (stress result)
Constraint:
{Edge AI lowest authority}
Deployment Interaction Model
UserDevice
↓ HTTPS
CloudRun_API
↓
CloudSQL_PostgreSQL
CloudRun_API → CloudStorage (signed URL)
CloudRun_API → WeatherAPI
MediaWorker → CloudStorage
MediaWorker → CloudSQL
CloudScheduler → CloudRun_API
Security Boundaries
We model:
Internet Boundary
-----------------
UserDevice ↔ CloudRun_API (HTTPS only)
Internal GCP Network
--------------------
CloudRun_API ↔ CloudSQL (private IP)
MediaWorker ↔ CloudSQL
No public DB exposure.

===== PAGE 96 =====
Scalability Model
API Scaling
CloudRun_API:
- Horizontal scaling
- Stateless containers
- Multiple instances
Partition locking ensures no duplicate replay.
Worker Scaling
MediaWorker:
- Independent autoscaling
- Does not affect API performance
Database Scaling
CloudSQL:
- Connection pooling
- Read replica possible
- Future partitioning possible
Critical Deployment Constraints
{No shared memory between services}
{All state persisted in DB}
{Event system guarantees idempotency across instances}
{Replay atomic per partition}
{MediaWorker isolated from CTIS mutation}
Why This Deployment Diagram Is Strong
Because it:

===== PAGE 97 =====
✔ Shows physical separation
✔ Shows autoscaling model
✔ Shows worker isolation
✔ Shows network boundaries
✔ Shows DB as single source of truth
✔ Shows external API boundary
✔ Shows future migration readiness
✔ Matches Cloud Run + Cloud SQL stack exactly
UML Structural Modeling Completed
You now have:
✔ Use Case Diagrams (layered)
✔ Class Diagrams (layered architecture)
✔ Sequence Diagrams (industrial runtime flows)
✔ State Machine Diagrams (behavioral integrity)
✔ Component Diagram (logical architecture)
✔ Deployment Diagram (physical architecture)
SECTION 8.1 — Activity Diagram
Deterministic Replay Engine Workflow
This models the complete replay algorithm, including:
• Lock acquisition
• Snapshot optimization
• Action ordering
• Stress integration
• Risk computation
• Drift enforcement
• Snapshot creation
• Failure rollback

===== PAGE 98 =====
Activity Nodes (Structured Flow)
Start
[Start]
↓
Receive ReplayTriggered Event
Acquire Lock
Acquire row lock on CropInstance
(SELECT FOR UPDATE)
Decision:
Lock Acquired?
Yes → Continue
No → Retry Later
Load Snapshot
Load latest snapshot (if exists)
Decision:
Snapshot Exists?
Yes → Initialize state from snapshot
No → Initialize baseline from sowing_date
Load Ordered Actions
Fetch ActionLogs
ORDER BY (effective_date, monotonic_counter)
Loop: Apply Actions
For each action:

===== PAGE 99 =====
Validate chronological invariant
↓
Apply action to state
↓
Update stage offset
Decision:
Invariant violated?
Yes → Abort → Mark CropInstance = RecoveryRequired
No → Continue
Stress Integration
Collect stress signals:
- Backend ML
- Weather Risk
- Deviation Penalty
- Edge Signal
Apply confidence weighting
Apply exponential smoothing
Bound stress ∈ [0,1]
Risk Computation
Compute risk_index
Apply stage sensitivity factor
Clamp to [0,1]
⚖ Drift Enforcement
Decision:
|stage_offset| > max_drift?
Yes → Clamp to boundary
No → Continue
Snapshot Decision
Decision:

===== PAGE 100 =====
ActionCount % 50 == 0 OR StructuralAction?
Yes → Create Snapshot
No → Skip
Persist Updated State
UPDATE CropInstance
COMMIT TRANSACTION
Release Lock
Failure Path
If exception occurs:
ROLLBACK TRANSACTION
Increment retry_count
Return control to EventDispatcher
End
Mark Event = Processed
[End]
Textual UML Flow Representation
[Start]
↓
[Acquire Lock]
↓
[Load Snapshot?]
↙ ↘
[Yes] [No]
↓ ↓
[Initialize State]
↓
[Load Actions]
↓
[Loop Apply Actions]
↓
[Integrate Stress]
↓

===== PAGE 101 =====
[Compute Risk]
↓
↓
↓
↓
↓
↓
[End]
[Enforce Drift]
[Create Snapshot?]
[Persist State]
[Release Lock]
[Mark Event Processed]
Failure branch from any processing node → Rollback → DeadLetter logic.
Critical Workflow Constraints
{Replay must be atomic}
{No state mutation outside replay}
{Partition isolation enforced}
{All mutations event-driven}
{Failure must rollback completely}
SECTION 8.2 — Activity Diagram
Stress Integration & Signal Arbitration Workflow
Purpose
Compute a safe, bounded, confidence-weighted stress score using:
• Backend ML (CNN)
• Weather risk
• Deviation penalty
• Edge AI

===== PAGE 102 =====
• Previous stress memory
Without violating:
• Stage drift constraints
• Daily jump limit
• Authority hierarchy
Activity Flow Structure
Start
[Start]
↓
ReplayEngine requests stress update
Collect Signals
Collect:
backend_ml_stress
weather_risk
deviation_penalty
edge_signal
previous_stress
Decision:
Any signal missing?
Yes → substitute fallback value
No → continue
⚖ Authority Weight Assignment
Assign weights:
w_backend > w_weather > w_deviation > w_edge
Constraint:
{Edge AI lowest authority}

===== PAGE 103 =====
Compute Weighted Signal
stress_signal =
w_backend * backend_ml +
w_weather * weather_risk +
w_deviation * deviation_penalty +
w_edge * edge_signal
Apply Confidence Propagation
effective_signal =
stress_signal * confidence_score
Decision:
confidence_score < threshold?
Yes → reduce influence further
No → continue
Apply Exponential Smoothing
new_stress =
alpha * effective_signal +
(1 - alpha) * previous_stress
This prevents oscillation.
Daily Jump Limit Enforcement
Decision:
|new_stress - previous_stress| > max_daily_delta?
Yes → clamp to boundary
No → continue
Stage-Sensitive Adjustment
Apply stage_sensitivity_factor
Example:

===== PAGE 104 =====
• Germination → high sensitivity
• Mature stage → lower sensitivity
Bound Stress
if new_stress < 0 → set 0
if new_stress > 1 → set 1
Constraint:
{stress ∈ [0,1]}
Return to ReplayEngine
Return new_stress
[End]
Textual UML Flow Representation
[Start]
↓
[Collect Signals]
↓
[Assign Authority Weights]
↓
[Compute Weighted Signal]
↓
[Apply Confidence]
↓
[Apply Smoothing]
↓
[Check Daily Jump]
↙ ↘
[Clamp] [Continue]
↓
[Stage Sensitivity Adjustment]
↓
[Bound to 0–1]
↓
[Return Stress]
↓
[End]

===== PAGE 105 =====
Explicit Safety Constraints
{ML cannot directly modify stage}
{Confidence must reduce influence}
{Daily delta must be bounded}
{Stress must remain probabilistic}
{Authority order strictly enforced}
Why This Diagram Is Architecturally Strong
Because it:
✔ Encodes multi-signal arbitration
✔ Enforces authority hierarchy
✔ Models smoothing memory
✔ Prevents oscillation
✔ Prevents shock spikes
✔ Encodes bounded probabilistic output
✔ Makes ML assistive, not dominant
SECTION 8.3 — Activity Diagram
Trust Score Recalculation Workflow
Trigger
Triggered by:
• ServiceRequest → Completed
• ServiceReview submitted
• Admin override
• Scheduled decay update

===== PAGE 106 =====
Activity Flow Structure
Start
[Start]
↓
Trigger TrustScoreEngine.recompute(provider_id)
Fetch Provider Data
Load:
completion_count
total_requests
complaint_count
review_ratings
last_completion_date
verification_status
escalation_flags
Decision:
Provider suspended?
Yes → trust_score = 0
No → Continue
Compute Completion Rate (CR)
CR = completion_count / total_requests
Edge case:
If total_requests == 0 → CR = default (0.5)
⚠ Compute Complaint Ratio (CPR)
CPR = complaint_count / total_requests
Bound CPR to [0,1].

===== PAGE 107 =====
Compute Average Rating (AR)
normalized_rating = average_rating / 5
If no reviews:
normalized_rating = neutral_value
Apply Escalation Penalty (EP)
Decision:
Active escalation flags?
Yes → EP = penalty_weight
No → EP = 0
✔ Apply Verified Bonus (VB)
If verification_status == TRUE:
VB = bonus_weight
Else:
VB = 0
Compute Base Trust Score
trust_base =
w1 * CR +
w2 * (1 - CPR) +
w3 * normalized_rating +
w4 * VB -
w5 * EP
Constraint:
w1 > w2 > w3
Apply Trust Decay
days_inactive = today - last_completion_date
decay_factor = exp(-λ * days_inactive)
trust_after_decay = trust_base * decay_factor

===== PAGE 108 =====
Purpose:
Prevents long-term dominance without activity.
🕵 Fraud Adjustment (Optional)
Decision:
Fraud flag detected?
Yes → reduce trust_after_decay
No → Continue
Important:
{Fraud flag does NOT auto-suspend}
Clamp Final Score
if trust_after_decay < 0 → 0
if trust_after_decay > 1 → 1
Persist Updated Trust
UPDATE ServiceProvider
SET trust_score = final_value
COMMIT
End
Return final trust_score
[End]
Textual UML Flow Representation
[Start]
↓
[Load Provider Data]
↓
[Suspended?]
↙ ↘

===== PAGE 109 =====
[Yes=0] [No]
↓
[Compute CR]
↓
[Compute CPR]
↓
[Compute Rating]
↓
[Apply Escalation Penalty]
↓
[Apply Verified Bonus]
↓
[Compute Base Score]
↓
[Apply Trust Decay]
↓
[Fraud Flag?]
↙ ↘
[Reduce] [Continue]
↓
[Clamp 0–1]
↓
[Persist]
↓
[End]
Explicit Architectural Constraints
{Trust score independent of exposure ranking}
{ExposurePenalty applied outside this workflow}
{Trust must remain probabilistic [0,1]}
{Decay must apply continuously over inactivity}
{Fraud flag influences score but not state directly}
Why This Diagram Is Strong
Because it:
✔ Encodes mathematical trust foundation
✔ Separates trust from exposure fairness
✔ Applies decay over time
✔ Encodes fraud influence properly

===== PAGE 110 =====
✔ Prevents monopoly accumulation
✔ Maintains probabilistic bounds
✔ Ensures governance auditability
SECTION 8.4 — Activity Diagram
ML Training & Governance Lifecycle
Trigger
Triggered by:
• Admin initiates training
• Scheduled retraining job
• Performance drift detection
Activity Flow Structure
Start
[Start]
↓
Admin / Scheduler triggers model training
Dataset Preparation
Collect:
historical CropInstances
YieldRecords (MLTruth only)
stress history
weather data
deviation profiles
Constraint:

===== PAGE 111 =====
{No PII included}
Feature Engineering
Normalize numeric features
Encode categorical features
Standardize distributions
Assign Feature Schema Version
feature_schema_version = "vX.Y"
Constraint:
{Feature schema must be stored in MLTrainingAudit}
Model Training
Train:
OR
LogisticRegression (RiskPredictor)
CNN (StressCNN)
Record:
training_start
training_end
Evaluation Phase
Compute:
Accuracy
Precision
Recall
F1 Score
AUC
Decision:
Performance >= threshold?
Yes → Continue
No → Reject Model

===== PAGE 112 =====
If rejected:
Log failure in MLTrainingAudit
[End]
🗂 Store Model Artifact
Upload model file to CloudStorage
Register model metadata in ModelRegistry
Create MLTrainingAudit Entry
INSERT MLTrainingAudit:
dataset_reference
feature_schema_version
evaluation_metrics
training timestamps
Constraint:
{Every model must have audit record}
Drift Check (Optional)
If retraining triggered by drift:
Compare new metrics vs active model metrics
Decision:
Significant improvement?
Yes → Continue
No → Keep old model
Admin Review Gate
Admin reviews evaluation metrics
Decision:
Approve Activation?
Yes → Activate Model
No → Reject

===== PAGE 113 =====
Model Activation
Set new model.is_active = TRUE
Set old model.is_active = FALSE
Constraint:
{Activation must not retroactively mutate CropInstances}
Rollback Capability
If new model later underperforms:
Admin triggers rollback
Restore previous active model
Log rollback in audit
End
[End]
Textual UML Flow Representation
[Start]
↓
[Prepare Dataset]
↓
[Feature Engineering]
↓
[Assign Feature Schema]
↓
[Train Model]
↓
[Evaluate Performance]
↓
[Meets Threshold?]
↙ ↘
[Reject] [Store Artifact]
↓
[Create Audit Entry]

===== PAGE 114 =====
↓
[Admin Review]
↓
[Activate Model]
↓
[End]
Critical Governance Constraints
{No model activation without audit}
{Model version must be traceable}
{Feature schema must match inference schema}
{No retroactive mutation on activation}
{Rollback must be possible}
Why This ML Lifecycle Diagram Is Strong
Because it:
✔ Encodes reproducibility
✔ Separates training from activation
✔ Includes performance gating
✔ Includes admin review gate
✔ Includes audit logging
✔ Includes drift detection
✔ Includes rollback path
✔ Prevents silent ML upgrades
SECTION 8.5 — Activity Diagram
Offline Sync & Bulk Event Reconciliation
Workflow

===== PAGE 115 =====
Scenario
Farmer performs actions offline:
• Logs irrigation
• Logs fertilizer
• Uploads delayed media
When connection restored → Bulk sync occurs.
Activity Flow Structure
Start
[Start]
↓
Device regains network connectivity
Collect Unsynced Actions
Retrieve local unsynced actions
Each action includes:
action_effective_date
local_monotonic_counter
payload
Constraint:
{local_monotonic_counter strictly increasing per device}
Send Bulk Sync Request
Device → BackendAPI : POST /offline-sync
{array of unsynced actions}

===== PAGE 116 =====
Validate Sync Payload
BackendAPI:
Validate JWT
Validate payload structure
Decision:
Duplicate event_hash?
Yes → Skip insertion
No → Continue
Idempotency preserved.
Insert ActionLogs & Events
For each valid action:
INSERT ActionLog
INSERT Event(type="ActionLogged")
All inside one transaction per action.
Constraint:
{ActionLog + Event must be atomic}
Order Verification
Ensure sorting by:
(action_effective_date,
local_monotonic_counter)
Decision:
Out-of-order detected?
Yes → Reject payload
No → Continue
Chronological integrity enforced.

===== PAGE 117 =====
Debounce Replay
Multiple ActionLogged events may exist.
Instead of triggering replay per action:
Check:
Has ReplayTriggered been inserted recently?
Decision:
Replay already queued within debounce window?
Yes → Do not insert new ReplayTriggered
No → Insert ReplayTriggered
Constraint:
{Prevent replay storm}
Dispatcher Processing
Replay flow executes as defined in Section 8.1.
Partition locking ensures:
{Only one replay per crop_instance at a time}
Sync Completion Response
BackendAPI → Device :
Sync success
Return updated crop state
Failure Branch
If invariant violation detected during replay:
Mark CropInstance = RecoveryRequired
Notify Admin
Return error to device

===== PAGE 118 =====
Textual UML Flow Representation
[Start]
↓
[Collect Unsynced Actions]
↓
[Send Bulk Sync]
↓
[Validate Payload]
↓
[Duplicate?]
↙ ↘
[Skip] [Insert ActionLog + Event]
↓
[Order Verification]
↓
[Out-of-order?]
↙ ↘
[Reject] [Continue]
↓
[Debounce Replay]
↓
[Dispatcher Handles Replay]
↓
[Return Updated State]
↓
[End]
Critical Safety Constraints
{Idempotency via event_hash}
{Chronological ordering strictly enforced}
{local_monotonic_counter prevents reordering}
{Replay must remain atomic}
{Bulk sync must not cause replay storm}
{Offline actions cannot bypass invariants}
Why This Offline Workflow Is Strong
Because it:

===== PAGE 119 =====
✔ Preserves deterministic replay
✔ Prevents duplicate state mutation
✔ Enforces strict chronological ordering
✔ Prevents replay storm
✔ Maintains partition safety
✔ Integrates seamlessly with Event system
✔ Solves real-time vs offline contradiction
SECTION 9 — Requirements Traceability
Matrix (RTM)
Purpose of RTM
The RTM establishes bidirectional traceability between:
1. SRS Requirements
2. UML Models
3. TDD Implementation
4. Database Schema / Modules
It ensures:
✔ No requirement is ignored
✔ No feature is undocumented
✔ No design exists without requirement
✔ No implementation exists without design
Traceability Structure
We map:
SRS ID
→ Requirement Description
→ UML Artifact(s)
→ TDD Section
→ Implementation Component

===== PAGE 120 =====
RTM — Functional Requirements
Crop Lifecycle Management
SRS
TDD
ID Requirement UML Reference
Reference Implementation
FR-
1 User shall create CropInstance Use Case: Farmer
DB Schema
Subsystem
2.3.1 crop_instances table
FR-
User shall log actions
2
chronologically Sequence 4.1 CTIS Replay
Section action_logs table
FR-
System shall recompute state
3
deterministically Activity 8.1 Replay
Algorithm ReplayEngine
FR-
4 System shall enforce drift bounds Activity 8.1 Drift
Enforcement ReplayEngine
FR-
5 System shall compute stress score Activity 8.2 Stress Formula StressEngine
FR-
6 System shall compute risk index Activity 8.2 Risk
Calculation RiskCalculator
FR-
System shall support What-If
7
simulation Use Case: Farmer TDD What-If ReplayEngine
(clone mode)
FR-
8 System shall support offline sync Activity 8.5 Event
Dispatcher event_log table
Media & ML
SRS
ID Requirement UML Reference TDD Reference Implementation
FR-
9 User shall upload media Sequence 4.2 Media Section media_files table
FR-
System shall analyze stress via
10
ML Activity 8.2 ML Integration StressCNN
FR-
ML must provide confidence
11
score Activity 8.2 Confidence
Propagation ConfidencePropagator
FR-
ML shall not directly mutate
Sequence 4.2
Authority
12
stage
Constraint
Hierarchy ReplayEngine
FR-
13 Yield must be verified Sequence 4.3 Yield Verification YieldVerifier
FR-
Regional cluster must update
14
prospectively Sequence 4.3 ClusterEngine ml_models + cluster
logic

===== PAGE 121 =====
Service Orchestration (SOE)
SRS
ID Requirement UML
Reference TDD Reference Implementation
FR-
15 Farmer shall request service Sequence 4.4 SOE Section service_requests
table
FR-
16 System shall rank providers fairly Activity 8.3 Trust &
Exposure TrustScoreEngine
FR-
17 Trust shall decay over inactivity Activity 8.3 Trust Decay TrustScoreEngine
FR-
18 System shall detect fraud Sequence 4.4 Fraud Logic FraudDetector
FR-
High severity complaint shall
19
escalate State 5.3 EscalationPolicy service_reviews
table
FR-
20 SOE must not affect CTIS state Constraint in
Isolation
Architectural
3.4
Guarantee
constraint
Event System
SRS
ID Requirement UML Reference TDD
Reference Implementation
FR-
All state mutation must be event-
21
driven Event Class 3.3 Event Section event_log table
FR-
22 Events must be idempotent Event State
Machine 5.2 Hash logic UNIQUE(event_hash)
FR-
23 System shall retry failed events Event State
Machine Retry Logic retry_count field
FR-
24 System shall isolate partitions Sequence 4.1 Locking
SELECT FOR
Strategy
UPDATE
FR-
System shall support dead-letter
25
handling
Event State
Machine
DeadLetter
Logic status='DeadLetter'
🏛 Governance & ML Lifecycle
SRS
ID Requirement UML Reference TDD Reference Implementation
FR-
Admin shall activate rule
26
templates Sequence 4.5 Rule Governance rule_version_applied
FR-
Rule changes must be
27
prospective State 5.1 Rule Activation CropRuleTemplate
FR-
28 Admin shall activate ML model Activity 8.4 ModelRegistry ml_models table
FR-
29 ML must be auditable Activity 8.4 MLTrainingAudit ml_training_audit
table

===== PAGE 122 =====
FR-
30
Admin shall retry dead-letter
events
Event State
Machine ManualRetry event_log table
Non-Functional Requirements Traceability
Performance
NFR ID Requirement UML Evidence Implementation
NFR-1 Replay must be efficient Snapshot in 8.1 crop_instance_snapshots
NFR-2 System must scale to 600 users Deployment Diagram Cloud Run autoscaling
NFR-3 Media processing must not block API Sequence 4.2 MediaWorker isolation
Security
NFR ID Requirement UML Evidence Implementation
NFR-4 JWT Authentication Component Diagram AuthController
NFR-5 DB ACID consistency Sequence 4.1 PostgreSQL
NFR-6 No direct state mutation Constraints across diagrams Event-driven model
NFR-7 Idempotent event processing Event Lifecycle event_hash UNIQUE
Reliability
NFR ID Requirement UML Evidence Implementation
NFR-8 Crash recovery support Event Lifecycle Reset Processing → Created
NFR-9 Dead-letter isolation Event State Machine status='DeadLetter'
NFR-10 Partition isolation Sequence 4.1 crop_instance_id partition
Traceability Coverage Summary
Total Functional Requirements: 30
Total Non-Functional Requirements: 10
Coverage:
✔ 100% traced to UML
✔ 100% traced to TDD
✔ 100% traced to implementation artifacts

===== PAGE 123 =====
No orphan requirements.
No undocumented design elements.
\end{verbatim}
\section{Tech Stack}
\begin{verbatim}
# FILE: CultivaX Tech stack.pdf
# PAGES: 6



===== PAGE 1 =====
Teck Stack for CultivaX
STRATEGIC STACK DECISION (Finalized
for CultivaX Phase 1)
I will now lock a stack that is:
• Conservative
• Clean
• Scalable to 500–600 users without rewrite
• Migration-friendly to distributed architecture
• Research-ready later
1. BACKEND (Conservative but Powerful)
Python 3.11 + FastAPI
Why this is safe:
• Mature
• Async support
• High performance for moderate load
• Native ML integration
• Easy worker separation
• Clean OpenAPI auto docs
Scaling to 500–600 users?
Yes, easily.
Cloud Run can scale horizontally.
FastAPI async handles moderate concurrency well.
No rewrite required.

===== PAGE 2 =====
2. DATABASE
PostgreSQL (Cloud SQL)
Why:
• ACID (important for replay)
• Row-level locking
• JSONB support
• Strong indexing
• Partition support later
Scaling plan to 600 users:
• Add read replica
• Add connection pooling
• Add proper indexes
No schema rewrite needed.
3. EVENT SYSTEM (Conservative but
Migration-Ready)
In-Process Partitioned Dispatcher
But with:
• Events persisted in DB
• Partition key = crop_instance_id
• Clear dispatcher abstraction interface
Important:
We design dispatcher like this:
class EventDispatcherInterface:
def publish(event)
def process(event)
Later migration path:
Replace implementation with Pub/Sub
No business logic change.

===== PAGE 3 =====
That’s critical.
4. ML INTEGRATION (Strong but
Controlled)
Risk Model → Scikit-learn
CNN Backend → PyTorch
Edge AI → TensorFlow Lite (future-ready)
Why this hybrid is good:
• Logistic regression easy
• CNN flexible
• Edge deployable
• No external ML microservice needed
Scaling to 600 users?
Yes.
ML inference lightweight for this scale.
Heavy training can be offline batch.
5. MEDIA PROCESSING
Separate Cloud Run Worker
Backend API → queues job → worker processes
Even for 200 users:
Video worker separated prevents overload.
Scaling to 600:
Increase worker concurrency.
No rewrite needed.

===== PAGE 4 =====
6. FRONTEND
Next.js (React-based PWA)
Why conservative:
• Industry standard
• Easy hosting on Vercel or Cloud Run
• PWA ready
• Easy scaling
• Admin panel easy
No future overhaul needed.
7. DEPLOYMENT
Cloud Run Services:
Service 1 → Backend API
Service 2 → Media Worker
Service 3 → Video Worker (optional)
Cloud SQL → DB
Cloud Storage → Media
Cloud Scheduler → Cron jobs
Scaling to 600 users?
• Increase max instances
• Tune DB connections
• Add read replica
Still no major rewrite.
8. AUTHENTICATION
JWT-based auth (internal)
Later:
Can integrate Firebase Auth or OAuth
Without rewriting backend logic.

===== PAGE 5 =====
9. MONITORING
Cloud Logging + Cloud Monitoring
Add alert thresholds early.
Scaling safe.
WHY THIS IS CORRECT FOR YOUR
GOAL
Because you are:
Not building Netflix
Not building Uber
Not building 100k user system
You are building:
A logic-heavy, state-sensitive, ML-bounded agricultural system.
The bottleneck is logic integrity — not request throughput.
MIGRATION READINESS PLAN
(Important)
We embed these decisions NOW to avoid overhaul later:
1. No direct in-memory state reliance
2. All events persisted before processing
3. Dispatcher interface abstracted
4. ML separated logically from CTIS
5. Workers containerized
6. No synchronous heavy ML in API path
7. Feature flags implemented early
This ensures:

===== PAGE 6 =====
Moving to:
• Pub/Sub
• Redis
• Read replicas
• Horizontal scaling
Requires config change, not rewrite.
🏗 WHAT THIS STACK CAN HANDLE
Users Status
100 Extremely safe
200 Comfortable
500 Safe with tuning
600 Manageable with DB tuning + scaling
1000+ Requires Pub/Sub + DB replica
Perfect alignment with your roadmap.
\end{verbatim}
\section{MSDD Enhancements}
\begin{verbatim}






CultivaX
MSDD — CONTENT ENHANCEMENT ANALYSIS
Deep-Dive Review of All 13 Sections + Missing Architecture Modules





February 2026

Preamble: Scope & Methodology
This document is a comprehensive content enhancement analysis of the CultivaX Master System Design Document (MSDD) and its accompanying Tech Stack specification. Both documents were read in full before a single enhancement was written. The review identifies two categories of gaps:
	•	Intra-section gaps — places where a section exists but a specific sub-specification is absent, undefined, or logically incomplete.
	•	Missing modules — entire architectural concerns that the MSDD does not address at all but that are critical to the system's integrity, usability, or defensibility.
Every enhancement below is purely content-driven. No formatting or structural suggestions are made. Each enhancement specifies exactly what needs to be defined and why it matters for the system.


PART I: Missing Architecture Modules
The following are entirely absent from the MSDD. These are not minor omissions — they represent systems that every user of CultivaX will interact with or that are foundational to the system's correctness.

Missing Module 1 — Notification Architecture (Proposed Section 14)
Alerts are mentioned in Sections 1, 3, 4, 6, and 9 at least 30 times. CTIS transitions to AtRisk state and 'triggers alert.' The media pipeline triggers an alert if a threshold is crossed. Weather updates trigger alerts. But nowhere in the 188-page MSDD is there a specification for how any of these alerts actually work.

14.1 The Alert Emission Problem
Right now, the document describes events like ReplayTriggered, MediaAnalyzed, WeatherUpdated — all of which may result in an alert — but the alert itself is not modelled as a first-class entity. The result is that alert logic will be scattered unpredictably across CTIS, the event dispatcher, and the media pipeline during implementation.
A dedicated notification system must be defined with the following components:

14.2 Alert Entity Model
Each alert must be a structured entity stored in an alerts table:
Field
Type
Notes
alert_id
UUID
Primary key
farmer_id
FK
Target recipient
crop_instance_id
FK nullable
Relevant crop, if any
alert_type
TEXT
E.g. AtRisk, IrrigationDue, HarvestWindow, WeatherWarning, PestProbable
urgency_level
ENUM
Low / Medium / High / Critical
source_event_id
FK
Event that generated this alert
message_key
TEXT
Translation key (not raw string)
delivery_channels
JSONB
Which channels are active
status
TEXT
Pending / Delivered / Acknowledged / Dismissed
created_at
TIMESTAMP
Creation time
delivered_at
TIMESTAMP
First successful delivery
acknowledged_at
TIMESTAMP
Farmer explicitly acknowledged
expires_at
TIMESTAMP
After which alert is stale

14.3 Alert Types and Their Triggers
The following mapping must be explicit in the specification:
Alert Type
Triggering Event
Urgency
AtRisk transition
State machine → AtRisk
Critical
IrrigationDue
Rule engine window open
Medium
FertilizerDue
Rule engine window open
Medium
HarvestApproaching
Days to harvest < threshold
High
WeatherWarning
WeatherUpdated → risk spike
High
PestProbable
pest_density_index > threshold
High
StageDeviation
Delay threshold crossed
Medium
SyncRequired
Device offline > X days
Low
ServiceAvailable
SOE provider matched
Low
ReplayFailed
Replay engine error
Critical (admin only)

14.4 Notification Throttling Policy
Without throttling, the notification system becomes noise. The MSDD never defines this. The following rules must be enforced:
	•	No more than 3 alerts per crop instance per 24-hour window, except Critical urgency which bypasses this limit.
	•	Duplicate alerts of the same alert_type for the same crop instance must be suppressed if one is already in Pending or Delivered status and not yet expired.
	•	IrrigationDue and FertilizerDue alerts must be suppressed once the relevant action has been logged (ActionLogged of matching type resets the alert).
	•	Quiet hours must be configurable per farmer (default 10PM–6AM local time). Low and Medium urgency alerts are batched until quiet hours end. Critical alerts are never suppressed.

14.5 Delivery Channels & Fallback Chain
The MSDD defines WhatsApp and web/PWA as interfaces but never specifies how alerts travel through them. The delivery chain must be:
	•	In-app push notification (PWA service worker) — primary for online users.
	•	WhatsApp notification message — for users in WhatsApp mode or when offline threshold exceeded.
	•	SMS fallback — for Critical alerts only, if WhatsApp delivery not confirmed within X minutes.
Each channel must update the alert entity with delivery_channel_used and delivered_at. If all channels fail, the alert must surface prominently on next login.

14.6 Alert and Event Dispatcher Interaction
Alerts must be generated by the Event Dispatcher, not by individual modules. Each module (CTIS, ML, media pipeline) emits its events as normal. The Event Dispatcher, upon processing certain events, checks alert rules and emits a separate AlertGenerated event. This keeps alert logic centralized and prevents each module from independently managing notifications — which would lead to inconsistency.
Isolation Rule
Alert generation logic must live only in the Event Dispatcher layer. CTIS, SOE, and the ML subsystems must never directly write to the alerts table. They emit events; the Dispatcher decides whether an alert is warranted based on configurable alert rules stored in an alert_rules configuration table.

Missing Module 2 — Recommendation Engine (Proposed Section 15)
The word 'recommendation' appears over 40 times in the MSDD. Yet there is no specification for how recommendations are generated, what data they consume, how they are formatted, or how they are ranked. This is not a minor detail — for a low-literacy farmer checking their phone in the morning, the quality and clarity of today's recommendation is the entire value proposition of CultivaX.

15.1 Recommendation vs Alert: The Missing Distinction
The MSDD conflates alerts (urgent, event-triggered, push-delivered) with recommendations (advisory, proactive, calendar-based). These are fundamentally different objects:
Attribute
Alert
Recommendation
Origin
Event-triggered
Time/stage driven + ML assisted
Urgency
Always notable
Advisory by default
Delivery
Push / WhatsApp
In-app card, updated daily
Persistence
Until acknowledged
Refreshed each morning
Volume
Low (exceptions)
1–3 per day per crop
Authority
High — something happened
Suggestive — farmer decides

15.2 Recommendation Data Model
A Recommendation entity must be stored and retrieved, not computed on every UI load:
	•	recommendation_id (UUID)
	•	crop_instance_id (FK)
	•	recommendation_type: irrigation / fertilizer / pest_check / harvest_prep / market_timing / general_advisory
	•	priority_rank: 1 = most important today (only 1–3 shown to farmer)
	•	message_key: translation key, never raw string
	•	message_parameters: JSONB — dynamic values injected into translation template (e.g. { days_delayed: 3 })
	•	basis: rule_engine / ml_risk / behavioral_adaptation / regional_cluster — source transparency
	•	valid_from / valid_until: the window during which this recommendation is current
	•	status: Active / Acted / Dismissed / Expired
	•	generated_at: timestamp of last computation

15.3 Recommendation Priority Formula
When multiple recommendations are valid simultaneously, ranking must be deterministic. The proposed formula:
Priority Score = (urgency_weight × stage_criticality) + (risk_index × 0.4) + (days_until_deadline × -0.1)
Where urgency_weight is derived from recommendation_type (harvest_prep > irrigation > fertilizer > general), stage_criticality is 1.0 in critical growth windows and 0.5 otherwise, and days_until_deadline penalizes recommendations for distant actions.
The top 3 scoring recommendations per crop instance are surfaced per day. All others are retained but suppressed unless the farmer explicitly requests 'show all.'

15.4 Recommendation and Behavioral Adaptation Integration
The MSDD defines the behavioral adaptation model (Section 4.2, Layer 3) — the system tracks that a farmer consistently irrigates 2 days late. But it never says how this changes the recommendation. The specification must define: if a farmer's historical deviation offset for irrigation is +2 days, the recommendation must surface 2 days earlier than the rule engine's window opens. This offset must be stored as behavioral_recommendation_offset per action_type per farmer and must be applied at recommendation generation time, not at rule engine time.

Missing Module 3 — WhatsApp Conversation State Management (Add to Section 7)
Section 7.12 defines WhatsApp's structured menu format. But a critical question is left entirely unanswered: when a farmer replies '1 — Completed' to a WhatsApp message, how does the system know which crop and which action this reply refers to? The document assumes structured input but never designs the session state that makes structured input unambiguous.

7.12.1 WhatsApp Session Entity
A whatsapp_sessions table must be defined:
	•	session_id (UUID)
	•	farmer_phone (indexed)
	•	session_type: action_confirmation / alert_response / crop_selection / media_request
	•	context_data: JSONB — stores the pending action, crop_instance_id, expected options
	•	initiated_at: timestamp
	•	expires_at: timestamp (sessions expire after a defined window, e.g. 30 minutes)
	•	status: Pending / Responded / Expired / Cancelled

7.12.2 Session Lifecycle
When the system sends a structured WhatsApp menu to a farmer, it must create a session record before sending. The farmer's reply is matched against the active session by phone number. If no session exists for that phone, the reply is invalid and the system must respond with a 'session expired' message plus an instruction to return to the app.
Session expiry is critical: if a farmer leaves a menu unanswered and responds 30 minutes later, the contextual data (which crop, which action window) may be stale. The system must not silently accept stale replies as current actions. Expired sessions must be logged and the farmer notified to re-initiate.

7.12.3 Multi-Crop Disambiguation
If a farmer manages more than one active crop, every WhatsApp interaction that is crop-specific must either be initiated by the system with crop_instance_id baked into the session context, or the farmer must first complete a crop selection step before any action is offered. The system must never assume which crop a generic reply refers to.

Missing Module 4 — Farmer Data Ownership & Portability (Add to Section 11)
Section 11 defines threats from attackers, admins, and providers. But it never addresses the farmer as a data subject with rights over their own data. For a system that captures a farmer's entire agricultural lifecycle — every irrigation decision, every yield, every deviation — data ownership is not a feature, it is a trust foundation.

11.19 Farmer Data Rights Model
The following rights must be defined and technically implemented:

Right to Export: A farmer must be able to request a full export of their data including: all crop instance metadata, all action logs with effective dates, all yield records (reported only — not ml_yield_value), all recommendation history, all media analysis results (extracted features, not raw images). The export must be in a human-readable format (JSON or CSV). This must be an API endpoint: GET /api/v1/farmers/{id}/export.

Right to Deletion: When a farmer requests account deletion, the system must: soft-delete all crop instances and action logs, hard-delete all personal identifiers (name, phone), anonymize all ML training data associated with that farmer, delete behavioral adaptation profiles, and preserve anonymized regional contribution. A farmer's deletion must not corrupt regional clusters already computed from their data — those are prospective aggregates that cannot be retroactively removed.

Behavioral Profile Transparency: The behavioral adaptation model (Section 4.2, Layer 3) silently tracks farmer patterns. The farmer has a right to know this is happening. The UI must surface a 'Your Farming Patterns' view showing what behavioral offsets are applied to their recommendations. The farmer must be able to request a reset of their behavioral profile, which clears deviation_profiles and restarts adaptation from zero.

Missing Module 5 — Seasonal Window Calendar Definition (Add to Section 1)
Section 1.9 defines seasonal_window_category as {Early, Optimal, Late} and states that each has separate risk multipliers and yield sensitivity. But nowhere in the 188-page document is it specified how the system determines which category applies to a given crop instance. This is a foundational gap — the category drives stress sensitivity, yield projection, and risk threshold, yet its assignment algorithm is never specified.

1.9.2 Seasonal Window Assignment Model
seasonal_window_category must be assigned at crop creation time using the following logic:
	•	Each CropRuleTemplate must include a regional_calendar section per crop_type per region. This defines Optimal sowing window as a date range, e.g. Wheat in Haryana: November 1 – November 25. Early is defined as any sowing date before the optimal start. Late is any sowing date after the optimal end.
	•	At CropInstance creation, compare sowing_date against the regional_calendar for that crop_type and region.
	•	Assign seasonal_window_category accordingly and persist it. It must never change after assignment — even if the farmer modifies the sowing date, the original category is retained unless admin explicitly overrides it.

1.9.3 Regional Calendar Data Source
Regional calendars must be stored in a dedicated regional_sowing_calendars table with fields: crop_type, region, optimal_start (DATE), optimal_end (DATE), version_id, effective_from_date. This table is admin-managed and must follow the same versioning philosophy as CropRuleTemplates. When calendar data is unavailable for a specific region/crop combination, the system must fall back to a national-level default and flag the instance as calendar_fallback = true.

Missing Module 6 — CropRuleTemplate Governance & Authoring (Add to Section 1)
The MSDD treats CropRuleTemplates as pre-existing artifacts that are versioned and applied. But it never defines who creates them, through what interface, with what review process, or how errors in them are detected and mitigated. This is a serious governance gap — a badly authored template could silently give wrong irrigation windows to every wheat farmer in Haryana.

1.4.1 Template Authoring Workflow
CropRuleTemplates must only be creatable by Admin role. The authoring workflow must include:
	•	Draft stage: Admin creates template with draft status. Draft templates are never applied to active CropInstances.
	•	Simulation stage: Admin triggers a simulation using Rule Change Impact Preview (already mentioned in Section 1.4 but not specified). The simulation must show: projected stage timelines for 3 representative crop instances (Early, Optimal, Late seasonal windows), deviation from currently applied templates, and estimated yield impact for the next season.
	•	Approval stage: A second admin must approve the template before it moves to Active status. This dual-approval requirement must be enforced at the API layer. Single-admin environments (prototype) may skip this but must log a warning.
	•	Activation: Once approved, the template receives an effective_from_date. It applies only to CropInstances created after that date.

1.4.2 Template Validation Rules
Before any CropRuleTemplate is saved as Active, the system must enforce:
	•	All stage_definitions must be in chronological order with no gaps.
	•	irrigation_windows, fertilizer_windows, and harvest_windows must fall within valid stage date ranges.
	•	Total stage duration must match the known agronomic lifecycle for that crop type (validated against a reference range stored in a crop_reference_data table).
	•	risk_parameters must be numeric and within [0, 1] bounds.
Validation failures must return a structured list of errors, not a generic rejection.

Missing Module 7 — Market Intelligence Data Layer (Expand Section 4.7)
Section 4.7 specifies ARIMA for short-term market forecasting and seasonal bands for long-term forecasting. But it never defines the data layer that feeds these models. The market intelligence model cannot be built, tested, or defended without a specification of its data sources, refresh cadence, and failure behavior.

4.7.2 Market Data Source Specification
The market model must integrate with a configurable primary market data provider (e.g., Agmarknet for India, or a commodity exchange feed). The integration must define:
	•	Commodity code mapping: A market_commodity_map table must link crop_type to the external provider's commodity code (e.g. 'wheat' → 'WHEAT_HARYANA_MSP'). This table is admin-managed.
	•	Refresh cadence: Market price data must be fetched daily via scheduled job (Cloud Scheduler). Intraday updates are unnecessary and computationally wasteful for this system scale.
	•	Data format: Incoming data must be normalized to a standard structure: { crop_type, region, price_per_unit, unit, date, source_provider } and stored in a market_prices table.
	•	Retention: Historical price data must be retained permanently to feed ARIMA training. Only the last 2 years are actively used in model training.

4.7.3 Market API Failure Handling
Section 4.8 explicitly defines weather API failure fallback. Market data has no equivalent. The following must be defined: If the market API fails, use the most recent market_prices entry for that crop_type and region. Mark the forecast as market_data_stale = true. Display to farmer: 'Market prices may not be current.' Never use stale data older than 30 days without admin notification. If no market data exists for a crop_type, disable market intelligence for that crop and display: 'Market data not available for this crop.'

Missing Module 8 — Edge AI Model Distribution & Update Mechanism (Add to Section 6)
Section 6.4 specifies that Edge AI uses a lightweight CNN (MobileNet-based) running on-device as TensorFlow Lite. But nothing specifies how the model arrives on the device, how it is updated when a new version is trained, how version compatibility is enforced, or what happens when the on-device model is outdated relative to backend expectations.

6.4.4 Edge AI Model Distribution Architecture
The PWA must implement a model update mechanism using the following design:
	•	The model binary (.tflite file) must be hosted in Cloud Storage with a versioned path, e.g. gs://cultivax-models/edge/v{N}/stress_classifier.tflite.
	•	The backend must expose a public endpoint: GET /api/v1/edge-model/manifest — which returns { current_version, download_url, checksum, min_compatible_backend_version, file_size_bytes }.
	•	The PWA must check this manifest on startup and after reconnecting from offline. If the cached model version does not match current_version, the PWA must download and cache the new model in the background. The old model continues to serve until the new one is fully downloaded and checksum-verified.
	•	A max_model_age_days threshold must be defined. If the device has been offline longer than this threshold and reconnects, backend inference must be preferred over the potentially stale edge model until the new model is cached.

6.4.5 Edge-Backend Version Compatibility Contract
Each edge model version must declare min_compatible_backend_version. The backend must declare the minimum edge model version it accepts results from (min_acceptable_edge_version). If an edge inference result arrives with an edge_model_version below the backend's minimum threshold, the result must be discarded and backend re-inference triggered. This prevents semantic mismatch between what the edge model classifies as 'high stress' versus what the backend expects.


PART II: Intra-Section Content Enhancements
The following enhancements address gaps within sections that already exist in the MSDD. Each one identifies a specific sub-specification that is absent or logically incomplete.

Section 1 (CTIS) — Enhancements

Enhancement: Stage-Aware Alpha in Stress Smoothing
Section 1.10 defines the stress smoothing formula as new_stress = α × current_signal + (1−α) × previous_stress, with α as a configurable constant (example: 0.4). This is correct but insufficiently specified.
A fixed alpha is inappropriate across all growth stages. In early-stage crops (germination, seedling), the plant has low biomass and is genuinely fragile — stress signals should be weighted more heavily (higher alpha, more reactive). In late-stage crops (grain fill, ripening), the plant has committed to its trajectory and stress signals are less actionable — a lower alpha (more stable, less reactive) better reflects agronomic reality.
The specification must define a stage_alpha_profile per CropRuleTemplate: a mapping from growth_stage to alpha value, allowing different stress sensitivity at different life stages. The formula remains identical; only the alpha parameter becomes stage-dependent.

Enhancement: What-If Engine — Output Contract & Comparison
Section 1.14 defines What-If Engine mechanics (deep copy, isolation, no mutation). But it defines no output contract. What exactly does the farmer see after running a simulation? The specification must define:
	•	What-If output must include: projected_stage_after_N_days, new_risk_index, new_stress_score, projected_harvest_window_shift (in days, positive or negative), estimated_yield_impact_percentage (range, not precise value), and a plain-language assessment key for the UI.
	•	Farmers must be able to run up to 3 What-If scenarios per crop instance per day (to prevent replay engine abuse through excessive simulation).
	•	Scenario comparison: if a farmer runs 2 What-If scenarios, the UI must support side-by-side comparison of the two outputs. The comparison data must be ephemeral — stored only in frontend state, never persisted to the database.

Enhancement: Sowing Date Modification — Replay Scope
Section 1.17 states that sowing date modification must trigger full replay. But full replay means replaying from sowing_date, which is now different. The specification must clarify: the replay must use the new sowing_date as the temporal anchor. All relative day numbers (current_day_number, baseline_day_number) must be recalculated relative to the new sowing_date. Actions that pre-date the new sowing_date must be flagged with a pre_sowing_date_action warning and logged but not applied in timeline recalculation. If the modification shifts the seasonal_window_category, admin must be notified — seasonal category changes are rare and potentially significant.

Section 2 (SOE) — Enhancements

Enhancement: Provider Specialization & Crop Alignment
Section 2.4 defines service_types as a general list (equipment / labor / transport / storage). But in agriculture, a tractor for wheat tillage is not the same as a harvester for rice. The service provider model needs crop-type alignment to be practically useful.
The service_providers entity must include a crop_specializations field (JSONB, list of crop_types the provider has experience with). The marketplace filtering must allow farmers to filter by crop_type, not just service_type. TrustScore computation should include a specialization_match_bonus when a provider's crop_specializations overlaps with the requesting farmer's crop_type — this rewards relevant experience.

Enhancement: Provider Response Time Metric
Section 2.8.1 defines ResponseTime as 'normalized response latency' but never specifies what this means. Response time from what to what? The specification must define: ResponseTime = time between ServiceRequest created and provider first contact attempt, where 'contact attempt' is defined as the provider acknowledging the request through the platform. This requires a provider_acknowledged_at field on ServiceRequest. Without this, ResponseTime cannot be computed and the trust formula has an undefined input.

Section 3 (Event Dispatcher) — Enhancements

Enhancement: Emergency Preemption for Critical Weather Events
Section 3.8 defines a priority model (High/Medium/Low). But the model applies 'only within same partition.' This creates a critical gap: if a weather API update indicates an imminent frost event in the next 24 hours, this is a time-critical cross-crop concern. A low-priority trust-score recalculation in the same partition should not delay it, but a long-running replay in that partition currently would.
The specification must define an Emergency priority tier above High. Emergency-priority events must be able to interrupt a currently running replay (pause it, save interim state as a partial snapshot, process the emergency event, then resume replay). This interrupt mechanism requires a checkpoint within the replay algorithm every N actions. The Emergency tier must be used exclusively for: extreme weather alerts, admin-initiated kill switches, and system health degradation events.

Enhancement: Event Expiry Policy
The event_log table has created_at and processed_at but no expires_at. For certain event types, processing a stale event is worse than not processing it. For example, a WeatherUpdated event that was queued during a 6-hour system outage should not be processed as current weather data when the system recovers. The specification must define: each event_type must have a max_processing_age_minutes in the alert_rules configuration. Events older than this threshold must be moved to dead_letter_events with failure_reason = 'expired_before_processing' rather than being executed with stale data.

Section 4 (ML Architecture) — Enhancements

Enhancement: Behavioral Adaptation Season Reset
Section 4.2 Layer 3 defines the behavioral adaptation model: the system tracks that a farmer consistently deviates (e.g., irrigates 2 days late) and shifts recommendations accordingly. But it defines no reset mechanism between growing seasons.
This is a serious logical gap. A farmer may have consistently delayed irrigation in Season 1 because their pump was broken — not because they prefer late irrigation. If this behavioral offset persists into Season 2 after the pump is repaired, the system will give systematically wrong early recommendations. The specification must define: behavioral offsets are season-scoped. At the start of each new CropInstance for the same farmer and crop_type, behavioral offsets from the previous season must carry a decay weight (e.g., 0.5). After two seasons without the pattern recurring, the offset must reset to zero. Farmers must also be able to manually reset their behavioral profile.

Enhancement: Regional Cluster Confidence Interval
Section 5.6.2 defines regional_clusters with avg_delay, avg_yield, and sample_size. The regional learning model in Section 4.2 Layer 4 uses these clusters to adjust recommendations. But there is a fundamental statistical problem: a cluster based on 5 samples and a cluster based on 500 samples both influence recommendations with equal weight. This is incorrect.
The regional_clusters table must add: std_dev_delay (FLOAT), std_dev_yield (FLOAT), and confidence_interval_95 (JSONB storing {lower, upper} for yield). The regional adaptation offset applied to recommendations must be weighted by a data_confidence_factor = min(1.0, sample_size / minimum_confident_sample_size). This means small-sample clusters have attenuated influence. Large-sample clusters have full influence. This protects farmers in sparsely populated regions from being over-influenced by a handful of atypical cases.

Enhancement: ML Inference Caching Policy
The risk model runs for every crop instance and is influenced by: stress_score, seasonal_risk_index, delay_days, weather_risk, recent deviation count, and stage position. Most of these change slowly. If a farmer's stress_score hasn't changed and weather risk is stable, re-running the logistic regression on every API call is computationally redundant at scale.
The specification must define an ML inference cache: risk inference results must be cached per crop_instance with a last_inference_at timestamp. Inference must only be re-run when: stress_score changes by more than a delta threshold, weather_risk changes, a new action is logged, or the cache is older than X hours. This cache must be stored as a field on crop_instances (last_risk_probability, last_inference_at) — not as a separate service, to keep the architecture conservative.

Section 5 (Data Model) — Enhancements

Enhancement: action_logs Compression Strategy
Section 1.19 states that maximum supported actions per crop is 500–1000 and replay latency target is < 2 seconds for 500 actions. Section 5.3.3 defines the action_logs table. But with JSONB metadata per action and full timestamp data, 1000 action records could be substantial. The specification must define a metadata compression policy: for completed (Harvested or Closed) crop instances, metadata JSONB fields for Informational-category actions (which have no replay impact per Section 1.6.3) may be compressed to a summary hash after a defined archival period. This preserves replay determinism while reducing storage footprint at scale.

Enhancement: Multi-Tenant Land Parcel Model
The current data model assumes one farmer owns one contiguous land area. The land_area field on crop_instances is a single FLOAT. But many Indian farmers cultivate multiple non-contiguous parcels — a farmer might have 3 acres in one village and 2 acres in another, both growing the same crop but in slightly different micro-climates.
The specification must introduce a land_parcels table: { parcel_id, farmer_id, parcel_name, region, sub_region, land_area, soil_type (optional JSONB), gps_coordinates (optional JSONB) }. CropInstances must optionally reference a parcel_id. When parcel data is available, regional risk calculations must use the parcel's sub_region, not just the broader region. This does not require GPS hardware — a farmer can describe their parcel by village name.

Section 6 (Media Intelligence) — Enhancements

Enhancement: Voice Note Intent Parsing
Section 6.8 defines voice media handling: converted to text, stored as transcript, parsed for keywords, used as contextual metadata. But 'parsed for keywords' is significantly underspecified. If a farmer says 'aaj meine pani diya' (today I irrigated), this should ideally be parseable as an irrigation action suggestion, not just metadata.
The specification must define a voice_intent_extraction layer: after speech-to-text, the transcript must be passed through an intent classifier that outputs a structured intent_proposal: { suggested_action_type, confidence, extracted_date_reference }. This proposal must be presented to the farmer for confirmation before any action is logged — voice can suggest, but the farmer must confirm. If intent confidence is below threshold, the transcript is stored as metadata only. The confirmation step protects chronological integrity.

Section 7 (Accessibility) — Enhancements

Enhancement: Accessibility Settings Synchronization
Section 7.14 stores accessibility_settings as JSONB in the farmers table. Section 5.3.1 shows preferred_language as a separate TEXT field in the same farmers table. This creates a synchronization risk: language preference is stored in two locations (as a top-level field and potentially within accessibility_settings if a language toggle is added), which could diverge.
The specification must consolidate: preferred_language must be moved into accessibility_settings JSONB as a primary field. A migration script must be defined to populate existing records. All API endpoints that return farmer data must expose a single source of truth for language preference. The frontend must read only from accessibility_settings for all display customization including language.

Enhancement: Offline Accessibility Mode Persistence
Section 7.3 defines accessibility modes as toggles stored per-user. But if a farmer is offline, their accessibility settings should still apply to locally cached content. The spec never addresses this. The PWA's IndexedDB offline cache must include accessibility_settings as a cached record, refreshed on every successful sync. This ensures that Lite Mode, High Contrast, and Low Literacy Mode all function correctly when the farmer has no connectivity — which is precisely when they most need the app to work.

Section 8 (API Contracts) — Enhancements

Enhancement: GET /crops Response Pagination & Filtering
Section 8.4 defines API endpoints for crop instances but only specifies GET /crops/{id} (single instance). A farmer with multiple active crops needs a list endpoint. The specification must define: GET /api/v1/crops with query parameters: state (filter by CropInstance state), crop_type (filter by type), include_archived (boolean, default false), page (integer), page_size (integer, default 10, max 50). The response must include a pagination envelope: { total_count, page, page_size, results: [...] }. Without this endpoint, the frontend cannot render a multi-crop dashboard.

Enhancement: Optimistic UI Contract for Action Logging
Section 8.4.3 defines action logging as POST /crops/{id}/actions which triggers an ActionLogged event. In high-latency rural networks, waiting for the full event processing cycle before showing the farmer feedback is poor UX. The specification must define an optimistic UI contract: the API must immediately return a 202 Accepted response with { action_id, status: 'queued', estimated_processing_ms } as soon as the action is persisted and the event is queued. A separate GET /api/v1/crops/{id}/actions/{action_id}/status endpoint must allow the frontend to poll for processing completion. The UI must optimistically show the action as logged and update with the final state when the poll resolves.

Section 9 (Deployment) — Enhancements

Enhancement: Database Connection Exhaustion Guard
Section 9.5 recommends PostgreSQL with Cloud SQL. Section 9.12 targets 100–200 users. Cloud Run containers are ephemeral and can spin up rapidly. If each container instance opens its own PostgreSQL connection pool (e.g., 5 connections per instance) and Cloud Run scales to 20 instances under load, that is 100 simultaneous connections — well within Cloud SQL limits. But the replay engine, video worker, and background jobs each need their own connections. The specification must define a connection budget: max_connections per service must be declared and summed to ensure the total never exceeds 80% of the Cloud SQL instance's max_connections limit. PgBouncer or Cloud SQL's built-in connection pooler must be mandated as soon as user count exceeds 200.

Enhancement: Secrets Rotation Policy
Section 9.14 mentions secrets stored in Secret Manager. But there is no rotation policy. JWT signing keys, database passwords, and WhatsApp webhook secrets must all be rotated. The specification must define: JWT signing keys must be rotated every 90 days. Database passwords must be rotated every 180 days. WhatsApp webhook secret must be rotated after any suspected compromise. All services must support hot rotation — reading the new secret without a container restart — using Secret Manager's versioning with automatic latest-version resolution.

Section 10 (Lifecycle Walkthrough) — Enhancements

Enhancement: First Crop Onboarding Flow
Section 10.2 begins with 'Step 1: Farmer Registration' and immediately moves to 'Step 2: Create Crop Instance.' But between registration and first crop creation, there is a critical UX gap that the MSDD ignores entirely: how does a farmer who may have limited smartphone literacy know what to do next?
The specification must define a FirstCropOnboarding state on the farmer entity (is_onboarded: BOOLEAN, default false). Until is_onboarded = true, the UI must present a simplified guided flow: step 1 (language selection), step 2 (accessibility mode selection), step 3 (crop type selection with visual icons for each crop type), step 4 (land area input with voice option), step 5 (sowing date with a calendar picker and voice option), step 6 (confirmation screen). Only after the first crop is successfully created does is_onboarded = true. Subsequent crop creation uses the standard API flow.

Enhancement: Yield Submission Walkthrough for First-Time Farmers
Section 10.7 defines yield submission at harvest. But for a first-time farmer, 'yield' is an abstract concept. The specification must define that when state transitions to ReadyToHarvest, the system must trigger a contextual guidance sequence: explain what yield means, provide the unit (kg/acre or quintal/acre as per farmer's region preference), show the biological maximum for this crop as a reference, and present a simple numeric input. The guidance sequence must be skippable for experienced users (those who have submitted yield before).

Section 11 (Security) — Enhancements

Enhancement: ML Poisoning via Media Manipulation
Section 11.6 addresses yield-based ML poisoning (false yield inflation). But there is an unaddressed vector: a sophisticated user could systematically upload high-quality images of healthy crops while reporting high stress, or vice versa. Over time, if these manipulated media features are used in training, the CNN classifier's stress calibration drifts.
The specification must define a media-feature cross-validation gate: before any media feature is admitted to the ML training dataset, it must pass a consistency check: if extracted stress_probability is high but farmer-reported stress outcome is consistently low over multiple instances, flag the farmer's media contributions for review. This cross-validation uses the yield_verification_score concept from Section 4.3 — extend it to media contribution scoring. A media_contribution_reliability_score per farmer must gate their media features from entering training if the score drops below threshold.

Section 12 (Testing) — Enhancements

Enhancement: Recommendation Engine Test Suite
Section 12 defines testing for CTIS, ML, media, event ordering, replay, and security. But it has no test category for the recommendation engine — which, per Part I of this document, is a missing module that needs to be built. When the recommendation engine is specified and implemented, the test suite must include: priority ranking determinism tests (same inputs must always produce the same ranked recommendation list), behavioral offset application tests (offset = +2 days must move recommendation 2 days earlier), suppression tests (acted recommendation must not resurface), and multi-crop interference tests (recommendations for Crop A must not bleed into Crop B's recommendation list).

Enhancement: Notification System Test Suite
Similarly, the notification architecture (Missing Module 1) requires a dedicated test category: throttle enforcement tests (no more than 3 alerts per crop per 24 hours), deduplication tests (same alert_type cannot have two Pending entries for the same crop), quiet hours suppression tests, delivery channel fallback tests (simulate WhatsApp failure → verify SMS or in-app fallback), and session expiry tests for WhatsApp (expired session reply must be rejected).

Section 13 (Roadmap) — Enhancements

Enhancement: Notification & Recommendation Layer in Phase Sequence
Section 13.2 through 13.10 define 10 development phases. The notification architecture and recommendation engine — both identified as missing modules in this analysis — are not placed anywhere in the roadmap. This creates an implementation gap: developers will build CTIS, the event dispatcher, and the UI without knowing when notifications and recommendations are built.
The notification architecture must be added as a deliverable in Phase 2 (Event Dispatcher stabilization), since alerts depend on the event system. The recommendation engine must be added as a deliverable in Phase 3 (media pipeline phase), since recommendations need stress and risk data to be meaningful. Both must be listed in the Phase 1 exit criteria as interfaces-only (data model defined, no delivery logic yet) to ensure the schema supports them from the start.

Enhancement: Explicit Data Migration Phase
Section 13 defines phases for feature development but has no phase for data migration and backfill. In a multi-phase build, the schema evolves: Phase 1 creates farmers, crop_instances, action_logs. Phase 2 adds event_log. Phase 3 adds media_records. If a Phase 2 feature needs a field that was not in Phase 1's schema, existing records need migration. Without a formal migration policy in the roadmap, schema changes during development will cause ad-hoc SQL edits that violate Section 5.15's migration policy.
A Migration Discipline must be added to each phase: before any new phase begins, a migration script for the schema changes of that phase must be written, tested on a copy of staging data, and approved. This is not a new phase — it is a gate condition for each existing phase transition.


PART III: Cross-Cutting Architectural Concerns
The following enhancements address concerns that span multiple sections and cannot be attributed to a single section.

Cross-Cut 1 — Multi-Region Farmer: Region Consistency
The current data model places region on both the farmers table (as a primary region) and on each crop_instance (which can differ). The MSDD never addresses what happens when a farmer creates a crop in a region different from their registered region. This affects: seasonal_window_category assignment (which region's calendar applies?), regional clustering (which cluster does this crop belong to?), weather risk (which weather API zone?), and SOE provider discovery (which region's providers are shown?).
The specification must state clearly: all regional computations use the crop_instance's region field, not the farmer's registered region. The farmer's region field is used only for initial defaults and contact information. This rule must be explicit to prevent implementation ambiguity.

Cross-Cut 2 — Translation Key Architecture
Section 7.13 says 'All UI text stored in translation files' and 'No hard-coded strings.' The recommendation engine specification (Part I, Missing Module 2) introduces message_key and message_parameters as the mechanism for dynamic translated messages. The notification system introduces the same. But neither the MSDD nor this analysis has defined the translation key naming convention or the parameterized message format.
A translation key architecture must be defined: keys must follow the format {module}.{category}.{identifier}, e.g. ctis.alert.at_risk, soe.review.submitted, recommendation.irrigation.due_tomorrow. Parameterized messages must use a named placeholder convention: { { crop_type } } irrigation is due in { { days } } days. The translation file format must be defined (JSON recommended). All code that emits user-facing messages must use only message_key + parameters, never raw strings — enforced by a linting rule during development.

Cross-Cut 3 — Audit Log Coverage Completeness
The AdminAuditLog is well-specified for admin actions. But the MSDD has a subtle gap: farmer self-service destructive actions (delete crop, dismiss alert, reset behavioral profile, request data export) are not defined as audit events. If a farmer's data export is requested and later disputed, there is no record of the request.
The specification must define a farmer_action_audit_log table separate from AdminAuditLog, recording: farmer_id, action_type (export_request / crop_archive / alert_dismiss / profile_reset / accessibility_change), entity_id, timestamp, and ip_address. This is not for moderation — it is for dispute resolution and for regulatory compliance if CultivaX is ever audited for data handling practices.

Cross-Cut 4 — Graceful Degradation Completeness
The MSDD defines graceful degradation for ML failure (fall back to rules), media failure (don't update stress), weather API failure (use historical baseline), and market API failure (this analysis added a policy). But there is no unified graceful degradation policy — each subsystem independently falls back without a system-wide degradation mode.
A system degradation state must be defined: the backend must maintain a system_health table with one row per subsystem (ML, weather_api, market_api, media_worker, event_dispatcher) and a status field (Operational / Degraded / Down). On startup and on scheduled checks, each subsystem health is verified. When any subsystem is Degraded or Down, a DegradedModeActive system flag must be set. The frontend must display a system status indicator informing the farmer of what is currently limited. This transforms silent degradation into transparent degradation — which is far more appropriate for a system that farmers depend on for agricultural decisions.


Summary: Enhancement Inventory
The following table catalogues all enhancements in this document for quick reference:

#
Enhancement
Category
Priority
1
Notification Architecture (Sec 14)
Missing Module
Critical
2
Recommendation Engine (Sec 15)
Missing Module
Critical
3
WhatsApp Session State Management
Missing Module
Critical
4
Farmer Data Portability & Ownership
Missing Module
High
5
Seasonal Window Calendar Definition
Missing Module
Critical
6
CropRuleTemplate Governance & Authoring
Missing Module
High
7
Market Data Integration Architecture
Missing Module
Medium
8
Edge AI Model Distribution Mechanism
Missing Module
High
9
Stage-Aware Alpha in Stress Smoothing
Section 1 (CTIS)
Medium
10
What-If Engine Output Contract
Section 1 (CTIS)
High
11
Sowing Date Modification — Replay Scope
Section 1 (CTIS)
High
12
Provider Specialization & Crop Alignment
Section 2 (SOE)
Medium
13
Provider Response Time Metric Definition
Section 2 (SOE)
Medium
14
Emergency Event Preemption
Section 3 (Events)
Medium
15
Event Expiry Policy
Section 3 (Events)
Medium
16
Behavioral Adaptation Season Reset
Section 4 (ML)
High
17
Regional Cluster Confidence Interval
Section 4 (ML)
High
18
ML Inference Caching Policy
Section 4 (ML)
Low
19
action_logs Compression Strategy
Section 5 (Data)
Low
20
Multi-Tenant Land Parcel Model
Section 5 (Data)
Medium
21
Voice Note Intent Parsing
Section 6 (Media)
Medium
22
Accessibility Settings Synchronization
Section 7 (Access.)
Medium
23
Offline Accessibility Mode Persistence
Section 7 (Access.)
High
24
GET /crops Pagination & Filtering
Section 8 (API)
Critical
25
Optimistic UI Contract for Action Logging
Section 8 (API)
Medium
26
DB Connection Exhaustion Guard
Section 9 (Deploy)
Medium
27
Secrets Rotation Policy
Section 9 (Deploy)
High
28
First Crop Onboarding Flow
Section 10 (Lifecycle)
High
29
Yield Submission Guidance for First-Timers
Section 10 (Lifecycle)
Medium
30
Media-Based ML Poisoning Prevention
Section 11 (Security)
High
31
Notification & Recommendation Test Suite
Section 12 (Testing)
High
32
Notification & Rec. Layer in Phase Sequence
Section 13 (Roadmap)
High
33
Explicit Data Migration Phase Gate
Section 13 (Roadmap)
Medium
34
Multi-Region Farmer Consistency Rule
Cross-Cutting
High
35
Translation Key Architecture
Cross-Cutting
Medium
36
Farmer Self-Service Audit Log
Cross-Cutting
Medium
37
Unified System Degradation State
Cross-Cutting
High


This analysis identified 8 missing modules and 29 intra-section enhancements. The most critical category is the missing modules — these represent entire architectural concerns that must be specified before implementation begins in their respective phases. The intra-section enhancements are important refinements that improve the system's correctness, fairness, and defensibility.
\end{verbatim}
\section{MSDD & TDD Patch}
\begin{verbatim}
# FILE: MSDD & TDD patch.pdf
# PAGES: 70



===== PAGE 1 =====
CULTIVAX PATCH & ENHANCEMENT DOCUMENT
(MSDD + TDD Consolidation Patch)
Version: 1.0
Patch Type: Comprehensive Structural & Enhancement Patch
Applies To:
• MSDD v1
• TDD v1
• Event Dispatcher Design
• ML Architecture
• SOE Design
• Infrastructure Model
SECTION 1 — CRITICAL STRUCTURAL
PATCHES ( Mandatory Before Production
Build)
This section contains architectural corrections that must be implemented before any large-scale deployment
or handover to another engineering team.
Each patch will include:
• Executive Context
• Problem Statement
• Root Cause Analysis
• Technical Modification (DB / Code / Flow)
• Architectural Impact
• Risk if Not Implemented
• Implementation Scope
1.1 EVENT SCHEMA VERSIONING &
UPCASTER LAYER

===== PAGE 2 =====
Executive Context
CultivaX uses an immutable event log with deterministic replay.
However, there is currently no mechanism to handle schema evolution of events.
If the event structure changes in future versions, historical events may become incompatible with new replay
logic.
This is a structural integrity risk.
Problem Statement
Current event_log structure does NOT include:
• schema version tracking
• payload transformation logic
• compatibility handling for older events
Because replay reconstructs state from events, any change in:
• payload shape
• field naming
• business rule expectations
can break historical reconstruction.
Root Cause
The event system was designed for immutability but not for version-aware immutability.
Immutability alone is insufficient in long-lived systems.
Technical Modification
Database Change
ALTER TABLE event_log
ADD COLUMN schema_version INT DEFAULT 1;

===== PAGE 3 =====
ALTER TABLE event_log
ADD COLUMN correlation_id UUID;
Optional (recommended):
CREATE INDEX idx_event_schema_version
ON event_log(schema_version);
Dispatcher Layer Modification
Introduce Upcaster Layer between event fetch and event handler execution.
Current Flow:
Fetch event → Execute handler
Updated Flow:
Fetch event → Upcast if needed → Execute handler
Upcaster Implementation
Example structure:
class EventUpcaster:
def upcast(self, event):
if event.schema_version == 1:
return self._v1_to_current(event)
if event.schema_version == 2:
return self._v2_to_current(event)
return event
def _v1_to_current(self, event):
payload = event.event_payload
# Example transformation
if "stressValue" in payload:
payload["stress_probability"] = payload.pop("stressValue")
event.schema_version = CURRENT_SCHEMA_VERSION
event.event_payload = payload
return event

===== PAGE 4 =====
Versioned Event Handlers
Event handlers must assume current schema only.
Upcaster ensures:
• Handlers never deal with legacy structures.
• Replay remains deterministic.
Architectural Impact
✔ Preserves event immutability
✔ Enables safe long-term evolution
✔ Allows addition of new fields
✔ Prevents replay breakage
✔ Supports migration to Pub/Sub safely
Risk If Not Implemented
Risk Severity
Replay failure after deployment High
Historical crop state corruption High
Forced data migration High
Loss of trust in deterministic engine Critical
Implementation Scope
• Modify DB schema
• Modify dispatcher
• Add Upcaster class
• Update testing to include legacy event replay
• Add integration test: “Replay 6-month-old event log”
Estimated effort: 1–2 days.

===== PAGE 5 =====
1.2 CORRELATION ID PROPAGATION
(FULL TRACEABILITY)
Executive Context
Observability exists but lacks unified request-to-event trace linkage.
Cross-service debugging is incomplete.
Problem Statement
Currently:
• API generates response
• Events are created
• Replay happens
• ML inference occurs
• Media worker processes
But there is no guaranteed shared trace ID across:
• API request
• event_log
• replay logs
• ML logs
• media logs
Root Cause
Tracing was implemented partially via request_id in response meta but not enforced across all subsystems.
Technical Modification
API Entry Point
Generate:

===== PAGE 6 =====
request_id = uuid4()
Attach to:
• request context
• structured logs
When Publishing Event
dispatcher.publish(
event_type="ActionLogged",
payload=data,
crop_instance_id=crop_id,
correlation_id=request_id
)
Store in DB.
Replay Logging
Each replay log entry must include:
{
"correlation_id": "...",
"event_id": "...",
"replay_duration": 48,
"stage_before": "...",
"stage_after": "..."
}
ML Inference Logging
{
"correlation_id": "...",
"model_version": "risk_v1.0",
"confidence": 0.82
}
Media Worker Logging
When processing MediaUploaded:
• Read correlation_id from event payload

===== PAGE 7 =====
• Propagate in worker logs
Architectural Impact
✔ Enables distributed tracing
✔ Simplifies peer testing
✔ Strengthens audit trail
✔ Enables production debugging
Risk If Not Implemented
Risk Severity
Difficult bug tracing Medium
Poor observability under load Medium
Failure root cause ambiguity High
Implementation Scope
• Modify dispatcher publish signature
• Modify event_log schema
• Modify all structured logs
• Update tests
Estimated effort: 1 day.
1.3 SEASONAL WINDOW ASSIGNMENT
ENGINE
Executive Context
Seasonal windows influence risk multipliers but no assignment engine exists.

===== PAGE 8 =====
Problem Statement
Currently:
• Crop stores region
• Seasonal category concept exists
• But system does not calculate whether sowing was early/optimal/late.
This affects:
• Risk calculation
• Stress integration
• Drift sensitivity
Root Cause
Seasonal modeling defined but not operationalized.
Technical Modification
New Table
CREATE TABLE seasonal_windows (
id UUID PRIMARY KEY,
crop_type TEXT NOT NULL,
region TEXT NOT NULL,
optimal_start DATE NOT NULL,
optimal_end DATE NOT NULL,
risk_multiplier NUMERIC(5,4) NOT NULL
);
Crop Creation Logic
At CropCreated event:
window = find_seasonal_window(crop_type, region)
if sowing_date < window.optimal_start:
category = "Early"
elif sowing_date > window.optimal_end:
category = "Late"
else:

===== PAGE 9 =====
category = "Optimal"
crop_instance.seasonal_category = category
crop_instance.seasonal_risk_multiplier = window.risk_multiplier
Risk Integration Update
Risk index becomes:
risk_index = base_risk * seasonal_risk_multiplier
Architectural Impact
✔ Adds agronomic realism
✔ Strengthens research defensibility
✔ Reduces false anomaly alerts
Risk If Not Implemented
Risk Severity
Agronomic inaccuracy High
False risk alerts Medium
Weak viva defense Medium
Implementation Scope
• Create new table
• Seed seasonal data
• Modify CropCreated handler
• Update risk calculation
• Add tests
Estimated effort: 2–3 days (including data seeding).
1.4 CROP RULE TEMPLATE
GOVERNANCE MODEL

===== PAGE 10 =====
Executive Context
CropRuleTemplate defines:
• Growth stages
• Stage durations
• Drift bounds
• Stress sensitivity
• Structural rules
These templates directly influence deterministic replay.
Currently, there is no governance lifecycle around these templates.
That is dangerous.
Problem Statement
Current system assumes:
rule_version_applied INT
But it does NOT define:
• Who creates rule versions?
• How they are validated?
• Who approves activation?
• How rollback occurs?
• How regional override is handled?
• Whether rule change retroactively affects active crops?
This creates:
• Governance ambiguity
• Replay instability risk
• Administrative corruption risk
Root Cause
Rule templates were treated as static configuration, not as version-controlled domain policy.

===== PAGE 11 =====
Technical Modification
New Governance Table
CREATE TABLE crop_rule_templates (
rule_version INT PRIMARY KEY,
crop_type TEXT NOT NULL,
region TEXT NOT NULL,
template_json JSONB NOT NULL,
is_active BOOLEAN DEFAULT FALSE,
approved_by UUID REFERENCES users(user_id),
approved_at TIMESTAMPTZ,
created_at TIMESTAMPTZ DEFAULT NOW()
);
Rule Activation Flow
Current (Unsafe):
Admin edits template → system uses it
Updated (Safe):
Admin creates draft template
↓
Run simulation validation
↓
Admin approves
↓
Mark is_active = TRUE
↓
Future crops use this rule_version
Important Constraint
Rule changes must NOT retroactively affect existing crop_instances.
Each crop_instance stores:
rule_version_applied
Replay always uses stored version.
This preserves deterministic history.

===== PAGE 12 =====
Simulation Validation Endpoint
Before activation:
POST /api/v1/admin/rule-template/validate
System:
• Runs sample replay using template
• Checks invariants
• Ensures no infinite drift
• Validates stage continuity
If validation fails:
Activation blocked.
Architectural Impact
✔ Protects deterministic replay
✔ Prevents accidental rule corruption
✔ Enables region-specific policy
✔ Supports future agronomic research
Risk If Not Implemented
Risk Severity
Accidental replay corruption Critical
Retroactive state mutation Critical
Loss of historical consistency Critical
Implementation Scope
• DB schema creation
• Admin endpoints
• Simulation engine integration
• Test coverage
Estimated effort: 2–3 days.

===== PAGE 13 =====
1.5 REGIONAL CLUSTER STATISTICAL
CONFIDENCE WEIGHTING
Executive Context
Regional learning adapts crop behavior based on cluster analysis.
However, sample size confidence is not weighted.
This can produce statistically invalid behavior.
Problem Statement
Current logic (conceptual):
cluster_mean_offset = compute_cluster_mean()
apply cluster_mean_offset
If cluster_size = 5 farmers,
it affects behavior as much as cluster_size = 500 farmers.
That is mathematically wrong.
Root Cause
Cluster model implemented without confidence calibration.
Technical Modification
Modify Adaptation Formula
Replace:
applied_offset = raw_offset

===== PAGE 14 =====
With:
confidence_weight = min(1.0, cluster_size / minimum_required_size)
applied_offset = raw_offset * confidence_weight
Example:
Cluster Size Minimum Required Confidence Effect
5 30 0.16 Very weak adaptation
15 30 0.5 Moderate
30+ 30 1.0 Full
Log Confidence Value
Add to adaptation log:
{
"cluster_size": 12,
"confidence_weight": 0.4,
"applied_offset": 2.3
}
Architectural Impact
✔ Preserves statistical defensibility
✔ Prevents early model corruption
✔ Makes research publication stronger
Risk If Not Implemented
Risk Severity
Overfitting to small region High
Farmer trust degradation Medium
Academic defensibility weak High
Implementation Scope
• Modify regional cluster engine
• Update logs
• Add test cases for cluster_size thresholds
Estimated effort: 1 day.

===== PAGE 15 =====
1.6 BEHAVIORAL ADAPTATION RESET
POLICY
Executive Context
Deviation logic adapts crop stage based on behavioral patterns.
But adaptation persists indefinitely.
This creates cross-season contamination.
Problem Statement
If farmer consistently irrigates late in Season 1,
system adapts stage offset.
But in Season 2:
• Weather changes
• Soil changes
• Crop variety changes
Adaptation persists incorrectly.
Root Cause
No lifecycle boundary enforcement for deviation profiles.
Technical Modification
At Harvested Event
When processing:

===== PAGE 16 =====
CropHarvested
Add logic:
archive_deviation_profile(crop_instance_id)
reset_deviation_profile(crop_instance_id)
Archive Table
Optional but recommended:
CREATE TABLE deviation_profile_archive (
id UUID PRIMARY KEY,
crop_instance_id UUID,
archived_profile JSONB,
archived_at TIMESTAMPTZ
);
Architectural Impact
✔ Prevents systemic bias
✔ Maintains season independence
✔ Aligns with agronomic cycles
Risk If Not Implemented
Risk Severity
Bias accumulation High
False stage prediction next season Medium
Implementation Scope
• Modify Harvested handler
• Add archive table
• Add tests
Estimated effort: 1 day.

===== PAGE 17 =====
1.7 ORPHAN ACTION POLICY
Executive Context
SowingDateModified is a structural event.
If sowing date moves forward, historical actions may become invalid.
Replay will crash due to invariant violation.
Problem Statement
Invariant:
action_date >= sowing_date
If sowing date changes from Jan 1 → Jan 10,
and action on Jan 5 exists:
Replay fails.
Root Cause
Replay invariants assume chronological consistency.
Structural modification does not re-evaluate past actions.
Technical Modification
Add Column
ALTER TABLE action_logs
ADD COLUMN is_orphaned BOOLEAN DEFAULT FALSE;

===== PAGE 18 =====
On SowingDateModified
Algorithm:
for action in actions:
if action.action_effective_date < new_sowing_date:
action.is_orphaned = True
Replay Adjustment
Replay must filter:
valid_actions = [a for a in actions if not a.is_orphaned]
Admin Visibility
Admin endpoint:
GET /api/v1/admin/orphan-actions/{crop_id}
Architectural Impact
✔ Prevents replay crash
✔ Preserves historical audit
✔ Avoids destructive deletion
Risk If Not Implemented
Risk Severity
Replay engine crash Critical
Data corruption Critical
Implementation Scope
• DB migration
• Update SowingDateModified handler
• Update replay filter
• Add admin endpoint
• Add test case

===== PAGE 19 =====
SECTION 2 — HIGH FUNCTIONAL
COMPLETENESS PATCH ( )
These are not structural survival fixes like Section 1,
but they are necessary to make CultivaX operationally complete, defensible, and production-ready.
Each patch will include:
• Executive Context
• Problem Statement
• Technical Modification
• Architectural Impact
• Risk if Not Implemented
• Implementation Scope
2.1 CROP TERMINATION WORKFLOW
Executive Context
Currently, a crop can transition:
Active → AtRisk → Harvested → RecoveryRequired
However, there is no proper lifecycle path for crop failure such as:
• Flood
• Pest destruction
• Drought
• Farmer abandonment
• Land reallocation
Deleting crop data would violate audit integrity.
This is a lifecycle gap.
Problem Statement
Without a termination workflow:

===== PAGE 20 =====
• Farmers may delete crops improperly
• Historical data may be lost
• Risk modeling becomes incomplete
• Insurance/analysis use-case is broken
Technical Modification
Add New Event
CropTerminated
Payload:
{
"reason": "Flood",
"termination_date": "YYYY-MM-DD"
}
Extend Crop State
Update crop_instances:
ALTER TABLE crop_instances
ADD COLUMN termination_reason TEXT;
Add new allowed state:
Terminated
Valid transitions:
Active → Terminated
AtRisk → Terminated
No transition from Terminated back to Active.
Replay Behavior
Upon CropTerminated event:
• Stop future stage progression
• Stop stress accumulation

===== PAGE 21 =====
• Freeze risk_index
• Mark lifecycle closed
Architectural Impact
✔ Preserves historical truth
✔ Enables future insurance integration
✔ Supports yield failure modeling
✔ Maintains event integrity
Risk If Not Implemented
Risk Severity
Improper deletion High
Loss of agronomic history High
Weak lifecycle modeling Medium
Implementation Scope
• Add event handler
• Extend state enum
• Modify replay terminal logic
• Add admin visibility
Estimated effort: 1 day.
2.2 NOTIFICATION MODULE (ALERT
DELIVERY LAYER)
Executive Context
CTIS generates:
• Risk changes

===== PAGE 22 =====
• Stress spikes
• Replay results
• Dead letters
• Service updates
But there is no structured notification system.
Currently, UI must poll state manually.
This is a delivery gap.
Problem Statement
Without notifications:
• Farmers miss critical alerts
• UX suffers
• Risk warnings are passive
• System feels non-responsive
Technical Modification
New Table
CREATE TABLE notifications (
id UUID PRIMARY KEY,
user_id UUID REFERENCES users(user_id),
crop_instance_id UUID,
message TEXT NOT NULL,
notification_type TEXT,
status TEXT CHECK (status IN ('Unread','Read')),
created_at TIMESTAMPTZ DEFAULT NOW()
);
Notification Triggers
Create notifications when:
• risk_index > threshold
• crop enters AtRisk

===== PAGE 23 =====
• event enters DeadLetter
• service request accepted/completed
• replay completed (optional)
API Endpoints
GET /api/v1/notifications
POST /api/v1/notifications/{id}/mark-read
Future Extension
Optional:
• SMS
• WhatsApp integration
• Push notifications
But backend layer must exist first.
Architectural Impact
✔ Separates state from delivery
✔ Improves UX
✔ Enables asynchronous alerting
✔ Keeps CTIS pure
Risk If Not Implemented
Risk Severity
Poor user experience Medium
Missed critical risk alerts High
Implementation Scope
• DB schema
• Event-to-notification mapping
• API endpoints
• UI integration

===== PAGE 24 =====
Estimated effort: 2 days.
2.3 GEO-POLYGON SUPPORT (POSTGIS
INTEGRATION)
Executive Context
Currently crop_instances store:
region TEXT
This is too coarse for:
• Hyperlocal weather modeling
• GPS validation of media
• Regional clustering precision
Problem Statement
District-level region is insufficient.
Weather variation can occur within 2–5 km.
Technical Modification
Enable PostGIS
CREATE EXTENSION postgis;
Add Geo Column
ALTER TABLE crop_instances
ADD COLUMN geo_polygon GEOGRAPHY(POLYGON, 4326);

===== PAGE 25 =====
Weather Query Change
Instead of:
get_weather(region)
Use:
SELECT weather
FROM weather_grid
WHERE ST_Intersects(geo_polygon, grid_polygon)
Architectural Impact
✔ Enables hyperlocal modeling
✔ Enables geo-validation
✔ Improves ML regional clustering
✔ Future IoT-ready
Risk If Not Implemented
Risk Severity
Coarse risk modeling Medium
Inability to validate media GPS High
Implementation Scope
• DB extension
• Migration
• Update weather logic
• Update cluster grouping
Estimated effort: 2–3 days.
2.4 MEDIA GEO-VERIFICATION

===== PAGE 26 =====
Executive Context
A farmer could upload:
• Random internet image
• Old image
• Different field image
Without geo-verification, ML can be poisoned.
Problem Statement
Media currently validated only by:
• file type
• size
• quality
No spatial verification.
Technical Modification
Extract EXIF GPS
During worker processing:
gps_data = extract_exif_gps(image)
Validate
if not ST_Contains(crop.geo_polygon, gps_point):
flag_media_as_unverified()
Add Field
ALTER TABLE media_files
ADD COLUMN geo_verified BOOLEAN DEFAULT FALSE;

===== PAGE 27 =====
Architectural Impact
✔ Prevents media fraud
✔ Protects ML from poisoning
✔ Strengthens audit
Risk If Not Implemented
Risk Severity
ML corruption High
Fake stress manipulation Medium
Implementation Scope
• EXIF parser
• Geo validation query
• UI indicator
Estimated effort: 1–2 days.
2.5 MARKET API FALLBACK POLICY
Executive Context
Market integration assumed stable.
No fallback defined.
Problem Statement
If market API:
• Times out
• Returns error
• Returns partial data

===== PAGE 28 =====
System may produce null risk or crash.
Technical Modification
Cache Last Known Market Price
Store:
CREATE TABLE market_cache (
crop_type TEXT,
region TEXT,
last_price NUMERIC,
updated_at TIMESTAMPTZ
);
Fallback Logic
If live API fails:
price = cached_price
confidence *= 0.7
log_event("MarketFallbackUsed")
Architectural Impact
✔ Graceful degradation
✔ Prevents crash
✔ Improves robustness
Risk If Not Implemented
Risk Severity
API dependency fragility Medium
Runtime errors Medium

===== PAGE 29 =====
Implementation Scope
• Cache table
• API wrapper
• Retry logic
• Fallback event
SECTION 3 — ML & INTELLIGENCE
ENHANCEMENTS ( Recommended / Some
Upgradeable)
This section strengthens:
• ML defensibility
• Farmer trust
• Model lifecycle management
• Data integrity
• Research readiness
These are not replay survival fixes, but they determine whether CultivaX is:
A basic rule engine with ML garnish
or
A defensible Human-AI agricultural system.
We proceed patch-by-patch.
3.1 ACTIVE LEARNING FEEDBACK LOOP
(Human-in-the-Loop)
Executive Context
Current ML flow:
Data → Train → Predict → Apply bounded influence
But there is no structured mechanism for:

===== PAGE 30 =====
• Farmer disagreement
• Model correction signals
• Feedback prioritization
Without feedback loop, ML stagnates.
Problem Statement
When ML predicts:
Risk = High
Farmer may:
• Ignore
• Override
• Say “Crop is fine”
This disagreement is a high-value training signal.
Currently, this signal is lost.
Root Cause
System logs predictions but not explicit disagreement events.
Technical Modification
Add New Event Type
PredictionRejected
Payload example:
{
"prediction_type": "Risk",
"predicted_value": 0.82,
"confidence": 0.74,
"user_response": "Healthy",

===== PAGE 31 =====
"crop_stage": "Tillering"
}
Add Feedback Table
CREATE TABLE ml_feedback (
id UUID PRIMARY KEY,
crop_instance_id UUID REFERENCES crop_instances(crop_instance_id),
prediction_type TEXT,
predicted_value NUMERIC,
confidence_score NUMERIC,
user_feedback TEXT,
created_at TIMESTAMPTZ DEFAULT NOW()
);
Training Pipeline Change
During retraining:
priority_samples = feedback_samples + random_samples
Feedback samples receive higher weight.
Architectural Impact
✔ Enables adaptive learning
✔ Strengthens human-AI co-pilot model
✔ Supports research publication
✔ Improves long-term model performance
Risk If Not Implemented
Risk Severity
Model stagnation Medium
Farmer distrust Medium
Reduced learning efficiency High

===== PAGE 32 =====
Implementation Scope
• Event type addition
• DB schema
• Update training pipeline
• Add UI “Dismiss Alert” option
Estimated effort: 2 days.
3.2 EXPLAINABILITY TOKENS (WHY,
NOT JUST WHAT)
Executive Context
Current ML output:
{
"risk_probability": 0.72,
"confidence": 0.81
}
This is opaque.
Farmers trust explanations, not numbers.
Problem Statement
Without explainability:
• ML appears arbitrary
• Overrides increase
• Trust decreases
• Research defensibility weakens

===== PAGE 33 =====
Technical Modification
Modify Risk Model Output
Add:
{
"risk_probability": 0.72,
"confidence": 0.81,
"top_contributing_factors": {
"rain_deficit": 0.45,
"deviation_trend": 0.21,
"late_sowing": 0.12
}
}
Implementation Method
For Logistic Regression:
coefficients = model.coef_
feature_importance = feature_value * coefficient
Select top 3 absolute contributors.
Store in Inference Log
{
"model_version": "risk_v1.0",
"factors": {...}
}
Architectural Impact
✔ Improves farmer trust
✔ Supports academic defense
✔ Enables transparency
✔ Facilitates model debugging

===== PAGE 34 =====
Risk If Not Implemented
Risk Severity
ML seen as black box Medium
Low adoption Medium
Implementation Scope
• Modify risk inference function
• Update structured logging
• Update UI display
Estimated effort: 1 day.
3.3 EDGE MODEL VERSIONING
CONTRACT
(Upgradable to if edge deployment becomes primary)
Executive Context
System uses:
• Backend CNN (PyTorch)
• Optional Edge AI (TFLite)
Currently:
No strict model version enforcement across edge/backend.
This can create inference inconsistency.
Problem Statement
Edge device may run:
stress_model_v1
Backend may use:

===== PAGE 35 =====
stress_model_v2
Without version tracking:
• Results inconsistent
• Debugging difficult
• Drift detection impossible
Technical Modification
Extend media_files Table
ALTER TABLE media_files
ADD COLUMN model_version TEXT;
Worker Must Emit Version
{
"stress_probability": 0.42,
"confidence": 0.78,
"model_version": "cnn_v1.2"
}
Replay Log Must Store Version Used
So that historical inference reproducible.
Version Compatibility Rule
If edge_version != backend_version:
• Reduce confidence weight slightly
• Log compatibility event

===== PAGE 36 =====
Architectural Impact
✔ Enables model rollback
✔ Improves observability
✔ Supports federated learning future
Risk If Not Implemented
Risk Severity
Inference inconsistency Medium
Debugging impossible Medium
Implementation Scope
• DB migration
• Update worker payload
• Update replay logging
Estimated effort: 1 day.
3.4 SMALL DATASET SAFEGUARD
ENFORCEMENT
(Upgrade to if production ML enabled)
Executive Context
Early deployment will have:
• Small dataset
• Sparse feedback
• Weak statistical foundation
If ML fully enabled early:
System will overfit.

===== PAGE 37 =====
Problem Statement
Without threshold enforcement:
• Logistic regression unstable
• Cluster adaptation unreliable
• Stress model biased
Technical Modification
Define Dataset Threshold
Example:
MIN_DATASET_SIZE = 200
MIN_CLUSTER_SIZE = 30
Training Guard
if dataset_size < MIN_DATASET_SIZE:
disable_ml_risk = True
Replay Adjustment
If ML disabled:
skip ML component in stress integration
use rule-based risk only
Admin Panel Indicator
Show:
ML Status: Disabled (Insufficient Data)

===== PAGE 38 =====
Architectural Impact
✔ Prevents overfitting
✔ Protects farmer trust
✔ Improves research integrity
Risk If Not Implemented
Risk Severity
Early model instability High
False risk signals High
Implementation Scope
• Add dataset size check
• Add feature flag integration
• Add admin visibility
Estimated effort: 1 day.
3.5 CONFIDENCE PROPAGATION
HARDENING
Executive Context
Stress integration currently weighted by confidence.
But confidence aggregation not formally standardized.
Problem Statement
Multiple confidence sources exist:
• Model certainty

===== PAGE 39 =====
• Image quality
• Frame variance
• Cluster stability
• Sample size adequacy
But no unified formula enforces bounds.
Technical Modification
Unified Confidence Model
final_confidence = (
w_model * model_confidence +
w_quality * image_quality +
w_cluster * cluster_confidence
)
final_confidence = clamp(final_confidence, 0, 1)
Log Confidence Breakdown
{
"confidence_components": {
"model": 0.8,
"quality": 0.7,
"cluster": 0.5
},
"final_confidence": 0.72
}
Architectural Impact
✔ Formalizes ML authority boundaries
✔ Supports research publication
✔ Prevents hidden signal amplification

===== PAGE 40 =====
Risk If Not Implemented
Risk Severity
Inconsistent stress behavior Medium
Hard-to-debug model impact Medium
Implementation Scope
• Refactor stress integration module
• Update logging
• Add unit tests
SECTION 4 — SECURITY HARDENING &
ABUSE PREVENTION
This section strengthens:
• Replay integrity
• Offline safety
• ML poisoning resistance
• Marketplace fraud resistance
• Data sovereignty
• Audit defensibility
Some items here move from “good design” → “defensible production system”.
We proceed patch-by-patch.
4.1 OFFLINE SYNC ABUSE DETECTION
Executive Context
Offline mode allows:
• Bulk action logging
• Delayed synchronization
• Chronological reconstruction

===== PAGE 41 =====
This is powerful — but dangerous.
Offline mode can be exploited to:
• Inject fabricated historical actions
• Manipulate stage progression
• Overload replay engine
• Fake agricultural compliance
Problem Statement
Current safeguards:
• local_monotonic_counter
• idempotency key
• chronological invariant enforcement
But missing:
• Behavioral anomaly detection
• Density validation
• Time plausibility validation
Threat Model
Malicious user syncs:
50 actions
all with action_effective_date = same timestamp
all within 2 minutes
This is impossible human behavior.
Technical Modification
Action Density Validator
Add validation before replay:

===== PAGE 42 =====
def validate_action_density(actions):
grouped = group_by_date(actions)
for date, daily_actions in grouped.items():
if len(daily_actions) > MAX_ALLOWED_ACTIONS_PER_DAY:
raise AbuseDetectedException()
Recommended:
MAX_ALLOWED_ACTIONS_PER_DAY = 20
Timestamp Pattern Detection
Detect identical timestamp clusters:
if identical_timestamp_count > threshold:
flag_account()
Abuse Log Table
CREATE TABLE abuse_flags (
id UUID PRIMARY KEY,
user_id UUID,
crop_instance_id UUID,
abuse_type TEXT,
details JSONB,
flagged_at TIMESTAMPTZ DEFAULT NOW()
);
Automatic Temporary Lock
If abuse severity high:
Account temporarily suspended pending admin review
Architectural Impact
✔ Protects deterministic timeline
✔ Prevents synthetic replay manipulation
✔ Enhances fraud detection
✔ Strengthens peer audit

===== PAGE 43 =====
Risk If Not Implemented
Risk Severity
Replay manipulation High
Timeline fraud High
ML poisoning Medium
Implementation Scope
• Add validation logic
• Add abuse table
• Integrate with auth lock mechanism
• Admin dashboard display
Estimated effort: 2 days.
4.2 REPLAY ANOMALY DETECTION
LAYER
Executive Context
Replay enforces invariants, but it does not currently monitor suspicious patterns like:
• Excessive structural changes
• Frequent sowing date modifications
• Oscillating yield submissions
These patterns may indicate manipulation.
Problem Statement
A user repeatedly modifies sowing date to manipulate stage.
Replay allows it (with orphan policy).
But frequent structural edits indicate abuse.

===== PAGE 44 =====
Technical Modification
Structural Event Frequency Monitor
Add:
if structural_events_in_last_30_days > threshold:
flag_account()
Example threshold:
MAX_STRUCTURAL_EVENTS_30_DAYS = 3
Yield Oscillation Detection
If:
abs(current_yield - previous_yield) > biological_limit
Flag inconsistency.
Replay Integrity Log Extension
Add:
{
"structural_event_count": 4,
"yield_variance_flag": true
}
Architectural Impact
✔ Detects manipulation patterns
✔ Enhances fraud detection
✔ Protects ML training dataset

===== PAGE 45 =====
Risk If Not Implemented
Risk Severity
Slow timeline manipulation Medium
Undetected yield fraud Medium
Implementation Scope
• Add monitoring logic
• Add structured logging
• Admin dashboard integration
Estimated effort: 1–2 days.
4.3 TRUST SCORE FRAUD HARDENING
(Marketplace Protection)
Executive Context
SOE trust score includes:
• Completion rate
• Complaint ratio
• Rating
• Exposure fairness
But advanced collusion patterns not fully addressed.
Threat Model
• Two accounts repeatedly review each other
• IP clustering abuse
• Rapid rating spikes
• Coordinated fake complaints

===== PAGE 46 =====
Technical Modification
IP Similarity Detection
Log:
{
"reviewer_ip": "...",
"provider_ip": "...",
"ip_similarity_flag": true
}
If similarity frequency high:
Reduce trust weight.
Sudden Rating Spike Detection
if rating_std_dev_last_7_days > threshold:
flag_review_pattern()
Trust Volatility Cap
Prevent trust score from changing more than:
MAX_TRUST_DELTA_PER_DAY = 0.15
Architectural Impact
✔ Marketplace integrity
✔ Prevents trust inflation
✔ Supports scalability
Risk If Not Implemented
Risk Severity
Review manipulation Medium
Marketplace unfairness Medium

===== PAGE 47 =====
Implementation Scope
• Extend trust engine
• Add volatility cap
• Add monitoring
Estimated effort: 1 day.
4.4 ML DATA POISONING DEFENSE
UPGRADE
Executive Context
System already:
• Clamps yield values
• Uses z-score validation
But can be hardened further.
Problem Statement
Attack vector:
User submits extreme yield repeatedly to shift regional cluster mean.
Technical Modification
Outlier Exclusion in Training
During dataset preparation:
if abs(z_score) > 3:
exclude_from_training()

===== PAGE 48 =====
Weighted Sample Contribution
Weight by:
sample_weight = yield_verification_score
Cluster Drift Monitor
Track:
if regional_mean_change > threshold:
freeze_cluster_update()
Architectural Impact
✔ Protects ML stability
✔ Prevents statistical hijacking
✔ Enhances research defensibility
Risk If Not Implemented
Risk Severity
Slow regional drift manipulation Medium
Implementation Scope
• Update training pipeline
• Add cluster drift log
• Add freeze logic
Estimated effort: 1 day.
4.5 DATA EXPORT & SOVEREIGNTY
PROTOCOL
(Important for long-term compliance)

===== PAGE 49 =====
Executive Context
Privacy laws require:
• Data portability
• Right to deletion
• User access transparency
Currently not formally implemented.
Technical Modification
Data Export Endpoint
GET /api/v1/users/{id}/export-data
Generate:
• JSON file
• All crop history
• All action logs
• Yield records
• Media metadata (not raw media unless requested)
Soft Delete Enforcement
Upon user deletion:
is_active = FALSE
Preserve audit history.
Export Format
Zip file containing:
crops.json
actions.json

===== PAGE 50 =====
yield.json
services.json
media_metadata.json
Architectural Impact
✔ Compliance readiness
✔ Builds farmer trust
✔ NGO/government collaboration ready
Risk If Not Implemented
Risk Severity
Legal non-compliance (future) Medium
Trust erosion Medium
Implementation Scope
• Export service
• File generation
• Rate limiting
• Admin audit log
SECTION 5 — INFRASTRUCTURE &
DEVOPS HARDENING
This section strengthens:
• Deployment safety
• Cost efficiency
• Disaster resilience
• Environment isolation
• Scaling discipline
• Operational robustness
These are not domain logic changes, but they determine whether CultivaX is:
A deployable academic system
or
A cloud-native, production-defensible platform.

===== PAGE 51 =====
We proceed patch-by-patch.
5.1 SPOT / PREEMPTIBLE INSTANCE
STRATEGY (COST OPTIMIZATION)
Executive Context
Video processing is:
• CPU-heavy
• Memory-heavy
• Delay-tolerant
• Non-blocking to CTIS
Currently deployed on standard Cloud Run instances.
This is unnecessarily expensive.
Problem Statement
Without cost optimization:
• Video inference costs scale rapidly
• Pilot deployment becomes financially inefficient
• Academic demo budgets may be exceeded
Technical Modification
Deploy Video Worker on Preemptible Compute
If using Cloud Run:
• Use lower min instances
• Configure CPU allocation per request
• Allow instance scale-to-zero

===== PAGE 52 =====
If migrating to GCE:
• Use Spot VMs (preemptible instances)
Ensure Job Idempotency
Because preemptible instances may terminate mid-job:
Video worker must:
if media_processing_started:
resume_or_retry_safely()
Event status reset to Created if crash occurs (already defined in dispatcher).
Cost Impact
Expected savings:
• 50–80% for heavy video workloads
Architectural Impact
✔ Cost-efficient scaling
✔ No impact on CTIS
✔ Worker isolation preserved
Risk If Not Implemented
Risk Severity
High operational cost Medium
Budget constraints Medium
Implementation Scope
• Adjust Cloud Run config
• Validate retry safety
• Add monitoring alert for preemption spikes

===== PAGE 53 =====
Estimated effort: 1 day.
5.2 BACKUP & DISASTER RECOVERY
POLICY
Executive Context
Cloud SQL backups enabled by default, but no formal restore validation process exists.
Backups are useless unless restore is tested.
Problem Statement
Without restore validation:
• Backup corruption may go unnoticed
• Disaster recovery unverified
• Audit compliance incomplete
Technical Modification
Quarterly Restore Simulation
Procedure:
1. Restore Cloud SQL snapshot to staging instance
2. Run replay integrity check
3. Verify:
a. crop_instances consistency
b. event_log integrity
c. no orphaned chain
4. Log restore test result

===== PAGE 54 =====
Backup Verification Log Table
CREATE TABLE backup_verification_logs (
id UUID PRIMARY KEY,
restore_timestamp TIMESTAMPTZ,
verification_status TEXT,
notes TEXT,
verified_by UUID
);
Alerting
If backup age > 48 hours:
Trigger monitoring alert.
Architectural Impact
✔ Ensures resilience
✔ Strengthens production readiness
✔ Supports NGO/government partnerships
Risk If Not Implemented
Risk Severity
Data loss unrecoverable Critical
Compliance failure High
Implementation Scope
• Create restore SOP document
• Add log table
• Configure monitoring
Estimated effort: 2 days.

===== PAGE 55 =====
5.3 ENVIRONMENT SEPARATION
HARDENING
Executive Context
System currently supports:
• Development
• Production
But no strict enforcement of separation.
Problem Statement
Without strict environment separation:
• Dev data may leak into prod
• Accidental migration possible
• Test rules may affect real crops
Technical Modification
Enforce Three Environments
DEV
STAGING
PROD
Each must have:
• Separate Cloud SQL instance
• Separate Cloud Storage bucket
• Separate Secret Manager entries
• Separate Cloud Run services

===== PAGE 56 =====
Environment Guard
Add at app startup:
if ENVIRONMENT not in ["DEV","STAGING","PROD"]:
raise RuntimeError("Invalid environment configuration")
Database Guard
Prevent production DB connection from dev build.
Add:
if ENVIRONMENT == "DEV" and DATABASE_URL contains "prod":
abort_startup()
Architectural Impact
✔ Prevents catastrophic configuration errors
✔ Enables safe CI/CD
✔ Professional-grade DevOps discipline
Risk If Not Implemented
Risk Severity
Production data corruption Critical
Accidental rule activation High
Implementation Scope
• Config validation
• Environment-specific secrets
• Infrastructure separation
Estimated effort: 1–2 days.

===== PAGE 57 =====
5.4 DEPLOYMENT SAFETY GUARDRAILS
Executive Context
CI/CD pipeline exists but lacks safety checks for:
• Schema migrations
• Rule activation
• Model deployment
Problem Statement
Uncontrolled deployment may:
• Introduce incompatible schema
• Break replay
• Activate unvalidated ML
Technical Modification
Pre-Deployment Migration Check
In CI pipeline:
Run alembic migration dry-run
Run replay integrity test
Replay Regression Test
Before deploy:
• Load sample historical event log
• Run replay
• Verify expected state output

===== PAGE 58 =====
ML Model Deployment Guard
New model cannot be activated unless:
admin_approval = TRUE
evaluation_metrics above threshold
Architectural Impact
✔ Prevents silent corruption
✔ Ensures deterministic safety
✔ Strengthens research reproducibility
Risk If Not Implemented
Risk Severity
Broken replay after deploy High
Model instability Medium
Implementation Scope
• Extend CI pipeline
• Add regression dataset
• Add admin activation gate
Estimated effort: 2 days.
5.5 SCALING GUARDRAILS &
MONITORING AUTOMATION
Executive Context
Scaling plan exists conceptually, but guardrails not automated.

===== PAGE 59 =====
Problem Statement
Without automated scaling alerts:
• Event backlog may silently grow
• Replay latency may degrade
• DB connection exhaustion unnoticed
Technical Modification
Monitoring Thresholds
Set alerts:
Event backlog > 50
Replay time > 300ms
DB connections > 80%
CPU > 80%
Backpressure Automation
If backlog high:
reduce replay concurrency
Auto-Scaling Config
Cloud Run:
• Max instances dynamic
• Concurrency limit tuned
Architectural Impact
✔ Self-protecting system
✔ Predictable scaling to 600 users
✔ No surprise performance collapse

===== PAGE 60 =====
Risk If Not Implemented
Risk Severity
Performance degradation Medium
Unexpected outages Medium
Implementation Scope
• Configure Cloud Monitoring
• Add performance logging
• Adjust concurrency controls
SECTION 6 — OPTIONAL RESEARCH &
FUTURE EXTENSIONS ( )
These are not required for MVP or Phase-1 production,
but they elevate CultivaX into a:
• Research-grade platform
• Publication-ready system
• Startup-ready AgTech architecture
• Long-term scalable intelligence ecosystem
These extensions must remain modular and must NOT break deterministic replay.
6.1 THERMAL TIME (GDD MODEL
INTEGRATION)
Executive Context
Current baseline stage progression uses:
days_since_sowing
Biologically inaccurate.

===== PAGE 61 =====
Plants grow based on heat accumulation, not calendar days.
Research Motivation
Growing Degree Days (GDD):
GDD=∑Tmax+Tmin2−Tbase GDD=∑2Tmax +Tmin −Tbase
This accounts for:
• Heatwaves
• Cold delays
• Climate variability
Problem Solved
Without GDD:
• Heatwave accelerations misclassified as anomalies
• Early harvest falsely flagged
• Biological realism limited
Technical Extension Design
Add Thermal Units Field
ALTER TABLE crop_instances
ADD COLUMN accumulated_gdd NUMERIC;
Stage Lookup Changes
Instead of:
baseline_stage = stage_lookup(days_since_sowing)
Use:
baseline_stage = stage_lookup(accumulated_gdd)

===== PAGE 62 =====
Weather Dependency
Requires:
• Daily temperature ingestion
• Weather data reliability
Architectural Impact
✔ Increases agronomic realism
✔ Supports climate adaptation modeling
✔ Research publication potential
Risk
• Adds external API dependency
• Complicates deterministic replay
• Requires temperature versioning
Recommendation
Defer to Phase 3 or research thesis integration.
6.2 SOCIAL VERIFICATION LAYER
(COMMUNITY GROUND TRUTH)
Executive Context
Farmers trust other farmers more than algorithms.
ML-only risk alerts may face skepticism.

===== PAGE 63 =====
Concept
If ML detects pest outbreak in Region X:
Trigger:
Community Poll:
“Are you observing aphids in your field?”
Collect 5–10 nearby farmer confirmations.
Technical Extension
Add:
CREATE TABLE community_validations (
id UUID PRIMARY KEY,
crop_instance_id UUID,
region TEXT,
response BOOLEAN,
created_at TIMESTAMPTZ
);
Confidence boost if majority positive.
Architectural Impact
✔ Human-in-the-loop reinforcement
✔ Improves ML confidence calibration
✔ Supports social agronomy research
Risk
• Requires user density
• Risk of herd bias
• Needs moderation layer

===== PAGE 64 =====
Recommendation
Pilot feature after 500+ active users.
6.3 PRODUCE PASSPORT (PRE-HARVEST
MARKET HOOK)
Executive Context
System knows:
• Projected harvest date
• Stress history
• Yield estimate
• Deviation pattern
This intelligence can be monetized.
Concept
15 days before harvest:
Auto-generate:
Digital Produce Passport
Contains:
• Crop history
• Risk history
• Estimated yield
• Quality confidence
• Harvest date
Shareable link for buyers.

===== PAGE 65 =====
Technical Extension
Add:
CREATE TABLE produce_passports (
id UUID PRIMARY KEY,
crop_instance_id UUID,
projected_yield NUMERIC,
quality_score NUMERIC,
generated_at TIMESTAMPTZ
);
Auto-triggered by scheduled job.
Architectural Impact
✔ Bridges cultivation to marketplace
✔ Enables forward contracts
✔ Startup monetization path
Risk
• Requires buyer ecosystem
• Market pricing integration
Recommendation
Post-MVP monetization layer.
6.4 FEDERATED LEARNING (EDGE
TRAINING)

===== PAGE 66 =====
Executive Context
Edge AI model currently inference-only.
Future: allow on-device model improvement without sharing raw images.
Concept
• Edge device trains locally
• Sends gradient updates
• Backend aggregates
No raw image upload required.
Architectural Prerequisites (Already Met)
✔ Model version tracking
✔ Feature schema versioning
✔ ML audit log
Technical Extension
Add:
CREATE TABLE federated_updates (
id UUID PRIMARY KEY,
model_version TEXT,
gradient_blob BYTEA,
contributed_at TIMESTAMPTZ
);
Aggregate server-side.
Architectural Impact
✔ Privacy-preserving ML
✔ Research-grade capability
✔ Low-bandwidth rural adaptation

===== PAGE 67 =====
Risk
• Complex orchestration
• Gradient poisoning risk
Recommendation
Research-grade extension only.
6.5 BLOCKCHAIN / CRYPTOGRAPHIC
HISTORY ANCHORING
Executive Context
Event log already contains:
• event_hash
• chain_hash (planned)
Can anchor root hash externally.
Concept
At end of crop lifecycle:
Generate Merkle root of event chain.
Anchor hash externally (optional public ledger).
No on-chain logic required.

===== PAGE 68 =====
Technical Extension
Add:
ALTER TABLE crop_instances
ADD COLUMN merkle_root_hash TEXT;
Store final root.
Architectural Impact
✔ Tamper-proof crop history
✔ Bank/insurance proof
✔ Compliance-grade audit
Risk
• Regulatory overhead
• Complexity creep
Recommendation
Optional research extension.
6.6 ADVANCED GEO-SPATIAL
INTELLIGENCE
Executive Context
With PostGIS enabled:
Possible to:
• Cluster farms spatially

===== PAGE 69 =====
• Detect microclimate zones
• Perform NDVI satellite integration
Technical Extension
Add:
CREATE INDEX idx_crop_geo
ON crop_instances
USING GIST(geo_polygon);
Enable spatial clustering algorithms.
Architectural Impact
✔ Research publication potential
✔ Microclimate modeling
✔ NGO-level intelligence dashboards
6.7 IOT SENSOR INTEGRATION MODEL
Executive Context
Future integration:
• Soil moisture sensors
• Temperature sensors
• Irrigation automation
Integration Model
IoT Device → API → Event → ReplayTriggered
No structural change required.

===== PAGE 70 =====
Example Event
{
"event_type": "SoilMoistureReading",
"moisture_value": 0.42
}
Replay integrates as bounded stress modifier.
Architectural Impact
✔ Fully autonomous agriculture support
✔ Smart irrigation foundation
✔ Future funding pathway
\end{verbatim}
\end{document}
